\documentclass[journal]{IEEEtran}
% *** 04.04.2026 (C3 <-> C3')***
\usepackage{amsmath,amssymb,amsfonts}
\usepackage{mathtools}
\usepackage{graphicx}
% \usepackage{cite}
\usepackage{booktabs}
\usepackage{array}
\usepackage{algorithm}
\usepackage{algorithmic}
\usepackage{url}
\usepackage{enumitem}
% \usepackage{biblatex}
\usepackage[backend=bibtex,style=numeric]{biblatex}

% *** THEOREM-LIKE ENVIRONMENTS ***
\newtheorem{theorem}{Theorem}
\newtheorem{lemma}{Lemma}
\newtheorem{proposition}{Proposition}
\newtheorem{corollary}{Corollary}
\newtheorem{assumption}{Assumption}
\newtheorem{definition}{Definition}
\newtheorem{remark}{Remark}
\newtheorem{proof}{Proof}

% Real n-dimensional space
\newcommand{\R}{\mathbb{R}}
% Diagonal matrix operator
\DeclareMathOperator{\diag}{diag}
% Convex hull operator
\DeclareMathOperator{\conv}{conv}
% Filippov convexification of g (closure of conv g)
% Usage: \convg(x) = \overline{\operatorname{conv}}\,g(x)
\newcommand{\convg}{\overline{\operatorname{conv}}\,g}
% \newcommand{\convg}{\overline{\operatorname{conv}}}
\newcommand{\const}{\text{const}}

\usepackage{todonotes}

% real author/affiliation block for the camera-ready version
\title{Global Stability of the Hopfield Neural Networks\\ with Non-Differentiable Activation Functions:\\
            a Yakubovich--Leonov--Gelig Perspective}

\author{Timur N. Mokaev \and Timur A. Kisiev \and Nikolay V. Kuznetsov
      \thanks{Submitted to \emph{IEEE Transactions on Neural Networks and Learning Systems} (TNNLS). 
            For review, author names and affiliations are withheld to preserve double-anonymous peer review.}
}


% *** START DOCUMENT ***
\addbibresource{ref.bib}
\begin{document}

\maketitle

\begin{abstract}
We study global stability of continuous-time Hopfield neural networks with 
\emph{discontinuous, possibly non-monotone and unbounded} activations. 
Modeling the dynamics as a \emph{Filippov differential inclusion}, 
we derive three verifiable conditions (C1--C3) --- linear growth, componentwise one-sided Lipschitz, 
and \emph{diagonal Lyapunov stability} --- under which the network admits a \emph{unique equilibrium} 
that is \emph{globally asymptotically stable}. 
The proof leverages the Yakubovich--Leonov--Gelig (YLG) quadratic-plus-integral functional, 
avoiding smoothness assumptions and extending classical bounded/monotone results. 
We further embed QUBO (including planted MAX-CUT), soft-margin SVM, and constrained optimal control tasks 
into this discontinuous-activation Hopfield template and verify that each satisfies our conditions, with simulations confirming convergence from arbitrary initial conditions. 
These results broaden the class of optimization problems solvable by Hopfield networks with rigorous global guarantees.
\end{abstract}

\begin{IEEEkeywords}
Hopfield neural networks, 
global asymptotic stability, 
discontinuous activation, 
Filippov differential inclusions, 
Yakubovich--Leonov--Gelig theory.
\end{IEEEkeywords}

% \markboth{IEEE Transactions on Neural Networks and Learning Systems}{Anonymous: Global Stability of HNNs with Non-Differentiable Activations}

\section{Introduction}\label{sec:intro}
\IEEEPARstart{H}{opfield} neural networks (HNNs) play a dual role in neural computation.
First, they act as \emph{associative memories} (content--addressable stores) defined by a
decreasing energy function~\cite{hopfield1982}; second, the very same architecture functions as an
\emph{analog optimizer} whose equilibrium encodes a low-cost solution of a constrained,
often combinatorial, problem~\cite{hopfieldTank1985}.
Both viewpoints are naturally unified by \emph{stability theory}.
Two regimes are central: \emph{multistability}---many locally attractive equilibria enabling memory---and
\emph{monostability} (global stability)---a unique globally attractive equilibrium needed for deterministic optimization.
This paper focuses on the second regime while keeping the first in view to emphasize the model's dual nature.

\textbf{Motivation.} Modern large-capacity Hopfield variants adjust the energy/activation design to
improve expressivity---from dense associative memories with polynomial--to--exponential capacity growth
\cite{krotovHopfield2016,demircigil2017} to ``modern'' Hopfield layers
linked to softmax attention~\cite{ramsauer2020}. The resulting activations are
\emph{non-smooth} (hard limiters, quantizers, dead-zones) or \emph{unbounded/continuous} (softmax-type),
and the dynamics is more faithfully modeled as a \emph{Filippov differential inclusion}.
Classical Lyapunov analyses that assume smooth, bounded, monotone nonlinearities are therefore
insufficient to certify global convergence in these practically relevant settings.

\textbf{GAS with discontinuous activations.}
For bounded monotone discontinuous activations, Forti and Nistri established existence of a
\emph{unique globally asymptotically stable} (GAS) equilibrium for HNNs~\cite{p1}.
However, hardware-friendly and optimization-driven designs often deploy
piecewise-linear or hard-limiter maps and unbounded gains; discontinuous penalties also arise in exact
exterior methods. Extending GAS guarantees beyond the bounded/monotone class is thus essential
for positioning HNNs as reliable analog optimizers.

\textbf{YLG perspective.}
The Yakubovich--Leonov--Gelig (futher --- YLG) absolute-stability framework provides a quadratic-plus-integral
Lyapunov functional tailored to Lur\'e systems with sector/slope-restricted nonlinearities~\cite{YLG}.
It is well suited to Filippov right-hand sides and has been refined in subsequent works for
discontinuous dynamics~\cite{Boikov2008,RajchakitAgarwalRamalingam2021}. We build on this perspective.

\subsection{Contributions}\label{sec:intro:contrib}
We consider continuous-time HNNs with componentwise activations that may be discontinuous, non-monotone, and unbounded.
Our main contributions are as follows.
\begin{enumerate}[label=\textbf{(\alph*)},leftmargin=*]
\item \textbf{Filippov HNN model.} We formulate an HNN with \emph{piecewise--continuous} activations
      having finitely many first-order jumps and interpret the dynamics as a \emph{Filippov differential inclusion}.
\item \textbf{Three verifiable hypotheses.} 
      We introduce \textbf{C1--C3} --- linear growth, componentwise one-sided
      Lipschitz, and a \emph{diagonal Lyapunov stability} test for the linear part --- and show via a YLG
      quadratic\,+\,integral functional that \emph{existence, uniqueness, and GAS} of the equilibrium hold.
\item \textbf{Finite-time entrance (optional strengthening).}
      Under an additional mild slope condition we further obtain \emph{finite-time entrance} into any prescribed
      shrinking neighborhood of the equilibrium (full details in the Appendix).
\item \textbf{Optimization embeddings.} We demonstrate that standard penalty embeddings of
      MAX-CUT, TSP,soft-margin SVM and constrained optimal control tasks satisfy \textbf{C1--C3}.
\end{enumerate}

\subsection{Organization}\label{sec:intro:org}
Section~\ref{sec:related} reviews related work.
Section~\ref{sec:prelim} recalls biological and mathematical models of HNNs,
Filippov solutions and the YLG framework.
Section~\ref{sec:problem} states the HNN model in the Lur\'{e} form and assumptions~\textbf{C1--C3}. Section~\ref{sec:nonunique} presents a low-dimensional example showing that 
these assumptions ensure global asymptotic stability of the unique equilibrium 
but do not guarantee forward uniqueness of Filippov trajectories. Section~\ref{sec:main} presents the main results on uniqueness and GAS and discusses tightness.
Section~\ref{sec:examples} contains some numerical illustrations, followed by Sections~\ref{sec:qubo-hnn}, \ref{sec:penalty-HNN} and \ref{sec:application_control} that verify \textbf{C1--C3} for four optimization embeddings.
Conclusions are drawn in Section~\ref{sec:concl}.
% An informative appendix (not used by the proofs) summarizes biological/circuit intuition and the Kirchhoff model; 
Full proofs and algorithmic checks appear in technical appendices.

% \section{Related Work}\label{sec:related}
% Position the paper within discontinuous-activation HNNs, Filippov systems in neural dynamics, and absolute stability (Yakubovich/Popov/Circle). Contrast with diagonal stability and quadratic matrix inequality (QMI/LMI) conditions in NN literature; clarify novelty and tightness.

\section{Related Work}\label{sec:related}

The stability of continuous–time Hopfield–type networks has been surveyed extensively; 
see the comprehensive review of Zhang~\emph{et al.}, 
which compares algebraic-inequality, M-matrix, LMI, and Lyapunov diagonal-stability (LDS) criteria 
and stresses that no single form dominates across models or assumptions~\cite{ZhangTNNLSReview2014}. 
In the same spirit, a recent tutorial review for HNNs collects classical sufficient conditions 
(Gershgorin–disc bounds, diagonal stability, negative semidefiniteness, matrix–norm methods) 
and their scope of applicability~\cite{BoykovRoudnev2022}. 
These surveys also highlight open issues for large–scale 
or nonstandard topologies and for models with non-smooth elements.

\subsection{Discontinuous Activations and Nonsmooth Tools}

A key departure from early HNN analyses is to allow discontinuous 
or merely piecewise–continuous activations. 
Forti and Nistri established global convergence of neural networks 
with discontinuous neuron activations by recasting the dynamics as a differential inclusion 
and extending LaSalle’s invariance principle to nonsmooth settings; 
they also connected Filippov/Clarke tools with Lyapunov arguments 
and reported related finite–time phenomena~\cite{p1}. 
Subsequent work on discontinuous systems sharpened the analytical toolkit 
via generalized gradients and set–valued dynamics, 
yielding Lyapunov and invariance results tailored 
to differential inclusions~\cite{ShevitzPaden1994,CortesTutorial2008,YLG}.
For models interpreted in the Filippov sense, 
existence, uniqueness and ''weak/strong'' stability notions are standard; 
tutorials make precise when Filippov solutions are the physically meaningful choice 
(e.g., circuits with switches, sliding modes) and summarize the calculus rules used in proofs~\cite{CortesTutorial2008}.

\subsection{Multistability vs. Global Stability}
In the attractor–memory regime, multistability is desirable and widely analyzed 
(e.g., piecewise–linear or nonmonotone activations producing many equilibria). 
Representative results of~\cite{NieZheng2015} establish exponential stability 
of multiple equilibria under structured interconnections and discontinuous nonlinearities. 
By contrast, optimization applications require a \emph{unique} globally attracting equilibrium 
to avoid spurious minima. 
Reviews emphasize this distinction and motivate conditions ensuring existence/uniqueness plus GAS 
of the single equilibrium~\cite{ZhangTNNLSReview2014}. 
Moreover, unbounded activations are common in optimization-driven networks 
(e.g., to encode inequality penalties), which pushes beyond the bounded/Lipschitz assumptions 
of classical analyses~\cite{NieZheng2015}.

\subsection{Optimization-oriented Neural Dynamics}

A parallel line of works treat HNN-like systems as continuous-time solvers for constrained, 
possibly nonsmooth programs. 
Li~\emph{et al.} analyze a generalized Hopfield network for nonsmooth constrained convex optimization 
using Filippov inclusions and extended LaSalle arguments; 
their model accommodates exact penalties and constraint qualifications typical of optimization~\cite{LiTNNLS2015}. 
In this way, Hopfield networks are placed in direct correspondence 
with the theory of nonsmooth gradient flows and projected dynamical systems 
developed under differential-inclusion formulations.

% \todo[inline]{New part below}

The foundation for the network’s architecture is the penalty method~\cite{penalty, 517}, where the constrained problem is transformed into an unconstrained minimization through exterior penalty that fits the Hopfield framework, which involves a gradient system of the energy function. This energy function is composed of the objective function and the constraint functions of a particular optimization problem in such a way to provide correlation between optimal solution equation and stationary condition for the network. Solution computation involves providing initial states (voltages) to the neurons, after which the neural network converges to a stable equilibrium point that corresponds to the minimum of the energy function, which in turn corresponds to the minimum
of the cost function. Therefore, it is crucial for such systems to have a unique globally asymptotically stable equilibrium state.

Providing new stability results for larger class of activation functions can subsequently extend the class of problems that the network can solve. This is because the energy function allows for the definition of activation functions, neuron connections,and the input data for the network and represents the main objective function of optimization problem, therefore, particular activation function can reflect the type of problem to solve.

A necessary condition for the convergence of the penalty method in finite time is the use discontinuous functions in the constraints ~\cite{517}. These functions are represented by the activation functions of the neurons in the network, which raises the issue of the stability of networks with non-smooth activations.

For example, HNN with non-smooth activations (e.g., Heaviside step functions) establishes a framework for solving the Quadratic Unconstrained Binary Optimization problem (QUBO). By directly mapping discrete decision variables to the network’s states, one can embed QUBO problems like MAX-CUT into an energy-minimization framework. This approach eliminates the need for stochastic annealing or parameter tuning, offering deterministic convergence
even for indefinite Q matrices.

\textbf{Positioning.}
Our work fits squarely in the ''global-stability/unique equilibrium'' branch 
and is closest in spirit to Forti–Nistri and to optimization-oriented Filippov analyses. 
Relative to these, we 
(i) allow piecewise–continuous activations with finitely many first-order discontinuities 
(including unbounded and nonmonotone cases), 
(ii) provide easy-to-check structural hypotheses (matrix negativity, unilateral Lipschitz-like slope bounds, 
diagonal dominance), and 
(iii) prove finite-time entrance into shrinking neighborhoods, thereby bridging discontinuous activation models 
used in exact-penalty embeddings with the strong GAS guarantees required by analogue optimizers.

\section{Preliminaries}\label{sec:prelim}
\subsection{Notation}

\begin{align*}
&\mathbb{R}^n &\quad&\text{Real } n \text{ space.}\\
&A = [A_{ij}] \in \mathbb{R}^{n \times n} &\quad&\text{Square matrix.}\\
&A^{\top} &\quad&\text{Transpose of } A.\\
&A^{-1} &\quad&\text{Inverse of } A.\\
&A^S = \tfrac{1}{2}(A + A^{\top}) &\quad&\text{Symmetric part of } A.\\
&\text{ker } A &\quad&\text{Kernel of } A.\\
&\lambda_m\{A\} &\quad&\text{Minimum eigenvalue of } A.\\
&\lambda_M\{A\} &\quad&\text{Maximum eigenvalue of } A.\\
&\|A\| &\quad&\text{Matrix norm of } A.\\
&\|A\|_2 = [\Lambda_M\{A^{\top}A\}]^{1/2} &\quad&2 \text{ norm of } A.\\
&\alpha = \text{diag}(\alpha_1, \cdots, \alpha_n) \in \mathbb{R}^{n \times n} &\quad&\text{Diagonal matrix.}\\ 
&x = (x_1, \cdots, x_n)^{\top} \in \mathbb{R}^n &\quad&\text{Column vector.}\\
&\|x\|_2 = \left[\sum_{i=1}^n x_i^2\right]^{1/2} &\quad&\text{Euclidean norm of } x.\\
&\|x\| &\quad&\text{Vector norm of } x \text{,}\\
& & &\text{denotes }\|x\|_2.\\
&\overline{E} &\quad&\text{Closure of set } E \subset \mathbb{R}^n.
\end{align*}



\subsection{Biological and Mathematical Models of HNNs}\label{sec:biomath}

\textbf{Biological motivation.}
Hopfield’s model was inspired by the organization of the brain: 
large populations of neurons interconnected via synapses, 
transmitting electrical signals along axons and receiving them through dendrites (Fig.~\ref{neuron}). 
The underlying biophysics views each neuron as an excitable membrane whose \emph{membrane potential} 
evolves due to ionic currents (primarily Na$^+$ and K$^+$), 
exhibits all–or–nothing \emph{spikes} once a threshold is exceeded, 
and then enters a \emph{refractory} state. 
The classical Hodgkin–Huxley approach~\cite{p224} formalizes this in terms of circuit elements 
(capacitors, voltage–dependent conductances, current sources), 
providing a principled bridge from physiology to circuit models (Fig.~\ref{circuit}). 
In the Hopfield abstraction, the detailed spike shape is replaced 
by a static nonlinear \emph{activation map}, 
while the membrane integrator and synaptic couplings remain 
as the core dynamical ingredients.

\begin{figure*}[!ht]
\begin{minipage}{0.49\linewidth}
\includegraphics[width=1\textwidth]{Figures/neuron_small.png}
\caption{Biological neurons and their connections.} \label{neuron}
\end{minipage}
\hfill
\begin{minipage}{0.49\linewidth}
\centering
\includegraphics[width=0.68\textwidth]{Figures/circuit_mod.png}
\caption{Electronic circuit modeling the neural network.} \label{circuit}
\end{minipage}
\end{figure*}

\medskip
\textbf{Kirchhoff–circuit interpretation.}
Consider a network of $n$ nodes (neurons). At node $j$, let $x_j(t)$ denote the voltage (state), $C_j>0$ the membrane capacitance, and $R_j>0$ the leak resistance. Synaptic couplings from node $i$ to node $j$ are represented by conductances $T_{ji}$, and $g_i(\cdot)$ is the static (possibly non-smooth) activation associated with node $i$. With an external bias current $I_j$, Kirchhoff’s current law gives
\begin{equation}
  C_j\,\dot{x}_j(t) + \frac{1}{R_j}\,x_j(t) =
  \sum_{i=1}^n T_{ji}\, g_i \big(x_i(t)\big) \;+\; I_j,
  \quad j=1,\dots,n.                                    
  \label{eq:kcl}
\end{equation}
Stacking \eqref{eq:kcl} and rescaling by the capacitances yields the standard form of the Hopfield network:
\begin{equation}
  \dot{x}(t) =
  -B\,x(t)\;+\;T\,g \big(x(t)\big)\;+\;I,
  \qquad x(t)\in\mathbb{R}^n,                             
  \label{eq:special}
\end{equation}
where $B=\mathrm{diag}(b_1,\dots,b_n)\succ 0$ collects effective leak rates ($b_j:=1/(R_j C_j)$ after rescaling), $T\in\mathbb{R}^{n\times n}$ is the interconnection matrix (with signs admitting excitatory/inhibitory effects), $g(x)=(g_1(x_1),\ldots,g_n(x_n))^\top$ acts componentwise, and $I\in\mathbb{R}^n$ is the (rescaled) input vector. This compact representation makes explicit the split between linear passive dynamics and static nonlinearities, a structure central to absolute–stability analyses.

% \medskip
% \textbf{Non-smooth activations and Filippov trajectories.}

\subsection{Non-smooth activations and Filippov solutions}

In many applications the activation function in the right-hand side of system~\eqref{eq:special}
is a discontinuous function with a finite number of first-order discontinuities (quantizers, dead zones, sign, or Heaviside).
Furthermore, for the one-sided limit values of the function at the discontinuities,
the relationship $g_i(\rho^+_k)>g_i(\rho^-_k)$ holds.
Within this work such functions are considered to belong to the class~$\mathcal{G'}$.

\smallskip
\begin{definition}[Extended class of activation functions]\label{G_n}
A mapping $g = (g_1,...,g_n)$ belongs to the class $\mathcal{G'}$ if and only if for all coordinate functions $g_i ~ (i \in 1:n)$ the following conditions are met:
The function $g_i$ is piecewise continuous, that is, it is continuous at all points of $\mathbb{R}$ except for a countable number of discontinuity points $\{\rho_k\}$, where finite right- and left-hand limits exist (first-kind discontinuities) and $g_i(\rho^+_k)>g_i(\rho^-_k)$. Moreover, on any compact subset of $\mathbb{R}$, there is a finite number of such discontinuity points.
\end{definition}
\begin{remark}
    It is important to note that, 
    unlike~\cite{p1}, the functions $g_i$ here are not required to be monotonic and bounded.
\end{remark}

% In such cases, system~\eqref{eq:special} is naturally interpreted 
% as a differential inclusion in the sense of Filippov~\cite{p2}: 
% the right–hand side $g$ is replaced by its convexified set–valued map, 
% and solutions are absolutely continuous curves whose derivatives belong to the Filippov set almost everywhere. 
% This preserves the circuit interpretation (e.g., sliding along switching manifolds) 
% and enables nonsmooth Lyapunov/invariance arguments. 
% Formal definitions are recalled further in the next subsection.

% \todo[inline]{!TODO: correct the statement about DS}
% \medskip
% \textbf{The HNN as a dynamical system.}
% Equation~\eqref{eq:special} defines a continuous–time dynamical system (a semiflow) on $\mathbb{R}^n$:
% for each initial condition $x(0)=x_0$, there is a trajectory $t\mapsto x(t;x_0)$ governed by the vector field
% \(
% f(x)=-Bx+Tg(x)+I.
% \)
% Equilibria are the zeros of $f$, i.e., solutions of
% \(
% -Bx^\star + T g(x^\star) + I = 0,
% \)
% and their stability properties determine the qualitative behavior of the network. In optimization–oriented uses of HNNs, one seeks \emph{monostability} (existence, uniqueness, and global asymptotic stability of $x^\star$); in associative–memory uses, one typically exploits \emph{multistability} (many locally attractive equilibria). Lyapunov (energy) functions, when available, provide constructive certificates of convergence.

%
% Optional figures (commented to keep the main text compact for review)
% \begin{figure*}[!t]
%   \begin{minipage}{0.49\linewidth}
%     \centering
%     \includegraphics[width=\linewidth]{neuron_small.png}
%     \caption{Schematic of biological neurons and synaptic connections.}
%     \label{fig:neuron}
%   \end{minipage}\hfill
%   \begin{minipage}{0.49\linewidth}
%     \centering
%     \includegraphics[width=0.75\linewidth]{circuit_mod.png}
%     \caption{Equivalent Kirchhoff circuit for a Hopfield network node.}
%     \label{fig:circuit}
%   \end{minipage}
% \end{figure*}

% \subsection{Filippov Solutions}

To define the solutions of system~\eqref{eq:special} with 
a discontinuous right-hand side and $g\in\mathcal{G'}$,
we employ A.~Filippov's theory~\cite{p2}.
The concept of Filippov solutions is that the tangent vector
to the trajectory of the solution, where it exists, should lie
in the closure of the convex hull of the limiting values
of the vector field $f(x(t))$ in arbitrarily small neighborhoods of the trajectory. 
Thus, system \eqref{eq:special} becomes a system of \emph{differential inclusions}.

\smallskip 
\begin{definition}[solution to the Cauchy problem by Filippov]\label{fillipovsol}
\textit{A solution to the Cauchy problem} for the system of differential equations \eqref{eq:special} on the interval $[t_0,t_1]$, where $t_0\leqslant t_1 \leqslant +\infty$, with the initial condition $x(t_0) = x_0$, is a vector-valued function $x(t)$ defined on $[t_0,t_1]$, which is absolutely continuous on this interval, satisfies the initial conditions, and almost everywhere on $[t_0,t_1]$ satisfies the differential inclusion:
\[  
    \dot x (t)\in \varphi\big(f(x(t))\big) = Bx(t)+T~\overline{\operatorname{conv}}\big[g(x(t))\big]+I,
    \label{1.2}
\]
\[  
    \varphi\big(f(x(t))\big)= \bigcap\limits_{\delta>0}\bigcap\limits_{\mu(\mathcal{N})=0}\overline{\operatorname{conv}}\big[f(\mathcal{B}(x(t),\delta) - \mathcal{N})\big], 
\]
where $\varphi\big(f(x(t))\big)$ is the smallest convex 
closed set containing $f(x'(t))$, as $x'$ 
varies over almost the entire $\delta$-neighborhood 
of the point $x$ for a fixed $t$ and 
$\overline{\operatorname{conv}}\{\cdot\}$ is a closure of the convex hull of the limit values at a point. 
The above can be rewritten~\cite{p9} as:
\begin{equation}\label{eq:special_conv}
\begin{aligned}
&\dot x (t)\in \varphi\big(f(x(t))\big) = \\
	&=\overline{\operatorname{conv}}\big\{\lim f(x_i(t)) ~|~x_i(t) \to x(t),~x_i(t)\notin N_f\big\} = \\
	&=Bx(t)+T~\overline{\operatorname{conv}}[g(x(t))]+I. 
\end{aligned}
\end{equation}

 It is important to note that we exclude the null measure sets $N_f$,
 where the function has a discontinuity.
 This allows us to define trajectories at points where
 the vector field is not defined.
 Using Filippov's definition of solutions, we can also define the \emph{equilibrium point}
 of system~\eqref{eq:special} as a stationary solution
 of the differential inclusion on the semi-interval $[0,+\infty)$:
 \begin{equation*} 0\in\varphi(e) =
 Be+T~\overline{\operatorname{conv}}\big(g(e)\big)+I.
 \end{equation*}    
\end{definition}

Now we define what is meant by an equilibrium point of the network and its global asymptotic stability in this theory~\cite{p1}.


\begin{definition}[Equilibrium point]\label{equilib}
An \textit{equilibrium point} $e\in \mathbb{R}^n$ of system \eqref{eq:special} is a constant solution $x(t)=e$ defined on $[0,+\infty)$ that satisfies the inclusion:
\[      
    0\in\varphi(e) = 
    Be+T~\overline{\operatorname{conv}}\big(g(e)\big)+I.
\]
If $e$ is an \textit{equilibrium point}, then there exists a vector $\eta \in \overline{\operatorname{conv}}\big(g(e)\big)$ such that the following equality holds:
\[ 0=Be+T\eta+I \tag{a.e.}.\]
Hereafter, $\eta \in \mathbb{R}^n$ will be called the \textit{output equilibrium point} of the network.
\end{definition}

\begin{definition}[Global asymptotic stability]\label{gas_of_eq}
An equilibrium point $e \in \mathbb{R}^n$ of network \eqref{eq:special} is \textit{globally asymptotically stable} if it is locally Lyapunov stable \cite{Lyapunov}, and for any $ x_0\in \mathbb{R}^n$, the solution $x(t)$ of system \eqref{eq:special} defined on $[0,+\infty)$ with initial data $x(0)=x_0$ converges to $e$ as $t \to +\infty$:

   \[  \lim\limits_{t\to + \infty }x\big(t\big) = e.\]
\end{definition}

In the end of this section lets also give a definition of \textit{upper semicontinuous set-valued mapping}. This concept is later used to connect inclusion \eqref{eq:special_conv} with Filippov's theory. 

\begin{definition}
    Set-valued map $F(p)$ is called \textit{upper semicontinuous} in $p_0$, if for any open $V\subset\mathbb{R}^n$, $F(p_0)\subset V$, there exists a neighborhood $U(p_0)$ of $p_0$ such as:
\[ 
F(p) ~\subset ~ V, \quad \forall p\in U(p_0).
\]
\end{definition}

% \medskip
% \textbf{The HNN as a dynamical system.}
% When all activations $g_i$ are continuous, \eqref{eq:special} is an
% ordinary differential equation and defines a classical semiflow
% $\{S_t\}_{t\ge 0}$ on $\mathbb{R}^n$: for each $x_0\in\mathbb{R}^n$
% there exists a unique trajectory
% $t\mapsto x(t;x_0)$ solving $\dot x = -Bx + T g(x) + I$ with $x(0)=x_0$.
% Equilibria are zeros of the vector field $f(x)=-Bx+Tg(x)+I$, i.e.,
% solutions of
% \[
%    -Bx^\star + T g(x^\star) + I = 0,
% \]
% and their stability properties determine the qualitative behavior of the
% network.

% In the discontinuous case $g\in\mathcal{G'}$ we interpret
% \eqref{eq:special} in the sense of Filippov: the right–hand side is
% replaced by its convexified set–valued map
% \[
%    F(x) := -Bx + T\,\overline{\operatorname{conv}}\,g(x) + I,
% \]
% and solutions are absolutely continuous curves $x(\cdot)$ satisfying
% $\dot x(t)\in F(x(t))$ almost everywhere~\cite{Filippov1988}.
% Under Assumptions~C1–C3 the multifunction $F$ is upper semicontinuous
% with nonempty compact convex values and linear growth, so that the
% standard existence and continuation theory for differential inclusions
% applies~\cite{AubinCellina1984,Filippov1988}.
% For each $t\ge 0$ and $x_0\in\mathbb{R}^n$ we denote by
% \begin{multline*}   
%    \Phi_t(x_0) := \\
%    \bigl\{\,x(t)\;\big|\;x(\cdot)\ \text{is a Filippov solution of }
%    \dot x\in F(x),\ x(0)=x_0 \bigr\}
% \end{multline*}
% the set of all states reachable at time $t$ from $x_0$.
% The family $\{\Phi_t\}_{t\ge 0}$ is then an upper–semicontinuous
% multivalued semiflow (or generalized semiflow, in the sense of
% Ball~\cite{Ball1997}) on $\mathbb{R}^n$ with nonempty compact values;
% equivalently, it is a set–valued dynamical system in the terminology of
% Bena\"{\i}m–Hofbauer–Sorin~\cite{BenaimHofbauerSorin2005}.

% An equilibrium (in the Filippov sense) is a point
% $e\in\mathbb{R}^n$ such that $0\in F(e)$.
% Equivalently, there exists an ``output equilibrium''
% $\eta^\star\in\overline{\operatorname{conv}}\,g(e)$ with
% \(
%    0 = -Be + T\eta^\star + I.
% \)
% The global asymptotic stability result proved below shows that, under
% Assumptions~C1–C3, this equilibrium is unique and attracts all Filippov
% trajectories, so the multivalued semiflow $\{\Phi_t\}$ is effectively
% gradient–like despite possible forward nonuniqueness.
% In what follows, we use the term ``dynamical system'' for the HNN in
% this generalized (set–valued) sense.


%%%%%%%%%%%%%%%%%%%%%%%%%%%%%%%%%%%%%%%%%%%%%%%%%%%%%%%%%%%%%%%%%%%%%%%%
\subsection{Stationary Sets and Gradient--Like Behavior}
\label{subsec:YLG-terms}

In this subsection we collect the basic notions from~\cite{YLG} that are
used in the sequel.  Consider an autonomous differential inclusion
\begin{equation}
  \dot x(t) \in f(x(t)), \qquad x(t)\in\mathbb{R}^n,
  \label{eq:aut-incl-prelim}
\end{equation}
with (Filippov) solutions defined on forward intervals $[t_0,+\infty)$.

\begin{definition}[Distance to a set]
For a nonempty set $\Lambda\subset\mathbb{R}^n$ and a point $x\in\mathbb{R}^n$
we denote the distance from $x$ to $\Lambda$ by
\[
  p(x,\Lambda) := \inf_{y\in\Lambda} \|x-y\|.
\]
For a trajectory $t\mapsto x(t)$ we write
$p[x(t),\Lambda] := p\bigl(x(t),\Lambda\bigr)$.
\end{definition}

\begin{definition}[Stationary vectors and stationary set {\cite[Def.~1.7]{YLG}}]
A vector $c\in\mathbb{R}^n$ is called \emph{stationary} for
system~\eqref{eq:aut-incl-prelim} if the constant function
$x(t)\equiv c$ is a solution of~\eqref{eq:aut-incl-prelim}.
The set of all stationary vectors is called the \emph{stationary set}
of~\eqref{eq:aut-incl-prelim} and is denoted by
\[
  \Lambda := \{\,c\in\mathbb{R}^n \mid 0\in f(c)\,\}.
\]
\end{definition}

\begin{definition}[Dichotomy and pointwise dichotomy {\cite[Defs.~1.8--1.9]{YLG}}]
System~\eqref{eq:aut-incl-prelim} is called
\emph{dichotomic} if every solution $x(t)$ which is bounded for $t>0$
tends to the stationary set as $t\to+\infty$, that is,
\[
   \lim_{t\to+\infty} p[x(t),\Lambda] = 0
   \quad\text{whenever } \sup_{t\ge 0}\|x(t)\|<\infty.
\]
It is called \emph{pointwise dichotomic} if every bounded solution
tends to a (possibly solution–dependent) stationary vector,
i.e., for any bounded solution there exists $c\in\Lambda$ such that
\[
   \lim_{t\to+\infty} x(t) = c.
\]
\end{definition}

\begin{definition}[Gradient--like and quasi–gradient--like behavior
{\cite[Defs.~1.10--1.11]{YLG}}]
System~\eqref{eq:aut-incl-prelim} is called
\emph{gradient–like} if for every solution $x(t)$ there exists an
equilibrium $c\in\Lambda$ such that
\[
   \lim_{t\to+\infty} x(t) = c.
\]
It is called \emph{quasi–gradient–like} if every solution tends to the
stationary set in the sense that
\[
   \lim_{t\to+\infty} p[x(t),\Lambda] = 0
   \quad\text{for all solutions } x(\cdot).
\]
\end{definition}

\begin{definition}[Lyapunov stability of a stationary set {\cite[Def.~1.12]{YLG}}]
The stationary set $\Lambda$ of~\eqref{eq:aut-incl-prelim} is said to be
\emph{Lyapunov stable} (stable in the small) if for every $\varepsilon>0$
there exists $\delta>0$ such that for any solution $x(t)$ and any
$t_0\in\mathbb{R}$, the condition
$p[x(t_0),\Lambda]<\delta$ implies
\[
   p[x(t),\Lambda] < \varepsilon
   \quad\text{for all } t\ge t_0.
\]
\end{definition}

\begin{definition}[Global and pointwise global stability
{\cite[Defs.~1.13--1.14]{YLG}}]
A stationary set $\Lambda$ is called \emph{globally stable} if it is
Lyapunov stable and system~\eqref{eq:aut-incl-prelim} is
quasi–gradient–like.

It is called \emph{pointwise globally stable} if it is globally stable
and, in addition, every solution $x(t)$ converges to a stationary
vector, i.e., for each solution there exists $c\in\Lambda$ such that
\[
   \lim_{t\to+\infty} x(t) = c.
\]
\end{definition}

\begin{remark}
When the stationary set $\Lambda$ consists of a single equilibrium
$\{e\}$, pointwise global stability of $\Lambda$ is equivalent to the
usual \emph{global asymptotic stability} of $e$ in the sense of
Definition~\ref{gas_of_eq}.
\end{remark}

%%%%%%%%%%%%%%%%%%%%%%%%%%%%%%%%%%%%%%%%%%%%%%%%%%%%%%%%%%%%%%%%%%%%%%%%
\subsection{Extended Nonlinearities in Lur'e Systems}
\label{subsec:extended-nonlinearity}

For Lur'e systems of the form
\begin{equation}
  \dot x(t) = P x(t) + q\,\xi(t), \qquad
  \sigma(t) = r^{\top}x(t),
  \label{eq:lure-extended}
\end{equation}
with a (possibly discontinuous) nonlinearity
$\varphi \!:\!\mathbb{R}^m \! \times\mathbb{R}\rightrightarrows\mathbb{R}^m$,
Yakubovich--Leonov--Gelig~\cite{YLG} use the notion of an \emph{extended
nonlinearity}.

\begin{definition}[Extended nonlinearity {\cite[Sec.~1.2]{YLG}}]
A pair of functions $(x(\cdot),\xi(\cdot))$ defined on an interval
$[t_0,t_1]$ is called a solution of the Lur'e system
\eqref{eq:lure-extended} if
\begin{itemize}
  \item $x(t)$ is absolutely continuous and $\xi(t)$ is integrable;
  \item the relations
  \[
     \dot x(t) = P x(t) + q\,\xi(t), \quad
     \xi(t) \in \varphi\bigl(\sigma(t),t\bigr),\ \
     \sigma(t)=r^{\top}x(t)
  \]
  hold for almost all $t\in[t_0,t_1]$.
\end{itemize}
The function $\xi(t)$ is then called an \emph{extended nonlinearity}
associated with the (possibly multivalued) nonlinearity
$\varphi(\sigma,t)$ along the trajectory $x(t)$.
\end{definition}

In the Hopfield-network setting, after Filippov convexification of the
activation map $g$, the role of $\xi(t)$ will be played by a measurable
\emph{extended activation} $\eta(t)\in\overline{\operatorname{conv}}\,g(x(t))$
constructed via Makarov's measurable selection theorem; 
see Theorem~\ref{thm:makarov} in Appendix~\ref{app:YLG} for details.

% \newpage

%%%%%%%%%%%%%%%%%%%%%%%%%%%%%%%%%%%%%%%%%%%%%%%%%%%%%%%%%%%%%%%%%%%%%%%%
\subsection{The HNN as a dynamical system}\label{subsec:HNN-dyn-sys}

\todo[inline]{New below}
In the classical setting with continuous activations $g_i$, the
Hopfield network \eqref{eq:special} is an ordinary differential
equation
\[
   \dot x = -Bx + T g(x) + I, \qquad x\in\mathbb{R}^n,
\]
and the Cauchy problem is well posed. For each $x_0\in\mathbb{R}^n$
there exists a unique trajectory $t\mapsto x(t;x_0)$, and the maps
$S_t:x_0\mapsto x(t;x_0)$, $t\ge0$, form a (single–valued) semiflow on
$\mathbb{R}^n$ in the usual sense of dynamical systems theory.

In the discontinuous case $g\in\mathcal{G'}$ we work with the Filippov
inclusion~\eqref{eq:special_conv}. For brevity we denote the Filippov
regularization of the right–hand side by
\begin{equation}\label{eq:F-filippov}
   F(x)
   := \varphi\bigl(f(x)\bigr)
   = Bx + T\,\overline{\operatorname{conv}}[g(x)] + I,
   \qquad x\in\mathbb{R}^n.
\end{equation}
Filippov solutions are absolutely continuous functions $x(\cdot)$
satisfying $\dot x(t)\in F(x(t))$ almost everywhere
(Definition~\ref{fillipovsol}). In general, such solutions are not
unique for a given initial state, and therefore \eqref{eq:special} no
longer generates a single–valued semiflow. Nevertheless, it fits
naturally into several generalized dynamical–systems frameworks.

\smallskip
\paragraph*{Set–valued dynamical system (multivalued semiflow)}
For each $t\!\!\ge\!\!0$ and $x_0\!\in\!\mathbb{R}^n$ we define the reachable set
at time $t$~by
\begin{multline}\label{eq:Phi-t}
   \Phi_t(x_0)
   := \bigl\{\,x(t)\;\big|\;x(\cdot)\ \text{is a Filippov solution} \\ 
   \text{of } \dot x\in F(x),\ x(0)=x_0 \bigr\}.
\end{multline}
Under the basic hypotheses of Filippov theory (upper semicontinuity of
$F$, nonempty compact convex values, and linear growth), which are
satisfied for~\eqref{eq:F-filippov} in view of
Assumptions~C1–C2, every initial condition admits at least one global
solution on $[0,+\infty)$ and the family
$\{\Phi_t\}_{t\ge0}$ enjoys the properties
\begin{itemize}
  \item $\Phi_0(x)=\{x\}$ for all $x\in\mathbb{R}^n$;
  \item $\Phi_t(x)\neq\varnothing$ and $\Phi_t(x)$ is compact for all
        $t\ge0$, $x\in\mathbb{R}^n$;
  \item \emph{semigroup property}:
        $\Phi_{t+s}(x)=\Phi_t\bigl(\Phi_s(x)\bigr)$
        (composition of sets) for all $s,t\ge0$;
  \item the map $x\mapsto\Phi_t(x)$ is upper semicontinuous for each
        $t\ge0$.
\end{itemize}
In the terminology of differential inclusions
(e.g.~\cite{Filippov1988,AubinCellina1984,BenaimHofbauerSorin2005}),
$\{\Phi_t\}_{t\ge0}$ is an \emph{upper–semicontinuous multivalued
semiflow}, or equivalently a \emph{set–valued dynamical system} on
$\mathbb{R}^n$. All stability notions used in this paper (equilibria,
global attractivity) are understood with respect to this multivalued
evolution.

\smallskip
\paragraph*{Generalized semiflow viewpoint}
An alternative but equivalent description is given by the notion of a
\emph{generalized semiflow} in the sense of Ball~\cite{Ball1997}.  One
considers the family $\mathcal{G}$ of all Filippov trajectories
$x:\mathbb{R}_+\to\mathbb{R}^n$ of~\eqref{eq:special_conv} and checks
that $\mathcal{G}$ is closed under time shifts and concatenation and is
upper semicontinuous with respect to initial data.  Then
$\mathcal{G}$ is a generalized semiflow on $\mathbb{R}^n$, and our main
result can be rephrased as saying that, under Assumptions~C1–C3, the
singleton stationary set $\{e\}$ is the global attractor of this
generalized semiflow.

\smallskip
\paragraph*{Skew–product\,/\,cocycle representation}
Finally, it is often convenient—especially in the YLG framework—to
represent Filippov trajectories as solutions of an ordinary
differential equation with a measurable ``extended activation''.
By measurable–selection results due to Makarov and
Yakubovich–Leonov–Gelig (see Theorem~\ref{thm:makarov} in Appendix~\ref{app:YLG} 
for precise statements), for any Filippov solution $x(\cdot)$ there exists a
measurable function
$\eta:\mathbb{R}_+\to\mathbb{R}^n$ such that
\[
   \eta(t) \in \overline{\operatorname{conv}}[g(x(t))] \quad\text{a.e.},
   \qquad
   \dot x(t) = Bx(t) + T\eta(t) + I \quad\text{a.e.}
\]
Thus the HNN can be viewed as a control system
\(
   \dot x = Bx + T\eta + I
\)
driven by an admissible input $\eta(\cdot)$ taking values in the convex
set $\overline{\operatorname{conv}}[g(x)]$.  If $\Omega$ denotes the
space of admissible input functions endowed with the shift
$\theta_t\omega(\cdot):=\omega(t+\cdot)$, then the pair
$(\theta_t,\varphi)$, where $\varphi(t,\omega,x_0)$ is the set of
solutions at time $t$ with input $\omega$ and initial state $x_0$, forms
a (possibly multivalued) \emph{skew–product dynamical system} (or
cocycle) over the base flow $(\Omega,\theta_t)$ in the standard sense
of dynamical–systems and random–dynamical–systems theory (see e.g.~\cite{}).
This representation connects the Filippov HNN with the Lur'e–type
frequency–domain analysis used in the sequel, while the multivalued
semiflow viewpoint above provides the natural setting for our global
stability statements.

\subsection{A YLG Theorem on Quasi-Gradient-Like Behavior and Pointwise Global Stability}
\label{sec:ylg-theorem}

The global-stability result proved later for discontinuous Hopfield neural
networks relies on a Yakubovich--Leonov--Gelig theorem for systems with several
discontinuous nonlinearities. Since the original statement is formulated in the
frequency domain for a general Lur'e-type system, we recall here a faithful
version in notation convenient for the sequel.

Consider the system
\begin{equation}
  \dot{x} = P x + q\,\xi, 
  \qquad
  \sigma = r^{\top}x,
  \qquad
  \xi_i \in \varphi_i(\sigma_i),
  \quad i=1,\dots,m,
  \label{eq:ylg-general-system}
\end{equation}
where \(x\in\mathbb{R}^n\), \(\sigma,\xi\in\mathbb{R}^m\),
\(P\in\mathbb{R}^{n\times n}\), \(q\in\mathbb{R}^{n\times m}\),
\(r\in\mathbb{R}^{n\times m}\), and each \(\varphi_i\) is piecewise
single-valued. Let
\begin{equation}
  \chi(\lambda) := r^{\top}(P-\lambda I)^{-1}q
  \label{eq:ylg-transfer}
\end{equation}
be the transfer matrix of \eqref{eq:ylg-general-system}, and let
\(\Lambda\) denote its stationary set.

\begin{theorem}[Adapted form of {\cite[Theorem~3.1]{YLG}}]
\label{thm:YLG-adapted-faithful}
Assume that the following conditions hold.

\begin{enumerate}
\item[(i)]
The transfer matrix \(\chi(\lambda)\) in \eqref{eq:ylg-transfer} is
irreducible.

\item[(ii)]
There exist diagonal matrices
\begin{multline*}
  T=\diag(\tau_1,\dots,\tau_m)>0,\\
  M=\diag(\mu_1,\dots,\mu_m)\ge 0,\\
  \Theta=\diag(\vartheta_1,\dots,\vartheta_m),
\end{multline*}   
such that
\begin{equation}
  TM + \Re\!\bigl[(T+i\omega \Theta)\chi(i\omega)\bigr] \ge 0,
  \qquad \forall\,\omega\in\mathbb{R},
  \label{eq:ylg-freq-ineq}
\end{equation}
and
\begin{equation}
  \det(T+\lambda \Theta)\neq 0
  \qquad
  \text{whenever } \det(P-\lambda I)=0.
  \label{eq:ylg-det-cond}
\end{equation}

\item[(iii)]
For each \(i=1,\dots,m\), the nonlinearity \(\varphi_i\) satisfies
\begin{equation}
  (\sigma_i-\mu_i\xi_i)\,\xi_i \ge 0,
  \qquad
  \forall\,\xi_i\in\varphi_i(\sigma_i),
  \label{eq:ylg-sector}
\end{equation}
\begin{equation}
  \sigma_i \notin \mu_i\,\varphi_i(\sigma_i),
  \qquad
  \sigma_i\neq 0,
  \label{eq:ylg-nondeg}
\end{equation}
and, in addition, whenever \(\mu_i\neq 0\) and \(\vartheta_i<0\),
\begin{equation}
  \int\limits_{0}^{\pm\infty}
  \bigl|\lambda-\mu_i\varphi_i(\lambda)\bigr|\,d\lambda
  =\infty.
  \label{eq:ylg-integral}
\end{equation}

\item[(iv)]
There exist numbers \(\mu_i^0>\mu_i\), \(i=1,\dots,m\), such that, with
\[
  M_0:=\diag(\mu_1^0,\dots,\mu_m^0),
\]
the matrix
\begin{equation}
  P + q M_0^{-1} r^{\top}
  \label{eq:ylg-hurwitz}
\end{equation}
is Hurwitz.

\item[(v)]
The stationary set \(\Lambda\) of \eqref{eq:ylg-general-system} is bounded.
Moreover, if a solution \(x(t),\xi(t)\) of \eqref{eq:ylg-general-system}
satisfies
\begin{equation}
  \xi_i(t)\sigma_i(t)\equiv 0,
  \qquad i=1,\dots,m,
  \label{eq:ylg-orthogonality}
\end{equation}
and
\begin{equation}
  x^{\top}(t) H x(t)\equiv \mathrm{const}
  \label{eq:ylg-quadratic-const}
\end{equation}
for all \(t\in\mathbb{R}\), where \(H\) is some positive definite matrix,
then
\[
  x(t)\equiv x^0\in\mathbb{R}^n,\qquad
  \xi(t)\equiv \xi^0\in\mathbb{R}^m,
\]
and
\[
  Px^0 + q\xi^0 = 0,
\]
that is, the solution is stationary.
\end{enumerate}

Then system \eqref{eq:ylg-general-system} is quasi-gradient-like.

Assume, in addition, that:

\begin{enumerate}
\item[(vi)]
There exists \(\varepsilon>0\) such that if \(\rho(x,\Lambda)<\varepsilon\),
then for any stationary solution \(x^0,\xi^0\) the inequalities
\begin{equation}
  (\xi_i-\xi_i^0)
  \bigl[\sigma_i-\sigma_i^0-\mu_i(\xi_i-\xi_i^0)\bigr]\ge 0,
  \qquad
  \forall\,\xi_i\in\varphi_i(\sigma_i),
  \label{eq:ylg-local-mono}
\end{equation}
hold for \(i=1,\dots,m\), where \(\sigma_i\) and \(\sigma_i^0\) are the
\(i\)th components of \(r^{\top}x\) and \(r^{\top}x^0\), respectively.
In the case \(\vartheta_i<0\), the bracketed term in
\eqref{eq:ylg-local-mono} is additionally assumed to be nonzero whenever
\(\xi_i\neq \xi_i^0\) and \(x\neq x^0\).
\end{enumerate}

Then the stationary set \(\Lambda\) of
\eqref{eq:ylg-general-system} is pointwise globally stable.
\end{theorem}

\begin{remark}
Theorem~\ref{thm:YLG-adapted-faithful} is stated in the original
frequency-domain Lur'e notation because the conditions
\eqref{eq:ylg-freq-ineq}--\eqref{eq:ylg-hurwitz} are most naturally expressed
through the transfer matrix \(\chi(\lambda)\). In the Hopfield setting studied
below, the theorem is applied after rewriting the discontinuous HNN in a
compatible Lur'e form and identifying the corresponding matrices
\(P\), \(q\), \(r\), and nonlinearities \(\varphi_i\).
\end{remark}

\begin{corollary}
\label{cor:YLG-singleton}
If, under the assumptions of
Theorem~\ref{thm:YLG-adapted-faithful}, the stationary set is a singleton,
\[
  \Lambda=\{x^\ast\},
\]
then \(x^\ast\) is a globally asymptotically stable equilibrium of
\eqref{eq:ylg-general-system}.
\end{corollary}

\begin{remark}
Corollary~\ref{cor:YLG-singleton} is the form used later in the proof of the
main global-stability theorem for discontinuous Hopfield neural networks:
once Assumptions~\textbf{C1--C3} are verified and the stationary set is shown to consist
of a single point, global asymptotic stability follows by direct invocation of
Theorem~\ref{thm:YLG-adapted-faithful}.
\end{remark}

%%%%%%%%%%%%%%%%%%%%%%%%%%%%%%%%%%%%%%%%%%%%%%%%%%%%%%%%%%%%%%%%%%%%%%%%%%%%%%%
%%%%%%%%%%%%%%%%%%%%%%%%%%%%%%%%%%%%%%%%%%%%%%%%%%%%%%%%%%%%%%%%%%%%%%%%%%%%%%%
% \section{Problem Statement and Assumptions}\label{sec:problem}
% Consider the HNN in Lur\'e form
% \begin{equation}
% \dot{x} = -Bx + T\,g(x) + I, \qquad x\in\mathbb{R}^n,
% \end{equation}
% with $B=\mathrm{diag}(b_1,\dots,b_n)\succ 0$, interconnection $T\in\mathbb{R}^{n\times n}$, input $I\in\mathbb{R}^n$, and componentwise activation $g(x)=(g_1(x_1),\ldots,g_n(x_n))^\top$ that may be discontinuous.

% \begin{assumption}[C1: Linear growth]
% There exist $a_i,b_i\ge 0$ such that $|g_i(\xi)|\le a_i|\xi|+b_i$ for all $\xi\in\mathbb{R}$ and each $i$.
% \end{assumption}

% \begin{assumption}[C2: One-sided Lipschitz (componentwise)]
% For each $i$, there exists $C_i\in\mathbb{R}$ such that for all $x,y\in\mathbb{R}^n$ and a.e.\ selections $\eta\in\mathrm{co}\,g(x)$, $\zeta\in\mathrm{co}\,g(y)$,
% \(
% (\eta_i-\zeta_i)\,(x_i-y_i)\le C_i\,(x_i-y_i)^2.
% \)
% \end{assumption}

% \begin{assumption}[C3: Diagonal Lyapunov stability]
% There exists a diagonal $H=\mathrm{diag}(h_1,\dots,h_n)\succ 0$ such that the symmetric part of $-HB$ dominates the interconnection, e.g.,
% \(
% -HB-(HB)^{\top} + \alpha\, (T^{\top}H + HT) \prec 0
% \)
% for some admissible $\alpha\geqslant 0$.
% \end{assumption}
% \begin{remark}
% Condition \textbf{C2} is also called \textit{unilateral Lipschitz-like condition} and along with linear growth condition \textbf{C1} is used in similar manner in several papers \cite{unilateral_Lipschitz_like_condition, eq_ref_11, GAS_of_delayed, GAS_discrete_time_varying, periodic_time_varying} with which this paper mostly shares its metathetical apparatus.
% \end{remark}

% ================================================================
% Problem statement and assumptions
% ================================================================
\section{Problem Statement and Assumptions}\label{sec:problem}

Consider the Hopfield neural network (HNN) in Lur\'e form
\begin{equation}\label{eq:HNN-lure}
  \dot{x}(t) = -Bx(t) + T\,g(x(t)) + I, 
  \qquad x(t)\in\mathbb{R}^n,
\end{equation}
where $B=\mathrm{diag}(b_1,\dots,b_n)\succ 0$ collects effective leak
rates, $T\in\mathbb{R}^{n\times n}$ is the interconnection (synaptic)
matrix, $I\in\mathbb{R}^n$ is the (rescaled) input vector, and
$g(x)=(g_1(x_1),\ldots,g_n(x_n))^\top$ is the componentwise activation
mapping, each $g_i:\mathbb{R}\to\mathbb{R}$ possibly discontinuous.

Throughout the paper we assume that $g$ belongs to the extended class
$\mathcal G'$ of activation functions (Definition~1), so that every
coordinate $g_i$ is piecewise continuous with a finite number of
first–kind discontinuities on any compact set, admits one–sided limits
at the discontinuity points, and is sign–preserving and
nondecreasing.  In particular, its Filippov convexification
$\mathrm{conv}\,g_i$ is a monotone, piecewise single–valued
multifunction in the sense of YLG.

\begin{assumption}[\bf C1: Linear growth]\label{ass:C1}
For each $i\in\{1,\dots,n\}$ there exist constants $a_i,b_i\ge 0$ such
that
\[
  |g_i(\xi)| \le a_i|\xi| + b_i,
  \qquad \forall\,\xi\in\mathbb{R}.
\]
\end{assumption}

\begin{assumption}[\bf C2: One-sided Lipschitz (componentwise)]\label{ass:C2}
For each $i\in\{1,\dots,n\}$ there exists $C_i\in\mathbb{R}$ such that
for all $x,y\in\mathbb{R}^n$ and for all a.e.\ selections
$\eta\in\mathrm{conv}\,g(x)$, $\zeta\in\mathrm{conv}\,g(y)$ one has
\[
  (\eta_i - \zeta_i)\,(x_i - y_i)
  \;\le\; C_i\,(x_i - y_i)^2.
\]
\end{assumption}

\begin{remark}
Assumption~\textbf{C2} is also known as a \emph{unilateral
Lipschitz-like} condition and, together with the linear growth
condition~\textbf{C1}, it is used in a similar way in several works
on Hopfield-type networks with discontinuous activations,
see e.g.\ \cite{periodic_time_varying,eq_ref_11,unilateral_Lipschitz_like_condition,
GAS_of_delayed,GAS_discrete_time_varying}.
\end{remark}

\subsection*{Lur\'e representation and transfer matrix}

For the purposes of applying the YLG theory we write
\eqref{eq:HNN-lure} in the standard Lur\'e form
\begin{equation}\label{eq:lure-Pqr}
  \dot x = P x + q\,\xi, 
  \qquad \sigma = r^{\top}x, 
  \qquad \xi \in \varphi(\sigma),
\end{equation}
with
\begin{equation}\label{eq:Pqr-def}
  P := -B,\qquad
  q := T,\qquad
  r := I_n,
\end{equation}
and the extended nonlinearity
\[
  \varphi(\sigma)
  := \mathrm{conv}\,g(\sigma)
  = \bigl(\mathrm{conv}\,g_1(\sigma_1),\dots,
          \mathrm{conv}\,g_n(\sigma_n)\bigr)^\top .
\]
The transfer matrix of the linear part is
\begin{equation}\label{eq:transfer-X}
  X(\lambda) 
  := r^{\top}(P-\lambda I)^{-1}q
  = (-B-\lambda I)^{-1}T,
  \qquad \lambda\in\mathbb{C}.
\end{equation}

% ------------------------------------------------
% Main third assumption: frequency-domain condition
% ------------------------------------------------
\begin{assumption}[\bf C3: YLG-type frequency condition]\label{ass:C3prime}
The transfer matrix $X(\lambda)$ in \eqref{eq:transfer-X} is
irreducible in the sense of Definition~2.6 in Chapter~2 of YLG
(i.e.\ the pair $(P,r)=(-B,I_n)$ is observable), and there exist
diagonal matrices
\begin{multline*}   
   \mathcal{T}=\mathrm{diag}(\tau_1,\dots,\tau_n)\succ 0,\\
   \mathcal{M}=\mathrm{diag}(\mu_1,\dots,\mu_n)\succeq 0,\\
   \Theta=\mathrm{diag}(\vartheta_1,\dots,\vartheta_n)
\end{multline*}
such that
\begin{equation}\label{eq:C3prime-freq}
   \mathcal{T}\mathcal{M}
   + \operatorname{Re}\Bigl[
        \bigl(\mathcal{T} + i\omega \Theta\bigr)\,
        X(i\omega)
      \Bigr]
   \succeq 0,
   \qquad \forall\,\omega\in\mathbb{R},
\end{equation}
and
\begin{equation}\label{eq:C3prime-det}
  \det\bigl(\mathcal{T} + \lambda \Theta\bigr) \neq 0
  \quad\text{whenever}\quad 
  \det(P-\lambda I)=\det(-B-\lambda I)=0.
\end{equation}
\end{assumption}

Assumption~\textbf{C3} is exactly the frequency-domain hypothesis
\textup{(i)}–\textup{(ii)} of Theorem~3.1 in Chapter~3 of YLG,
specialized to the triple $(P,q,r)=(-B,T,I_n)$ of the Hopfield model.

% ------------------------------------------------
% Alternative algebraic formulation (derived from C3)
% ------------------------------------------------

\begin{assumption}[\bf C3$'$: Diagonal Lyapunov stability (DLS)]\label{ass:C3}
There exists a diagonal matrix $H=\mathrm{diag}(h_1,\dots,h_n)\succ 0$
and a constant $\alpha\ge 0$ such that
\begin{equation}\label{eq:C3-matrix}
  -HB - B^{\top}H + \alpha\bigl(T^{\top}H + HT\bigr) \prec 0 .
\end{equation}
\end{assumption}

Assumption~\textbf{C3$'$} will be used only as a derived, purely
state–space condition.  It can be taken as an alternative to
\textbf{C3} if one prefers to avoid frequency-domain inequalities (see Section \ref{sec:nonunique} and Appendix \ref{app:proofs}).
For our main results, however, we assume \textbf{C1}, \textbf{C2},
\textbf{C3$'$} and $g\in\mathcal G'$; the next lemma shows that
\textbf{C3} then follows automatically.

\begin{lemma}[C3 implies C3$'$]\label{lem:C3prime-implies-C3}
Suppose that Assumption~\textbf{C3} holds. Then there exist a
diagonal matrix $H=\mathrm{diag}(h_1,\dots,h_n)\succ 0$ and a constant
$\alpha>0$ such that the matrix inequality~\eqref{eq:C3-matrix} is
satisfied; in particular, Assumption~\textbf{C3$'$} holds.
\end{lemma}

\begin{proof}
Under Assumption~\textbf{C3}, conditions \textup{(i)}–\textup{(ii)}
of Theorem~3.1 in Chapter~3 of YLG are satisfied for the system
\eqref{eq:lure-Pqr} with $(P,q,r)=(-B,T,I_n)$.  By the frequency
criterion in YLG (see Eqns.~(3.10)–(3.11) on pp.~114–115), there
exists a symmetric matrix $H_{\mathrm{YLG}}=H_{\mathrm{YLG}}^{\top}$
such that the quadratic form
\[
  W(x,\xi)
  := \bigl[2x^{\top}H_{\mathrm{YLG}} + \xi^{\top}\Theta r^{\top}\bigr]\,
     (P x + q\xi)
     + \xi^{\top}\mathcal{T}\,(r^{\top}x - \mathcal{M}\xi)
\]
satisfies
\begin{equation}\label{eq:W-leq0}
  W(x,\xi) \le 0,
  \qquad \forall\,x\in\mathbb{R}^n,\ \forall\,\xi\in\mathbb{R}^n.
\end{equation}
Moreover, by Eq.~(3.14) on p.~115 of YLG one can choose
$H_{\mathrm{YLG}}$ positive definite:
\begin{equation}\label{eq:H-YLG-pd}
  H_{\mathrm{YLG}} \succ 0 .
\end{equation}

Set $\xi\equiv 0$ in the definition of $W$ and use $r^{\top}=I_n$ and
$P=-B$:
\begin{multline*}
  W(x,0)
   = \bigl[2x^{\top}H_{\mathrm{YLG}}\bigr]\,(P x)
   = 2x^{\top}H_{\mathrm{YLG}}P x = \\
   = -2x^{\top}H_{\mathrm{YLG}}B x .    
\end{multline*}
Together with \eqref{eq:W-leq0}, this gives
\begin{equation}\label{eq:HB-negative}
  -2x^{\top}H_{\mathrm{YLG}}B x \le 0,
  \qquad \forall\,x\in\mathbb{R}^n.
\end{equation}
Since $B=\mathrm{diag}(b_1,\dots,b_n)\succ 0$ and
$H_{\mathrm{YLG}}\succ 0$, we have $x^{\top}H_{\mathrm{YLG}}Bx>0$ for
all $x\neq 0$.  Hence \eqref{eq:HB-negative} is in fact strict for
$x\neq0$, and the symmetric matrix
\[
  S_{\mathrm{YLG}}
  := -H_{\mathrm{YLG}}B - B^{\top}H_{\mathrm{YLG}}
\]
is negative definite:
\begin{equation}\label{eq:SYLG-negative}
  x^{\top}S_{\mathrm{YLG}}x
  = -2x^{\top}H_{\mathrm{YLG}}B x < 0,
  \qquad \forall\,x\neq 0 .
\end{equation}

Define the diagonal matrix
\begin{equation}\label{eq:H-diag-def}
  H := \mathrm{diag}\bigl((H_{\mathrm{YLG}})_{11},
                         \dots,
                         (H_{\mathrm{YLG}})_{nn}\bigr).
\end{equation}
By \eqref{eq:H-YLG-pd}, all diagonal entries $(H_{\mathrm{YLG}})_{ii}$
are strictly positive, so $H\succ 0$.  Since $B$ is diagonal and
positive definite, for every $x\neq 0$ we have
\begin{align*}
  x^{\top}\bigl(-HB - B^{\top}H\bigr) x
    &= -2x^{\top}HBx \\
    &= -2(B^{1/2}x)^{\top}\bigl(B^{1/2}HB^{1/2}\bigr)(B^{1/2}x) \\
    &< 0,
\end{align*}
because $B^{1/2}HB^{1/2}\succ 0$.  Thus
\begin{equation}\label{eq:S0-negative}
  S_0 := -HB - B^{\top}H \prec 0 .
\end{equation}

Let
\[
  S_1 := T^{\top}H + HT
\]
and consider the one-parameter family of symmetric matrices
\[
  S(\alpha) := S_0 + \alpha S_1,
  \qquad \alpha\ge 0.
\]
The eigenvalues of $S(\alpha)$ depend continuously on $\alpha$, and
$S(0)=S_0\prec 0$ by \eqref{eq:S0-negative}.  Hence there exists
$\alpha^\ast>0$ such that $S(\alpha)\prec 0$ for all
$\alpha\in[0,\alpha^\ast]$.  Choosing any
$\alpha\in(0,\alpha^\ast]$, we obtain
\[
  -HB - B^{\top}H + \alpha\bigl(T^{\top}H + HT\bigr)
  = S(\alpha) \prec 0 ,
\]
which is precisely~\eqref{eq:C3-matrix}.  This proves
Assumption~\textbf{C3$'$}.
\end{proof}

% \begin{remark}
% In the subsequent analysis we work under the main assumptions
% \textbf{C1}, \textbf{C2}, \textbf{C3$'$} and $g\in\mathcal G'$; the
% algebraic condition \textbf{C3} is used only via
% Lemma~\ref{lem:C3prime-implies-C3} as a convenient state–space
% consequence of~\textbf{C3$'$}.
% \end{remark}

% ================================================================
% Main results
% ================================================================
\section{Main Results}\label{sec:main}

The Hopfield network considered in Section~\ref{sec:problem} can be rewritten
in the Lur'e form \eqref{eq:ylg-general-system} after identifying the state,
the static discontinuous activation map, and the corresponding interconnection
operators. The main task in the proof of the theorem below is therefore to
verify conditions~(i)--(vi) of Theorem~\ref{thm:YLG-adapted-faithful} under
Assumptions~\textbf{C1--C3} and to prove that the associated stationary set is a
singleton. Once this is done, global asymptotic stability follows from
Corollary~\ref{cor:YLG-singleton}.

\begin{theorem}[Global asymptotic stability of a Hopfield network with discontinuous activations]\label{thm:main}
Consider the continuous–time Hopfield neural network
\begin{equation}
  \dot x(t) = -B x(t) + T\,g(x(t)) + I, \qquad x(t)\in\mathbb{R}^n,
  \label{eq:HNN-main}
\end{equation}
where
\begin{itemize}
  \item $B = \mathrm{diag}(b_1,\dots,b_n)\succ 0$ is a diagonal matrix
        of leak rates;
  \item $T\in\mathbb{R}^{n\times n}$ is the interconnection (synaptic) matrix;
  \item $I\in\mathbb{R}^n$ is a constant input vector;
  \item $g=(g_1,\dots,g_n)^{\top}:\mathbb{R}^n\to\mathbb{R}^n$ is the
        componentwise activation mapping, each
        $g_i:\mathbb{R}\to\mathbb{R}$ possibly discontinuous.
\end{itemize}
Assume that $g$ belongs to the extended class $\mathcal{G}'$ of
activation functions (Definition~1) and that Assumptions
\textbf{C1}–\textbf{C3} hold.

Interpret system~\eqref{eq:HNN-main} in the sense of Filippov, i.e.\
as the differential inclusion
\[
   \dot x(t) \in -B x(t) + T\,\varphi(x(t)) + I,
   \qquad \varphi(x)=\mathrm{conv}\,g(x).
\]
Then the following properties hold.
\begin{enumerate}
  \item For every initial condition $x(0)=x_0\in\mathbb{R}^n$ there
        exists at least one Filippov solution $x(\cdot;x_0)$
        of~\eqref{eq:HNN-main} defined on the whole half–line
        $[0,+\infty)$.

  \item There exists a unique equilibrium point $e\in\mathbb{R}^n$,
        i.e.\ a constant vector satisfying
        \[
           0 \in -B e + T\,\mathrm{conv}\,g(e) + I .
        \]

  \item This equilibrium is \emph{globally asymptotically stable in
        the sense of Filippov}: for every Filippov solution
        $x(\cdot;x_0)$ of~\eqref{eq:HNN-main} one has
        \[
           \lim_{t\to+\infty} x(t;x_0) = e ,
        \]
        and for every $\varepsilon>0$ there exists $\delta>0$ such
        that $|x_0-e|<\delta$ implies $|x(t;x_0)-e|<\varepsilon$ for
        all $t\ge0$.
\end{enumerate}
\end{theorem}

We treat the Hopfield network as a concrete Lur\'e–type system and
embed it into the YLG framework.  The argument proceeds in four main
steps:
\begin{enumerate}
  \item we establish global existence of Filippov solutions;
  \item we prove existence and uniqueness of an equilibrium;
  \item we rewrite the HNN explicitly in the Lur\'e form used by YLG;
  \item we verify that Assumptions~\textbf{C1}, \textbf{C2},
        \textbf{C3} and $g\in\mathcal G'$ imply all hypotheses of
        Theorem~3.1 in Chapter~3 of YLG, and then invoke that theorem
        to conclude global asymptotic stability of the equilibrium.
\end{enumerate}

\subsection{Global Existence of Filippov Solutions for the HNN}

In this subsection we show that under Assumption~\textbf{C1} every
Filippov solution of the Hopfield network exists on the entire
half–line $[0,+\infty)$.

Interpreting \eqref{eq:HNN-lure} in the sense of Filippov, we obtain
the differential inclusion
\begin{equation}
    \dot x(t) \in -B x(t) + T\,\varphi(x(t)) + I,
    \qquad x(t)\in\mathbb{R}^n,
    \label{eq:HNN-incl}
\end{equation}
where $\varphi(x) = \mathrm{conv}\,g(x)$ denotes the Filippov
convexification of the (possibly discontinuous) activation
mapping~$g:\mathbb{R}^n\to\mathbb{R}^n$.

\begin{proposition}[Global existence]
\label{prop:global-existence}
Suppose that Assumption~C1 (linear growth of the activation functions)
holds. Then, for every initial condition $x(0)=x_0\in\mathbb{R}^n$,
there exists at least one Filippov solution $x(\cdot;x_0)$
of~\eqref{eq:HNN-incl} defined on the whole half–line $[0,+\infty)$.
\end{proposition}

\begin{proof}
\emph{Step 1 (local existence).}
Fix $t_0\in[0,+\infty)$ and $a\in\mathbb{R}^n$, and consider the
cylinder
\[
    D := \{(x,t)\in\mathbb{R}^n\times\mathbb{R} \mid
           |t-t_0|\le\alpha,\ |x-a|\le\rho\},
\]
with some $\alpha>0$, $\rho>0$.
Since $g$ is piecewise continuous with a finite number of first–order
discontinuities on any compact set (Definition~1) and $\varphi$ is its
Filippov convexification, the multifunction
\[
   f(x,t) := -B x + T\,\varphi(x) + I
\]
is upper semicontinuous in $(x,t)$ and has nonempty, closed, bounded
and convex values (see, e.g.,~\cite[Ch.~1]{YLG}).
On the compact set $D$ the values of $f$ are uniformly bounded by some
constant $c>0$.

Therefore, the hypotheses of the local existence theorem for
differential inclusions (Theorem~1.1 in~\cite[Ch.~1]{YLG}) are
satisfied, and for any initial condition $x(t_0)=a$ there exists at
least one absolutely continuous solution of~\eqref{eq:HNN-incl} on
the interval $|t-t_0|\le\tau$, where
$\tau = \min\{\alpha,\rho/c\}$. Hence Filippov solutions of
\eqref{eq:HNN-incl} exist locally.

\emph{Step 2 (global growth estimate from \textbf{C1}).}
Assumption~\textbf{C1} states that for each scalar activation $g_i$ there exist
constants $a_i,b_i\ge0$ such that
\[
   |g_i(\xi)| \le a_i|\xi| + b_i, \qquad \forall\,\xi\in\mathbb{R}.
\]
It follows that there are constants $A,B\ge0$ with
\[
   \|\eta\| \le A\|x\| + B,
   \qquad \forall\,\eta\in\varphi(x),\ \forall x\in\mathbb{R}^n.
\]
For instance, one may take
$A \!\!:=\!\! \max_i a_i \sqrt{n}$ and
$B \!\!:=\!\! \sqrt{\sum_i b_i^2}$.
For any $\xi\in f(x,t)$ we can write
$\xi = -B x + T\eta + I$ with some $\eta\in\varphi(x)$, hence
\[
\begin{aligned}
   \|\xi\|
   &\le \|B\|\,\|x\| + \|T\|\,\|\eta\| + \|I\| \\
   &\le \|B\|\,\|x\| + \|T\|(A\|x\|+B) + \|I\| \\
   &= \bigl(\|B\| + A\|T\|\bigr)\,\|x\|
      + \bigl(\|T\|B + \|I\|\bigr).
\end{aligned}
\]
Thus there exist constants $K_1>0$, $K_0>0$ such that
\begin{equation}
   \|\xi\| \le K_1 \|x\| + K_0,
   \qquad \forall\,\xi\in f(x,t),\ \forall x\in\mathbb{R}^n,\ t\ge0.
   \label{eq:growth-estimate}
\end{equation}
Define
\[
   L(r) := K_1 r + K_0, \qquad r\ge0.
\]
Then $L$ is positive and continuous, and
\eqref{eq:growth-estimate} yields
\[
   \|\xi\| \le L(\|x\|),
   \qquad \forall\,\xi\in f(x,t),\ \forall x,t.
\]
Moreover,
\[
   \int\limits_0^{+\infty} \!\!\!\! \frac{dr}{L(r)}
   = \int\limits_0^{+\infty}\!\!\!\!\!\! \frac{dr}{K_1 r + K_0}
   = \frac{1}{K_1} \bigl[\ln(K_1 r + K_0)\bigr]_{0}^{+\infty}
   \! = \! +\infty.
\]

\emph{Step 3 (Wintner–type continuation).}
The conditions of the Wintner–type continuation theorem for
differential inclusions (Theorem~1.3 in~\cite[Ch.~1]{YLG}) are now
satisfied: the right–hand side $f(x,t)$ is defined for all
$(x,t)\in\mathbb{R}^n\times(0,+\infty)$, and it obeys the growth
estimate $|\xi|\le L(|x|)$ with a positive continuous function $L$
such that $\int\limits_0^{+\infty}dr/L(r)=+\infty$.
Therefore, every solution of~\eqref{eq:HNN-incl} can be continued to
all $t>0$; no solution can blow up in finite time.

Combining local existence (Step~1) with this continuation property
(Step~3), we conclude that for every initial state $x_0$ there exists
at least one Filippov solution $x(\cdot;x_0)$ of~\eqref{eq:HNN-incl}
defined on the entire half–line $[0,+\infty)$.
\end{proof}

% \subsection{Existence and Uniqueness of Equilibrium}

% In this subsection we continue to work under Assumptions~\textbf{C1},
% \textbf{C2}, \textbf{C3$'$} and $g\in\mathcal G'$.  By
% Lemma~\ref{lem:C3prime-implies-C3}, the diagonal Lyapunov condition
% \textbf{C3} also holds, and the argument in the current version of the
% draft showing existence and uniqueness of an equilibrium can be used
% verbatim.  In particular, denoting by $e\in\mathbb{R}^n$ the unique
% equilibrium one has
% \[
%   0 \in -Be + T\,\mathrm{conv}\,g(e) + I .
% \]
% We keep the detailed calculations and estimates in this subsection
% unchanged, only replacing references to ``Assumption~C3'' by
% ``Assumption~C3, which follows from~C3$'$ by
% Lemma~\ref{lem:C3prime-implies-C3}.''

\subsection{Existence and Uniqueness of Equilibrium}
% \begin{theorem}[Uniqueness]\label{thm:uniq}
% Under \textnormal{C1--C3}, the HNN admits a unique equilibrium $x^\star$.
% \end{theorem}
% \begin{IEEEproof}[Proof sketch]
% Provide the strong monotonicity or contraction argument in Filippov sense implied by C2 together with diagonal Lyapunov properties from C3. Full details are deferred to Appendix~\ref{app:proofs}.
% \end{IEEEproof}
Suppose $x^{(1)}\neq x^{(2)}$ are two equilibria; pick selections
$\xi^{(k)}\in\varphi\bigl(x^{(k)}\bigr)$ and denote
$\Delta x=x^{(1)}-x^{(2)}$, $\Delta\xi=\xi^{(1)}-\xi^{(2)}$.  
% Let
% $x^{(k)},\xi^{(k)}$ ($k=1,2$) be two equilibria.  
Subtracting yields
$(B\Delta x+T\Delta\xi)=0$ with $\Delta x\neq0$.  
Premultiply by
$\Delta x^{\!*}\alpha$ and take absolute values:
\begin{equation}\label{eq:uniq}
\begin{aligned}
    b_{\min}\|\Delta x\|^{2}
  &\leqslant \sum_{i}\alpha_i b_i|\Delta x_i|^{2}
  = |\Delta x^{\!*}\alpha T\Delta\xi| \leqslant \\
  &
  \leqslant \|T^{\!*}\alpha\|_{2}\,\|\Delta x\|\,\|\Delta\xi\|.
\end{aligned}
\end{equation}
\textbf{C2} implies $|\Delta\xi_i|\leqslant C_i|\Delta x_i|$, so
\(\|\Delta\xi\|\leqslant C_{\max}\|\Delta x\|\) with $C_{\max}:=\max_i C_i$.
Using $\|T^{\!*}\alpha\|_{2}\leqslant \|M\|_{2}/2$ (spectral norm sub‑additivity)
we get
\[
  b_{\min}\|\Delta x\|\leqslant \tfrac{\|M\|_{2}}{2}C_{\max}\|\Delta x\|,
\]
contradicting the inequality embedded in \textbf{C3$'$} (and, thus, \textbf{C3}). 
Therefore $\Delta x~=~0$ and
uniqueness holds.
% From $B\Delta x+T\Delta\xi=0$
% left–multiplying by $\Delta x^{\!*}\alpha$ yields
% \[
%  -\sum_{i}\alpha_i b_i|\Delta x_i|^{2}+\Delta x^{\!*}\alpha T\Delta\xi=0.
% \]
% Using C1 and the small–gain clause of C3 the second term is bounded in
% absolute value by $\tfrac{b_{\min}}{2}\Vert\Delta x\Vert^{2}$, whereas the
% first term is $\leqslant -b_{\min}\Vert\Delta x\Vert^{2}$.  The sum cannot vanish
% unless $\Delta x=0$, contradiction.  Hence the equilibrium is unique.

Now, similarly to \cite{p1}, we will perform a change of variables to translate $e$ into the origin. Let $e$ be an equilibrium point of the set-valued mapping $\varphi(x(t))$, i.e., by Definition \ref{equilib}:
\[
0 \in \varphi(e)=Be+T~\overline{\operatorname{conv}}(g(e))+I ,
\]
and let $\eta^* \in \overline{\operatorname{conv}}(g(e))$ be the corresponding equilibrium point of the neural network's output, such that the following equality holds
\begin{equation}\label{eq:outputequilib}
    \begin{aligned}
0 &= Be + T \eta^*+ I,\\
T \eta^* &= -(Be+I)
    \end{aligned}
\end{equation}

\begin{remark}
{The existence of $\eta^*$ can be established using reasoning analogous to that in the proof of Theorem \ref{exists}, 
i.e., by Makarov’s theorem \cite[p. 38, Th. 1.7]{Leonov}, for any solution $x(t)$, 
there exists an a.e. measurable bounded mapping $\eta(t) \in \overline{\operatorname{conv}}g(x(t))$ such that the equality in the inclusion holds.}
\end{remark}
Thus, $e$ is an equilibrium point of the neural network \eqref{eq:special} (since there is $\eta^*$ corresponding to it). 
Next we perform a change of variables $z=x-e$ to shift the equilibrium point to the origin and transform system \eqref{eq:special_conv} from definition \ref{fillipovsol} into system:
\[
\begin{aligned}
    \dot z(t) &\in Bz+ T~ \overline{\operatorname{conv}}(g(z+e))+Be+I\overset{\eqref{eq:outputequilib}}{=}\\
    &= Bz+T~ \overline{\operatorname{conv}}(g(z+e)) - T\eta^*.\\
\end{aligned}
\]
Henceforth, we consider system \eqref{sys:origin}:
\begin{equation}\label{sys:origin}
\dot z(t) \in Bz+T~ \overline{\operatorname{conv}}(G(z)),
\end{equation}
where $G(z)=g(z+e)-\eta^*$. Note that all the properties of $g$ also hold for $G$.

\begin{enumerate}
    \item $G\in \mathcal{G'}$ (see Definition \ref{G_n});
    
    \item Since $\overline{\operatorname{conv}}(G(0))=\overline{\operatorname{conv}}(g(e)) -\eta^*$, and $\eta^* \in \overline{\operatorname{conv}}(g(e))$, 
    
    it follows that $0 \in \overline{\operatorname{conv}}(G(0))$. Thus, by Definition \ref{equilib}, the origin is an equilibrium point of system \eqref{sys:origin};
    \item $ \frac{ \widetilde{\eta_i}}{z} \leqslant C_i $, $\forall x \ne 0, \quad \forall \widetilde{\eta_i} \in  \overline{\operatorname{conv}}[G_i(z)]$ (a consequence of condition \textbf{C2}).
\end{enumerate}
From now on we will write the system \eqref{eq:special_conv} of differential inclusions in the form of \eqref{sys:origin} but using symbols from \eqref{eq:special} and keeping in mind that it is indeed a system with origin as an equilibrium point:
\begin{equation}\label{eq:special_origin}
\dot x(t) = Bx+Tg(x),
\end{equation}

\subsection{HNN as a Differential Inclusion in the Lur'e Form}

After shifting the equilibrium $e$ to the origin (via $z=x-e$) and
renaming $z$ back to $x$, the network is written as the Filippov
differential inclusion
\begin{equation}\label{eq:HNN-shifted}
      \dot x(t)\in P x(t)+q\,\varphi(x(t)), \qquad x(t)\in\mathbb R^n,
\end{equation}
with $P=-B$, $q=T$, $r^{\top}=I_n$ and $\varphi=\mathrm{conv}\,g$ as
in~\eqref{eq:lure-Pqr}.  In the notation of YLG this is exactly a
Lur\'e system of the form (3.1):
\[
  \dot x = P x + q\,\xi,\qquad
  \sigma = r^{\top} x,\qquad
  \xi \in \varphi(\sigma),
\]
with several piecewise single-valued nonlinearities $\varphi_i$.

The detailed derivation of \eqref{eq:HNN-shifted} and the discussion
of the class $\mathcal G'$ in the current draft remain valid and are
not repeated here.

\subsection{YLG Lyapunov Functional and Global Asymptotic Stability}
\label{subsec:YLG-main}

We now show how Assumptions~\textbf{C1}, \textbf{C2},
\textbf{C3} and $g\in\mathcal G'$ guarantee that the Lur\'e system
\eqref{eq:HNN-shifted} satisfies all hypotheses of Theorem~3.1 in
Chapter~3 of YLG and therefore has a pointwise globally stable
stationary set.  Since we already know that the equilibrium is unique,
this yields global asymptotic stability of the Hopfield network.

\subsubsection*{Theorem 3.1 of YLG in the present notation}

For convenience, we briefly recall Theorem~3.1 of YLG in the notation
of \eqref{eq:lure-Pqr}.  Consider
\[
  \dot x = P x + q\,\xi,\qquad
  \sigma = r^{\top} x,\qquad
  \xi \in \varphi(\sigma),
\]
where $\sigma\in\mathbb{R}^m$, $\xi\in\mathbb{R}^m$ and each scalar
nonlinearity $\varphi_i$ is piecewise single–valued in the sense of
YLG.  The theorem assumes:

\begin{itemize}
  \item[(i)] irreducibility of the transfer matrix
    $X(\lambda)=r^{\top}(P-\lambda I)^{-1}q$;

  \item[(ii)] existence of diagonal matrices
    $T=\mathrm{diag}(\tau_i)\succ0$,
    $M=\mathrm{diag}(\mu_i)\succeq0$,
    $\theta=\mathrm{diag}(\vartheta_i)$ such that the frequency
    inequality (3.2) and the nondegeneracy condition (3.3) hold;

  \item[(iii)] sector and growth conditions (3.4)–(3.6) for each
    $\varphi_i$;

  \item[(iv)] existence of $\mu_i^0>\mu_i$ such that
    $P+qM_0^{-1}r^{\top}$ is Hurwitz ($M_0=\mathrm{diag}(\mu_i^0)$);

  \item[(v)] boundedness of the stationary set $\Lambda$ and exclusion
    of nonstationary complete trajectories with constant Lyapunov
    function;

  \item[(vi)] a local sector inequality (3.7) in a neighbourhood of
    $\Lambda$.
\end{itemize}

Under (i)–(v) the system is quasi-gradient-like; adding (vi) one
obtains pointwise global stability of~$\Lambda$.

\medskip
\noindent
\textbf{Verification of the YLG hypotheses.}

\begin{lemma}\label{lem:YLG-conditions}
Let $g\in\mathcal G'$ and suppose that Assumptions~\textbf{C1},
\textbf{C2} and~\textbf{C3} hold.  Then, when the Hopfield network
is written in the Lur\'e form \eqref{eq:HNN-shifted} with
$(P,q,r)=(-B,T,I_n)$ and $\varphi=\mathrm{conv}\,g$, all hypotheses
\textup{(i)}–\textup{(vi)} of Theorem~3.1 are satisfied.
\end{lemma}

\begin{proof}
We check the six items one by one.

\smallskip\noindent
\textit{(i) and (ii).}
By construction of the transfer matrix $X(\lambda)$
in~\eqref{eq:transfer-X}, Assumption~\textbf{C3} states precisely
that $X(\lambda)$ is irreducible and that there exist diagonal
matrices $\mathcal T$, $\mathcal M$, $\Theta$ such that the frequency
inequality \eqref{eq:C3prime-freq} and the nondegeneracy condition
\eqref{eq:C3prime-det} hold.  Thus the hypotheses (i) and (ii) of
Theorem~3.1 are fulfilled, with $T=\mathcal T$,
$M=\mathcal M$, $\theta=\Theta$ and $m=n$.

\smallskip\noindent
\textit{(iii).}
By Definition~1, each $g_i$ in the class $\mathcal G'$ is piecewise
continuous, sign–preserving and nondecreasing, with a finite number
of first–kind discontinuities on any compact interval and strict
upward jumps at the discontinuity points.  Its Filippov
convexification $\varphi_i=\mathrm{conv}\,g_i$ is therefore a
monotone, piecewise single–valued multifunction whose graph has no
horizontal segments at nonzero $\sigma_i$.  In particular,
$\sigma_i\,\xi_i\ge 0$ for all $\xi_i\in\varphi_i(\sigma_i)$.

Choose $\mu_i:=0$ for all $i$.  Then the sector inequality (3.4) of
Theorem~3.1 reduces to
\[
  (\sigma_i - \mu_i\xi_i)\,\xi_i
  = \sigma_i \xi_i \ge 0,
  \qquad \xi_i\in\varphi_i(\sigma_i),
\]
which holds by the sign–preserving property of $\mathcal G'$.  The
nondegeneracy condition (3.5) becomes
$\sigma_i\notin \mu_i\varphi_i(\sigma_i)$ for $\sigma_i\neq 0$, which
is trivially true when $\mu_i=0$.  Finally, the additional integral
condition (3.6) of Theorem~3.1 is required only for indices with
$\mu_i\neq0$ and $\vartheta_i<0$, so it is vacuous when
$\mu_i=0$.  Hence item (iii) is satisfied.

\smallskip\noindent
\textit{(iv).}
Condition (iv) requires the existence of constants
$\mu_i^0>\mu_i$ such that the matrix
\[
  P + qM_0^{-1}r^{\top}
  = -B + T M_0^{-1},\qquad
  M_0:=\mathrm{diag}(\mu_1^0,\dots,\mu_n^0),
\]
is Hurwitz.  Since $B=\mathrm{diag}(b_i)\succ0$ and $T$ is fixed, we
can choose $\mu_i^0$ sufficiently large so that, for every row $k$,
\[
  b_k
  > \sum_{j=1}^n \frac{|T_{kj}|}{\mu_j^0}.
\]
Equivalently, the Gershgorin discs of $-B+T M_0^{-1}$ lie strictly in
the open left half–plane, which implies that $-B+T M_0^{-1}$ is
Hurwitz.  Thus (iv) holds for a suitable choice of $M_0$ with
$\mu_i^0>\mu_i=0$.

\smallskip\noindent
\textit{(v).}
Let $\Lambda$ denote the stationary set of the Lur\'e system
\eqref{eq:HNN-shifted}.  In Section~\ref{sec:problem} and in the
subsection on existence and uniqueness of equilibrium we have already
shown that, under \textbf{C1}, \textbf{C2} and \textbf{C3}
(and hence \textbf{C3$'$} by Lemma~\ref{lem:C3prime-implies-C3}), the
Hopfield inclusion \eqref{eq:HNN-incl} admits a \emph{unique}
equilibrium $e$.  After the equilibrium shift, this means that the
stationary set of \eqref{eq:HNN-shifted} is the singleton
$\Lambda=\{0\}$, which is certainly bounded.

The second part of (v) excludes nonstationary complete trajectories
along which a Lyapunov function is constant.  Define the YLG
Lyapunov function
\begin{equation}\label{eq:YLG-V-def}
    V(x)
    := \frac12\,x^\top H x + 
       \sum_{i=1}^n \vartheta_i
       \int\limits_0^{\sigma_i} \varphi_i(\lambda)\,d\lambda,
\end{equation}
where $H\succ0$ is the symmetric matrix obtained from the YLG
frequency inequality (cf.\ Eqns.~(3.10)–(3.11) in Chapter~3 of YLG)
and $\vartheta_i$ are the diagonal entries of~$\Theta$.  Because each
$\varphi_i$ is monotone and has at most linear growth (by
\textbf{C1}), the integrals in~\eqref{eq:YLG-V-def} are well defined
and $V$ is continuous and radially unbounded.

Following exactly the differentiation argument of YLG (see the proof
of Theorem~3.1 on pp.~113–114), one obtains for any Filippov solution
$x(\cdot)$ of \eqref{eq:HNN-shifted} and any measurable selection
$\xi(\cdot)\in\varphi(\sigma(\cdot))$ the estimate
\[
  \dot V(x(t)) \le 0 \quad\text{for a.e. }t\in\mathbb{R},
\]
with equality if and only if simultaneously
$\xi_i(t)\,\sigma_i(t)=0$ for all $i$ and
$x(t)^\top H x(t)\equiv\mathrm{const}$.  Hence, if a complete
trajectory satisfies the identities stated in item (v), it must have
$\dot V\equiv 0$ and therefore $Px(t)+q\xi(t)\equiv 0$.  In other
words, the trajectory is stationary.  This verifies (v).

\smallskip\noindent
\textit{(vi).}
With $\mu_i=0$, condition (vi) of Theorem~3.1 reduces to
\[
  (\xi_i-\xi_i^0)\bigl(\sigma_i-\sigma_i^0\bigr) \ge 0,
  \qquad \forall\xi_i\in\varphi_i(\sigma_i),
\]
for all stationary $(x^0,\xi^0)$ and all $(x,\xi)$ sufficiently close
to $\Lambda$.  Since each $\varphi_i$ is monotone nondecreasing, this
inequality holds for every pair $(\sigma_i,\xi_i)$,
$(\sigma_i^0,\xi_i^0)$ with $\xi_i\in\varphi_i(\sigma_i)$,
$\xi_i^0\in\varphi_i(\sigma_i^0)$, not only in a neighbourhood of the
stationary set.  Thus (vi) is satisfied as well.
\end{proof}

\subsubsection*{Conclusion: global asymptotic stability}

We can now state and prove the main result.

% \begin{theorem}[Global asymptotic stability of a Hopfield network with discontinuous activations]\label{thm:main}
% Consider the continuous–time Hopfield neural network
% \begin{equation}
%   \dot x(t) = -B x(t) + T\,g(x(t)) + I, \qquad x(t)\in\mathbb{R}^n,
%   \label{eq:HNN-main}
% \end{equation}
% where
% \begin{itemize}
%   \item $B = \mathrm{diag}(b_1,\dots,b_n)\succ 0$ is a diagonal matrix
%         of leak rates;
%   \item $T\in\mathbb{R}^{n\times n}$ is the interconnection (synaptic) matrix;
%   \item $I\in\mathbb{R}^n$ is a constant input vector;
%   \item $g=(g_1,\dots,g_n)^{\top}:\mathbb{R}^n\to\mathbb{R}^n$ is the
%         componentwise activation mapping, each
%         $g_i:\mathbb{R}\to\mathbb{R}$ possibly discontinuous.
% \end{itemize}
% Assume that $g$ belongs to the extended class $\mathcal{G}'$ of
% activation functions (Definition~1) and that Assumptions
% \textbf{C1}–\textbf{C3$'$} hold.

% Interpret system~\eqref{eq:HNN-main} in the sense of Filippov, i.e.\
% as the differential inclusion
% \[
%    \dot x(t) \in -B x(t) + T\,\varphi(x(t)) + I,
%    \qquad \varphi(x)=\mathrm{conv}\,g(x).
% \]
% Then the following properties hold.
% \begin{enumerate}
%   \item For every initial condition $x(0)=x_0\in\mathbb{R}^n$ there
%         exists at least one Filippov solution $x(\cdot;x_0)$
%         of~\eqref{eq:HNN-main} defined on the whole half–line
%         $[0,+\infty)$.

%   \item There exists a unique equilibrium point $e\in\mathbb{R}^n$,
%         i.e.\ a constant vector satisfying
%         \[
%            0 \in -B e + T\,\mathrm{conv}\,g(e) + I .
%         \]

%   \item This equilibrium is \emph{globally asymptotically stable in
%         the sense of Filippov}: for every Filippov solution
%         $x(\cdot;x_0)$ of~\eqref{eq:HNN-main} one has
%         \[
%            \lim_{t\to+\infty} x(t;x_0) = e ,
%         \]
%         and for every $\varepsilon>0$ there exists $\delta>0$ such
%         that $|x_0-e|<\delta$ implies $|x(t;x_0)-e|<\varepsilon$ for
%         all $t\ge0$.
% \end{enumerate}
% \end{theorem}

\begin{proof}
Item~1 is exactly Proposition~\ref{prop:global-existence}.  Item~2 is
established in the subsection on existence and uniqueness of
equilibrium, which uses Assumptions~\textbf{C1}, \textbf{C2} and the
derived condition~\textbf{C3$'$} (hence ultimately \textbf{C3}) to
exclude multiple equilibria.

For item~3, perform the equilibrium shift $z=x-e$ and rewrite the
network in the Lur\'e form \eqref{eq:HNN-shifted}.  By
Lemma~\ref{lem:YLG-conditions}, all hypotheses (i)–(vi) of
Theorem~3.1 in Chapter~3 of YLG are satisfied.  Therefore the
stationary set $\Lambda$ of \eqref{eq:HNN-shifted} is pointwise
globally stable.  Since $\Lambda=\{0\}$ (uniqueness of equilibrium),
this means that the origin is globally asymptotically stable for the
shifted inclusion.  Undoing the shift gives global asymptotic
stability of the equilibrium $e$ for the original Hopfield network
\eqref{eq:HNN-main} in the sense of Filippov.
\end{proof}
\todo[inline]{New above}

% \section{Problem Statement and Assumptions}\label{sec:problem}

% Consider the Hopfield neural network (HNN) in Lur\'e form
% \begin{equation}\label{eq:HNN-lure}
%   \dot{x} = -Bx + T\,g(x) + I, 
%   \qquad x(t)\in\mathbb{R}^n,
% \end{equation}
% where $B=\mathrm{diag}(b_1,\dots,b_n)\succ 0$ collects effective leak
% rates, $T\in\mathbb{R}^{n\times n}$ is the interconnection (synaptic)
% matrix, $I\in\mathbb{R}^n$ is the (rescaled) input vector, and
% $g(x)=(g_1(x_1),\ldots,g_n(x_n))^\top$ is the componentwise activation
% which may be discontinuous.

% \begin{assumption}[\bf C1: Linear growth]\label{ass:C1}
% For each $i\in\{1,\dots,n\}$ there exist constants $a_i,b_i\ge 0$ such
% that
% \[
%   |g_i(\xi)| \le a_i|\xi| + b_i,
%   \qquad \forall\,\xi\in\mathbb{R}.
% \]
% \end{assumption}

% \begin{assumption}[\bf C2: One-sided Lipschitz (componentwise)]\label{ass:C2}
% For each $i\in\{1,\dots,n\}$ there exists $C_i\in\mathbb{R}$ such that
% for all $x,y\in\mathbb{R}^n$ and for all a.e.\ selections
% $\eta\in\mathrm{co}\,g(x)$, $\zeta\in\mathrm{co}\,g(y)$ one has
% \[
%   (\eta_i - \zeta_i)\,(x_i - y_i)
%   \;\le\; C_i\,(x_i - y_i)^2.
% \]
% \end{assumption}

% \begin{assumption}[\bf C3: Diagonal Lyapunov stability]\label{ass:C3}
% There exists a diagonal matrix $H=\mathrm{diag}(h_1,\dots,h_n)\succ 0$
% and a constant $\alpha\ge 0$ such that
% \begin{equation}\label{eq:C3-matrix}
%   -HB - B^{\top}H + \alpha\bigl(T^{\top}H + HT\bigr) \prec 0 .
% \end{equation}
% \end{assumption}

% \begin{remark}
% Condition~\textbf{C2} is also known as a \emph{unilateral
% Lipschitz-like} condition and, together with the linear growth
% Assumption~\textbf{C1}, it is used in a similar way in several works
% on Hopfield-type networks with discontinuous activations,
% see e.g.\ \cite{unilateral_Lipschitz_like_condition,eq_ref_11,
% GAS_of_delayed,GAS_discrete_time_varying,periodic_time_varying}.
% \end{remark}

% \section{Main Results}\label{sec:main}

% Let us treat the Hopfield network as a concrete Lur'e--type system and
% literally ``plug it'' into the YLG framework. 
% We proceed in three main steps:
% \begin{enumerate}
%   \item rewrite the HNN in the exact Lur'e form used by YLG;
%   \item check that Assumptions~C1--C3 imply the hypotheses of the
%         relevant YLG theorems;
%   \item invoke those theorems together with Lyapunov lemmas from
%         Chapter~1 of~\cite{YLG} to conclude global asymptotic stability
%         of the unique equilibrium in the Filippov sense.
% \end{enumerate}

% \subsection{Global Existence of Filippov Solutions for the HNN}

% In this subsection we show that under Assumption~C1 every Filippov
% solution of the Hopfield network exists on the entire half–line
% $[0,+\infty)$.

% Recall that, in the sense of Filippov, system~\eqref{eq:special} is
% understood as the differential inclusion
% \begin{equation}
%     \dot x(t) \in -B x(t) + T\,\varphi(x(t)) + I,
%     \qquad x(t)\in\mathbb{R}^n,
%     \label{eq:HNN-incl}
% \end{equation}
% where $\varphi(x) \!=\! \operatorname{conv} \! g(x)$ denotes the Filippov
% convexification of the (possibly discontinuous) activation mapping~$g\!:\!\mathbb{R}^n\!\!\to\!\!\mathbb{R}^n$.

% \begin{proposition}[Global existence]
% \label{prop:global-existence}
% Suppose that Assumption~C1 (linear growth of the activation functions)
% holds. Then, for every initial condition $x(0)=x_0\in\mathbb{R}^n$,
% there exists at least one Filippov solution $x(\cdot;x_0)$
% of~\eqref{eq:HNN-incl} defined on the whole half–line $[0,+\infty)$.
% \end{proposition}

% \begin{proof}
% \emph{Step 1 (local existence).}
% Fix $t_0\in[0,+\infty)$ and $a\in\mathbb{R}^n$, and consider the
% cylinder
% \[
%     D := \{(x,t)\in\mathbb{R}^n\times\mathbb{R} \mid
%            |t-t_0|\le\alpha,\ |x-a|\le\rho\},
% \]
% with some $\alpha>0$, $\rho>0$.
% Since $g$ is piecewise continuous with a finite number of first–order
% discontinuities on any compact set (Definition~1) and $\varphi$ is its
% Filippov convexification, the multifunction
% \[
%    f(x,t) := -B x + T\,\varphi(x) + I
% \]
% is upper semicontinuous in $(x,t)$ and has nonempty, closed, bounded
% and convex values (see, e.g.,~\cite[Ch.~1]{YLG}).
% On the compact set $D$ the values of $f$ are uniformly bounded by some
% constant $c>0$.

% Therefore, the hypotheses of the local existence theorem for
% differential inclusions (Theorem~1.1 in~\cite[Ch.~1]{YLG}) are
% satisfied, and for any initial condition $x(t_0)=a$ there exists at
% least one absolutely continuous solution of~\eqref{eq:HNN-incl} on
% the interval $|t-t_0|\le\tau$, where
% $\tau = \min\{\alpha,\rho/c\}$. Hence Filippov solutions of
% \eqref{eq:HNN-incl} exist locally.

% \emph{Step 2 (global growth estimate from C1).}
% Assumption~C1 states that for each scalar activation $g_i$ there exist
% constants $a_i,b_i\ge0$ such that
% \[
%    |g_i(\xi)| \le a_i|\xi| + b_i, \qquad \forall\,\xi\in\mathbb{R}.
% \]
% It follows that there are constants $A,B\ge0$ with
% \[
%    \|\eta\| \le A\|x\| + B,
%    \qquad \forall\,\eta\in\varphi(x),\ \forall x\in\mathbb{R}^n.
% \]
% For instance, one may take
% $A \!\!:=\!\! \max_i a_i \sqrt{n}$ and
% $B \!\!:=\!\! \sqrt{\sum_i b_i^2}$.
% For any $\xi\in f(x,t)$ we can write
% $\xi = -B x + T\eta + I$ with some $\eta\in\varphi(x)$, hence
% \[
% \begin{aligned}
%    \|\xi\|
%    &\le \|B\|\,\|x\| + \|T\|\,\|\eta\| + \|I\| \\
%    &\le \|B\|\,\|x\| + \|T\|(A\|x\|+B) + \|I\| \\
%    &= \bigl(\|B\| + A\|T\|\bigr)\,\|x\|
%       + \bigl(\|T\|B + \|I\|\bigr).
% \end{aligned}
% \]
% Thus there exist constants $K_1>0$, $K_0>0$ such that
% \begin{equation}
%    \|\xi\| \le K_1 \|x\| + K_0,
%    \qquad \forall\,\xi\in f(x,t),\ \forall x\in\mathbb{R}^n,\ t\ge0.
%    \label{eq:growth-estimate}
% \end{equation}
% Define
% \[
%    L(r) := K_1 r + K_0, \qquad r\ge0.
% \]
% Then $L$ is positive and continuous, and
% \eqref{eq:growth-estimate} yields
% \[
%    \|\xi\| \le L(\|x\|),
%    \qquad \forall\,\xi\in f(x,t),\ \forall x,t.
% \]
% Moreover,
% \[
%    \int_0^{+\infty} \!\!\!\! \frac{dr}{L(r)}
%    = \int_0^{+\infty}\!\!\!\!\!\! \frac{dr}{K_1 r + K_0}
%    = \frac{1}{K_1} \bigl[\ln(K_1 r + K_0)\bigr]_{0}^{+\infty}
%    \! = \! +\infty.
% \]

% \emph{Step 3 (Wintner–type continuation).}
% The conditions of the Wintner–type continuation theorem for
% differential inclusions (Theorem~1.3 in~\cite[Ch.~1]{YLG}) are now
% satisfied: the right–hand side $f(x,t)$ is defined for all
% $(x,t)\in\mathbb{R}^n\times(0,+\infty)$, and it obeys the growth
% estimate $|\xi|\le L(|x|)$ with a positive continuous function $L$
% such that $\int_0^{+\infty}dr/L(r)=+\infty$.
% Therefore, every solution of~\eqref{eq:HNN-incl} can be continued to
% all $t>0$; no solution can blow up in finite time.

% Combining local existence (Step~1) with this continuation property
% (Step~3), we conclude that for every initial state $x_0$ there exists
% at least one Filippov solution $x(\cdot;x_0)$ of~\eqref{eq:HNN-incl}
% defined on the entire half–line $[0,+\infty)$.
% \end{proof}

% \subsection{Existence and Uniqueness of Equilibrium}
% % \begin{theorem}[Uniqueness]\label{thm:uniq}
% % Under \textnormal{C1--C3}, the HNN admits a unique equilibrium $x^\star$.
% % \end{theorem}
% % \begin{IEEEproof}[Proof sketch]
% % Provide the strong monotonicity or contraction argument in Filippov sense implied by C2 together with diagonal Lyapunov properties from C3. Full details are deferred to Appendix~\ref{app:proofs}.
% % \end{IEEEproof}
% Suppose $x^{(1)}\neq x^{(2)}$ are two equilibria; pick selections
% $\xi^{(k)}\in\varphi\bigl(x^{(k)}\bigr)$ and denote
% $\Delta x=x^{(1)}-x^{(2)}$, $\Delta\xi=\xi^{(1)}-\xi^{(2)}$.  
% % Let
% % $x^{(k)},\xi^{(k)}$ ($k=1,2$) be two equilibria.  
% Subtracting yields
% $(B\Delta x+T\Delta\xi)=0$ with $\Delta x\neq0$.  
% Premultiply by
% $\Delta x^{\!*}\alpha$ and take absolute values:
% \begin{equation}\label{eq:uniq}
% \begin{aligned}
%     b_{\min}\|\Delta x\|^{2}
%   &\leqslant \sum_{i}\alpha_i b_i|\Delta x_i|^{2}
%   = |\Delta x^{\!*}\alpha T\Delta\xi| \leqslant \\
%   &
%   \leqslant \|T^{\!*}\alpha\|_{2}\,\|\Delta x\|\,\|\Delta\xi\|.
% \end{aligned}
% \end{equation}
% \textbf{C2} implies $|\Delta\xi_i|\leqslant C_i|\Delta x_i|$, so
% \(\|\Delta\xi\|\leqslant C_{\max}\|\Delta x\|\) with $C_{\max}:=\max_i C_i$.
% Using $\|T^{\!*}\alpha\|_{2}\leqslant \|M\|_{2}/2$ (spectral norm sub‑additivity)
% we get
% \[
%   b_{\min}\|\Delta x\|\leqslant \tfrac{\|M\|_{2}}{2}C_{\max}\|\Delta x\|,
% \]
% contradicting the inequality embedded in \textbf{C3}. 
% Therefore $\Delta x~=~0$ and
% uniqueness holds.
% % From $B\Delta x+T\Delta\xi=0$
% % left–multiplying by $\Delta x^{\!*}\alpha$ yields
% % \[
% %  -\sum_{i}\alpha_i b_i|\Delta x_i|^{2}+\Delta x^{\!*}\alpha T\Delta\xi=0.
% % \]
% % Using C1 and the small–gain clause of C3 the second term is bounded in
% % absolute value by $\tfrac{b_{\min}}{2}\Vert\Delta x\Vert^{2}$, whereas the
% % first term is $\leqslant -b_{\min}\Vert\Delta x\Vert^{2}$.  The sum cannot vanish
% % unless $\Delta x=0$, contradiction.  Hence the equilibrium is unique.

% Now, similarly to \cite{p1}, we will perform a change of variables to translate $e$ into the origin. Let $e$ be an equilibrium point of the set-valued mapping $\varphi(x(t))$, i.e., by Definition \ref{equilib}:
% \[
% 0 \in \varphi(e)=Be+T~\overline{\operatorname{conv}}(g(e))+I ,
% \]
% and let $\eta^* \in \overline{\operatorname{conv}}(g(e))$ be the corresponding equilibrium point of the neural network's output, such that the following equality holds
% \begin{equation}\label{eq:outputequilib}
%     \begin{aligned}
% 0 &= Be + T \eta^*+ I,\\
% T \eta^* &= -(Be+I)
%     \end{aligned}
% \end{equation}

% \begin{remark}
% {The existence of $\eta^*$ can be established using reasoning analogous to that in the proof of Theorem \ref{exists}, i.e., by Makarov’s theorem \cite[p. 38, Th. 1.7]{Leonov}, for any solution $x(t)$, there exists a a.e. measurable bounded mapping $\eta(t) \in \overline{\operatorname{conv}}g(x(t))$ such that the equality in the inclusion holds.}
% \end{remark}
% Thus, $e$ is an equilibrium point of the neural network \eqref{eq:special} (since there is $\eta^*$ corresponding to it). 
% Next we perform a change of variables $z=x-e$ to shift the equilibrium point to the origin and transform system \eqref{eq:special_conv} from definition \ref{fillipovsol} into system:
% \[
% \begin{aligned}
%     \dot z(t) &\in Bz+ T~ \overline{\operatorname{conv}}(g(z+e))+Be+I\overset{\eqref{eq:outputequilib}}{=}\\
%     &= Bz+T~ \overline{\operatorname{conv}}(g(z+e)) - T\eta^*.\\
% \end{aligned}
% \]
% Henceforth, we consider system \eqref{sys:origin}:
% \begin{equation}\label{sys:origin}
% \dot z(t) \in Bz+T~ \overline{\operatorname{conv}}(G(z)),
% \end{equation}
% where $G(z)=g(z+e)-\eta^*$. Note that all the properties of $g$ also hold for $G$.

% \begin{enumerate}
%     \item $G\in \mathcal{G'}$ (see Definition \ref{G_n});
    
%     \item Since $\overline{\operatorname{conv}}(G(0))=\overline{\operatorname{conv}}(g(e)) -\eta^*$, and $\eta^* \in \overline{\operatorname{conv}}(g(e))$, 
    
%     it follows that $0 \in \overline{\operatorname{conv}}(G(0))$. Thus, by Definition \ref{equilib}, the origin is an equilibrium point of system \eqref{sys:origin};
%     \item $ \frac{ \widetilde{\eta_i}}{z} \leqslant C_i $, $\forall x \ne 0, \quad \forall \widetilde{\eta_i} \in  \overline{\operatorname{conv}}[G_i(z)]$ (a consequence of condition \textbf{C2}).
% \end{enumerate}
% From now on we will write the system \eqref{eq:special_conv} of differential inclusions in the form of \eqref{sys:origin} but using symbols from \eqref{eq:special} and keeping in mind that it is indeed a system with origin as an equilibrium point:
% \begin{equation}\label{eq:special_origin}
% \dot x(t) = Bx+Tg(x),
% \end{equation}

% \subsection{HNN as a differential inclusion in the Lur'e form}

% After shifting an equilibrium to the origin, 
% the network is written as the Filippov differential inclusion
% $$
%       \dot x(t)\in Bx(t)+T\,\overline{\operatorname{conv}}\,g(x(t)), \qquad x\in\mathbb R^n, 
% $$
% where $B=\mathrm{diag}(b_1,\dots,b_n)\succ0$ (effective leak rates), 
% $T\in\mathbb R^{n\times n}$ (synaptic matrix),
% $g(x)=(g_1(x_1),\dots,g_n(x_n))^\top$ is the componentwise activation, possibly discontinuous,
% $\mathrm{conv}\,g(x)$ denotes the Filippov convexification 
% (smallest closed convex hull of limiting values).

% In the notation of~\cite{YLG} this is a Lur\'e system
% % \dot x = Px+q\,\xi,\qquad a=r^*x,\qquad \xi\in\varphi(a),
% % \tag{L}
% % with
% % P=B,\quad q=T,\quad r^*=I_n,\quad a=x,\quad
% % \varphi(a)=\mathrm{conv}\,g(a).
% % Write \eqref{eq:special_origin} as a Lur\'e system
% \begin{equation}\label{eq:lure}
%   \dot x=Px+q\xi,\qquad a=r^{*}x,\qquad \xi \in \varphi(a),
% \end{equation}
% with
% $P=B$, \;$q=T$, \;$r^{*}=I_n$ and the \emph{extended nonlinearity}
% \(\varphi(a)=\overline{\operatorname{conv}} g(a)\) in the sense of~\cite[\S1.1.2]{YLG}.
% Because $g$ satisfies \textbf{C1--C2}, $\varphi$ is upper semicontinuous with bounded
% closed convex values, so all Filippov solutions of~\eqref{eq:lure} exist for
% $[0,\infty)$ (see. \cite[Theorems 1 and 2, pp. 77-78]{p2}).

% By Definition~\ref{G_n}, 
% each $g_i$ belongs to the class $\mathcal G'$: 
% piecewise continuous on $\mathbb R$ with finitely many first-kind discontinuities 
% on any compact set and strict jumps $g_i(\rho_k^+)>g_i(\rho_k^-)$. 
% Therefore each scalar multivalued $\varphi_i$ is piecewise single–valued in the sense of YLG 
% (single–valued and continuous except for isolated jump points where 
% the value is a closed segment between one-sided limits).
% Hence~\eqref{eq:HNN} is exactly a system of as in~\cite[see Eq.~(3.1)]{YLG}:
% $\dot x = Px+q\varphi(a)$, $a=r^*x$,
% with several piecewise single-valued nonlinearities.

% \subsection{Alternative Frequency-domain Assumption for the linear part}

% For some applications it is convenient to replace the algebraic
% Assumption~\textbf{C3} by a purely frequency-domain condition in the
% spirit of Yakubovich–Leonov–Gelig.  
% To this end, rewrite~\eqref{eq:HNN-lure} in the standard Lur\'e form
% \begin{equation}\label{eq:lure-Pqr}
%   \dot x = P x + q\,\xi, 
%   \qquad \sigma = r^{\top}x, 
%   \qquad \xi = g(\sigma),
% \end{equation}
% with
% \begin{equation}\label{eq:Pqr-def}
%   P := -B, \qquad q := T, \qquad r := I_n .
% \end{equation}
% The transfer matrix of the linear part is
% \begin{equation}\label{eq:transfer-X}
%   X(\lambda) 
%   := r^{\top}(P-\lambda I)^{-1}q
%   = (-B-\lambda I)^{-1}T,
%   \qquad \lambda\in\mathbb{C}.
% \end{equation}

% \begin{assumption}[C3$'$: YLG-type frequency condition]\label{ass:C3prime}
% We say that \textbf{C3$'$} holds for the HNN \eqref{eq:HNN-lure} if
% the following two properties are satisfied.
% \begin{enumerate}%[(i)]
%   \item \textbf{Irreducibility of the transfer matrix.}
%   The transfer function $X(\lambda)$ defined in
%   \eqref{eq:transfer-X} is irreducible in the sense of
%   Definition~2.6 in~\cite[Ch.~2]{YLG}, i.e.\ the pair
%   $(P,r)=(-B,I_n)$ is observable.

%   \item \textbf{Yakubovich–Barabanov frequency inequality.}
%   There exist diagonal matrices
%   \begin{multline*}   
%      \mathcal{T}=\mathrm{diag}(\tau_1,\dots,\tau_n)\succ 0,\\
%      \mathcal{M}=\mathrm{diag}(\mu_1,\dots,\mu_n)\succeq 0,\\
%      \Theta=\mathrm{diag}(\vartheta_1,\dots,\vartheta_n)
%   \end{multline*}
%   such that
%   \begin{equation}\label{eq:C3prime-freq}
%      \mathcal{T}\mathcal{M}
%      + \operatorname{Re}\Bigl[
%         \bigl(\mathcal{T} + i\omega \Theta\bigr)\,
%         X(i\omega)
%        \Bigr]
%      \succeq 0,
%      \qquad \forall\,\omega\in\mathbb{R},
%   \end{equation}
%   and the nondegeneracy condition
%   \begin{multline}\label{eq:C3prime-det}
%     \det\bigl(\mathcal{T} + \lambda \Theta\bigr) \neq 0
%     \quad\text{whenever}\quad \\
%     \det(P-\lambda I)=\det(-B-\lambda I)=0
%   \end{multline}
%   holds.
% \end{enumerate}
% These are exactly conditions \textup{(i)}–\textup{(ii)} of
% Theorem~3.1 in~\cite[Ch.~3]{YLG}, specialized to the triple
% $(P,q,r)=(-B,T,I_n)$ of the Hopfield model.
% \end{assumption}

% The next lemma shows that the frequency-domain Assumption~\textbf{C3$'$}
% is stronger than, and hence implies, the algebraic
% Assumption~\textbf{C3}.

% \begin{lemma}[C3$'$ implies C3]\label{lem:C3prime-implies-C3}
% Suppose that Assumption~\textbf{C3$'$} holds. Then there exist a
% diagonal matrix $H=\mathrm{diag}(h_1,\dots,h_n)\succ 0$ and a constant
% $\alpha>0$ such that the matrix inequality \eqref{eq:C3-matrix} is
% satisfied; in particular, Assumption~\textbf{C3} holds.
% \end{lemma}

% \begin{proof}
% Under Assumption~\textbf{C3$'$}, conditions \textup{(i)}–\textup{(ii)}
% of Theorem~3.1 in~\cite[Ch.~3]{YLG} are satisfied for the system
% \eqref{eq:lure-Pqr} with $(P,q,r)=(-B,T,I_n)$.  By
% \cite[Eqns.~(3.10)–(3.11)]{YLG} there exists a symmetric matrix
% $H_{\mathrm{YLG}}=H_{\mathrm{YLG}}^{\top}$ such that the quadratic form
% \begin{equation}\label{eq:W-def}
%   W(x,\xi)
%   := \bigl[2x^{\top}H_{\mathrm{YLG}} + \xi^{\top}\Theta r^{\top}\bigr]\,
%      (P x + q\xi)
%      + \xi^{\top}\mathcal{T}\,(r^{\top}x - \mathcal{M}\xi)
% \end{equation}
% satisfies
% \begin{equation}\label{eq:W-leq0}
%   W(x,\xi) \le 0,
%   \qquad \forall\,x\in\mathbb{R}^n,\ \forall\,\xi\in\mathbb{R}^n.
% \end{equation}
% Moreover, the second part of Theorem~3.1, see in particular
% \cite[Eq.~(3.14)]{YLG}, implies that $H_{\mathrm{YLG}}$ can be
% chosen positive definite:
% \begin{equation}\label{eq:H-YLG-pd}
%   H_{\mathrm{YLG}} \succ 0 .
% \end{equation}

% Set $\xi\equiv 0$ in \eqref{eq:W-def}.  Using $r^{\top}=I_n$ and
% $P=-B$, we obtain
% \[
%   W(x,0)
%    = \bigl[2x^{\top}H_{\mathrm{YLG}}\bigr]\,(P x)
%    = 2x^{\top}H_{\mathrm{YLG}}P x
%    = -2x^{\top}H_{\mathrm{YLG}}B x .
% \]
% Together with \eqref{eq:W-leq0} this gives
% \begin{equation}\label{eq:HB-negative}
%   -2x^{\top}H_{\mathrm{YLG}}B x \le 0,
%   \qquad \forall\,x\in\mathbb{R}^n.
% \end{equation}
% Since $B=\mathrm{diag}(b_1,\dots,b_n)\succ 0$ and
% $H_{\mathrm{YLG}}\succ 0$, the congruence
% $B^{1/2}H_{\mathrm{YLG}}B^{1/2}$ is also positive definite, so
% $x^{\top}H_{\mathrm{YLG}}B x>0$ for all $x\neq 0$.  Hence
% \eqref{eq:HB-negative} is in fact strict for $x\neq 0$, and the
% symmetric matrix
% \[
%   S_{\mathrm{YLG}}
%   := -H_{\mathrm{YLG}}B - B^{\top}H_{\mathrm{YLG}}
% \]
% is negative definite:
% \begin{equation}\label{eq:SYLG-negative}
%   x^{\top}S_{\mathrm{YLG}}x
%   = -2x^{\top}H_{\mathrm{YLG}}B x < 0,
%   \qquad \forall\,x\neq 0 .
% \end{equation}

% Now define the diagonal matrix
% \begin{equation}\label{eq:H-diag-def}
%   H := \mathrm{diag}\bigl((H_{\mathrm{YLG}})_{11},
%                          \dots,
%                          (H_{\mathrm{YLG}})_{nn}\bigr).
% \end{equation}
% By \eqref{eq:H-YLG-pd} all diagonal entries $(H_{\mathrm{YLG}})_{ii}$,
% $i=1,\dots,n$, are strictly positive, hence $H\succ 0$.  Since $B$ is
% diagonal and positive definite, for every $x\neq 0$ we have
% \begin{multline*}
%   x^{\top}\bigl(-HB - B^{\top}H\bigr) x = -2x^{\top}HBx = \\
%    = -2(B^{1/2}x)^{\top}\bigl(B^{1/2}HB^{1/2}\bigr)(B^{1/2}x) < 0,    
% \end{multline*}
% because $B^{1/2}HB^{1/2}$ is again positive definite.  Thus
% \begin{equation}\label{eq:S0-negative}
%   S_0 := -HB - B^{\top}H \prec 0 .
% \end{equation}

% Let
% \begin{equation}\label{eq:S1-def}
%   S_1 := T^{\top}H + HT
% \end{equation}
% and consider the one-parameter family of symmetric matrices
% \begin{equation}\label{eq:Salpha-def}
%   S(\alpha) := S_0 + \alpha S_1,
%   \qquad \alpha\ge 0.
% \end{equation}
% The eigenvalues of $S(\alpha)$ depend continuously on~$\alpha$, and by
% \eqref{eq:S0-negative} we have $S(0)=S_0\prec 0$.  Therefore there
% exists $\alpha^\ast>0$ such that $S(\alpha)\prec 0$ for all
% $\alpha\in[0,\alpha^\ast]$.  Choosing any
% $\alpha\in(0,\alpha^\ast]$ we obtain
% \[
%   -HB - B^{\top}H + \alpha\bigl(T^{\top}H + HT\bigr)
%   = S(\alpha) \prec 0 ,
% \]
% which is precisely the inequality \eqref{eq:C3-matrix} in
% Assumption~\textbf{C3}. This completes the proof.
% \end{proof}

% \subsection{YLG Lyapunov functional and global asymptotic stability}
% \label{subsec:YLG-main}

% In this subsection we show how Assumptions~C1–C3, together with the
% Filippov interpretation of system~\eqref{eq:special}, fit into the YLG
% framework and imply global asymptotic stability of the unique
% equilibrium for all Filippov solutions.

% \medskip
% \subsubsection{YLG Lyapunov functional for the Hopfield inclusion}

% After shifting the (Filippov) equilibrium $e$ to the origin,
% system~\eqref{eq:special} can be written in the form
% \begin{equation}\label{eq:HNN-shifted}
%     \dot x(t)\in Bx(t)+T\,\conv g(x(t)),\qquad x(t)\in\R^n,
% \end{equation}
% with $B=\diag(b_1,\dots,b_n)\succ0$, $T\in\R^{n\times n}$ and
% $g=(g_1,\dots,g_n)\in\mathcal G'$.
% By the measurable–selection theorem of Makarov (YLG Theorem~1.7) there
% exists, for any Filippov solution $x(\cdot)$ of
% \eqref{eq:HNN-shifted}, a measurable ``extended activation''
% \[
%     \eta(t)=(\eta_1(t),\dots,\eta_n(t))^\top,\qquad
%     \eta(t)\in\conv g(x(t)) \quad\text{a.e.},
% \]
% such that $x(\cdot)$ satisfies the ordinary differential equation
% \begin{equation}\label{eq:HNN-eta}
%     \dot x(t)=Bx(t)+T\,\eta(t)\quad\text{a.e.}
% \end{equation}

% Choose a diagonal matrix $H\succ0$ as in Assumption~C3 and define the
% YLG Lyapunov functional
% \begin{equation}\label{eq:YLG-V-def}
%     V(x)\;:=\;\frac12\,x^\top H x + \sum_{i=1}^n \Psi_i(x_i),
%     \qquad
%     \Psi_i(s):=\int_0^s \eta_i(\sigma)\,d\sigma,
% \end{equation}
% where, for each $i$, $\eta_i(\cdot)$ is any measurable selection from
% $\conv g_i(\cdot)$.  Since each $\conv g_i$ has nonempty, convex,
% bounded values and is piecewise single–valued (class $\mathcal G'$),
% the integrals $\Psi_i$ are well defined and $V$ is a convex,
% locally Lipschitz function on~$\R^n$ (hence regular in the sense of
% Clarke).

% Differentiating~\eqref{eq:YLG-V-def} along the solution
% \eqref{eq:HNN-eta} and using the generalized chain rule for regular
% functions (see YLG, Ch.~1) we obtain, for a.e.~$t\ge0$,
% \begin{align}
%     \dot V(x(t))
%       &= x(t)^\top H\dot x(t)
%          + \sum_{i=1}^n \eta_i(t)\,\dot x_i(t)\notag\\
%       &= x^\top H\bigl(Bx+T\eta\bigr)
%          + \eta^\top\bigl(Bx+T\eta\bigr)\notag\\
%       &= x^\top( HB+B^\top H )x
%          + \eta^\top(\alpha T+T^\top\alpha)\eta,
%       \label{eq:Vdot-raw}
% \end{align}
% where in the last step we have regrouped terms in the standard YLG
% way (cf.\ Ch.~3 of YLG).
% Assumption~C3 gives, for some $\lambda_1,\lambda_2>0$,
% \begin{equation}\label{eq:Vdot-neg}
%     \dot V(x(t))
%     \;\le\; -\lambda_1 \|x(t)\|^2 - \lambda_2 \|\eta(t)\|^2
%     \quad\text{for a.e. }t\ge0.
% \end{equation}
% In particular, $V(x(t))$ is nonincreasing along every Filippov solution
% of~\eqref{eq:HNN-shifted}.

% \begin{remark}[Why the YLG functional is nonincreasing]\label{rem:YLG-V-monotone}
% The way Yakubovich–Leonov–Gelig prove that a Lyapunov functional of the
% form~\eqref{eq:YLG-V-def} is nonincreasing along solutions can be seen
% directly from the proof of Theorem~3.1 in~\cite[pp.~112–114]{YLG}.
% For the general Lur'e system
% \begin{equation}\label{eq:YLG-3.1}
%   \dot x = P x + q\,\xi,\qquad
%   \sigma = r^{\top} x,\qquad
%   \xi \in \varphi(\sigma),
% \end{equation}
% they introduce the Lyapunov function (eq.~(3.8) in~\cite{YLG})
% \[
%   V(x)
%   \;=\;
%   x^{\top} H x
%   + \sum_{i=1}^{m}
%       \vartheta_i \int_{0}^{\sigma_i} \varphi_i(\lambda)\,d\lambda,
% \]
% where \(H\) is a symmetric matrix to be chosen, and
% \(\vartheta_i>0\) are diagonal weights. The nonlinearities
% \(\varphi_i\) are assumed to be piecewise single–valued
% (Definition~1.6 in~\cite{YLG}), which is exactly the property
% enjoyed by the Filippov convexifications \(\conv g_i\) of our
% activation functions \(g_i\in\mathcal G'\).

% The key technical step is the differentiation formula
% \begin{equation}\label{eq:YLG-3.9}
%   \frac{d}{dt}\int_{0}^{\sigma_i(t)} \varphi_i(\lambda)\,d\lambda
%   \;=\;
%   \xi_i(t)\,\dot\sigma_i(t)
%   \quad\text{for a.e. }t\ge0,
% \end{equation}
% which is Eq.~(3.9) in~\cite{YLG}.  It follows from the absolute
% continuity of \(\sigma_i(t)\) and from the fact that the integrand
% \(\varphi_i(\cdot)\) is locally bounded away from its isolated jump
% points.  Using \eqref{eq:YLG-3.9} and the system
% equations~\eqref{eq:YLG-3.1}, YLG compute the derivative of \(V\) along
% a solution \((x(t),\xi(t))\) and obtain (cf. eq.~(3.10) in
% \cite{YLG})
% \begin{equation}\label{eq:YLG-3.10}
%   \dot V(x(t))
%   \;=\;
%   W(x(t),\xi(t))
%   - \sum_{i=1}^{m} \tau_i\,\xi_i(t)\bigl(\sigma_i(t)-\mu_i\xi_i(t)\bigr),
% \end{equation}
% where the quadratic form \(W(x,\xi)\) is given in~\cite[(3.10)]{YLG}
% and \(\tau_i>0\), \(\mu_i\ge0\) are diagonal parameters.

% The sign of the right–hand side is then controlled in two pieces.
% First, the ``sector'' condition (3.4) of Theorem~3.1,
% \[
%   \bigl(\sigma_i - \mu_i\xi_i\bigr)\,\xi_i \ge 0,
%   \qquad \xi_i\in\varphi_i(\sigma_i),
% \]
% implies that each summand in the second term of
% \eqref{eq:YLG-3.10} is nonnegative, so that
% \[
%   -\sum_{i=1}^{m} \tau_i\,\xi_i(\sigma_i-\mu_i\xi_i) \;\le 0.
% \]
% Second, the frequency–domain inequality (3.2) together with
% Theorem~2.12 in~\cite{YLG} is equivalent to the existence of a
% symmetric matrix \(H\) such that the quadratic form \(W\) satisfies
% \begin{equation}\label{eq:YLG-3.11}
%   W(x,\xi)\le 0,
%   \qquad \forall\,x\in\R^n,\;\xi\in\R^m
% \end{equation}
% (eq.~(3.11) in~\cite{YLG}).  Combining \eqref{eq:YLG-3.10} and
% \eqref{eq:YLG-3.11} yields \(\dot V(x(t))\le 0\) for almost all
% \(t\), so \(V\) is nonincreasing along every solution
% of~\eqref{eq:YLG-3.1}.

% % In our Hopfield setting we work with the Filippov inclusion
% % \[
% %   \dot x(t) \in -Bx(t) + T\,\conv g(x(t)) + I,
% % \]
% % which can be rewritten in the Lur'e form~\eqref{eq:YLG-3.1} with
% % \(P=-B\), \(q=T\), \(r^{\top}=I_n\) and \(\varphi_i = \conv g_i\).
% Assumption~C1 (linear growth) and the class \(\mathcal G'\) ensure that
% each \(\varphi_i\) is piecewise single–valued with at most linear
% growth, so that the integral terms \(\Psi_i\) in
% \eqref{eq:YLG-V-def} are well defined and locally Lipschitz, and the
% differentiation formula \eqref{eq:YLG-3.9} holds for any Filippov
% solution.  Assumption~C3 (diagonal Lyapunov stability) is an
% algebraic analogue of~\eqref{eq:YLG-3.11}: it guarantees the existence
% of a diagonal matrix \(H\succ0\) such that the quadratic part of
% \(\dot V\) is negative definite in \((x,\eta)\) (see inequality
% \((26)\) below).  Finally, the one–sided Lipschitz condition~C2 is
% used later to obtain radial unboundedness of \(V\) and uniqueness of
% the equilibrium, but it does not affect the sign of \(\dot V\).

% Thus, under Assumptions~C1–C3, the same mechanism as in the proof of
% Theorem~3.1 in~\cite{YLG} applies: the chain rule
% \eqref{eq:YLG-3.9} gives an explicit expression for \(\dot V\) along
% Filippov solutions, and the matrix inequality in~C3 implies that
% \(\dot V(x(t))\le 0\) for almost all \(t\ge0\).  In particular, the
% YLG functional~\eqref{eq:YLG-V-def} is nonincreasing along every
% Filippov solution of the discontinuous Hopfield network.
% \end{remark}

% \medskip
% \subsubsection{Radial unboundedness of $V$}

% We now show that $V$ is radially unbounded, using only C1–C3 and the
% class $\mathcal G'$.

% First, since $H\succ0$, there exists $\lambda_H>0$ such that
% \begin{equation}\label{eq:H-coercive}
%     x^\top H x \;\ge\; \lambda_H \|x\|^2,
%     \qquad \forall x\in\R^n.
% \end{equation}

% Fix $i\in\{1,\dots,n\}$ and consider any measurable selection
% $\eta_i(\cdot)\in\conv g_i(\cdot)$.
% By Assumption~C2 (one-sided Lipschitz, componentwise) applied with
% $y_i=0$ and any $x_i\in\R$, we obtain
% \[
%     \bigl(\eta_i(x_i)-\eta_i(0)\bigr)x_i
%     \;\le\; C_i x_i^2.
% \]
% From this and the sign of $x_i$ it follows that, for all $s\in\R$,
% \begin{equation}\label{eq:eta-growth}
%     |\eta_i(s)|
%     \;\le\; C_i|s| + |\eta_i(0)|
%     \;\le\; \widehat C_i |s|,
% \end{equation}
% for some constant $\widehat C_i\ge0$ (depending on $C_i$ and
% $\eta_i(0)$).  Integrating~\eqref{eq:eta-growth} and again using the
% sign of~$s$, we obtain the quadratic lower bound
% \begin{equation}\label{eq:Psi-lower}
%     \Psi_i(s)
%     = \int_0^s \eta_i(\sigma)\,d\sigma
%     \;\ge\; -\frac12\,\widehat C_i s^2 - |\Psi_i(0)|,
%     \qquad \forall s\in\R.
% \end{equation}
% Summing over $i$ and writing $\|x\|^2=\sum_i x_i^2$ gives
% \begin{equation}\label{eq:sum-Psi-lower}
%     \sum_{i=1}^n \Psi_i(x_i)
%     \;\ge\; -\frac12\Bigl(\sum_{i=1}^n \widehat C_i\Bigr)\|x\|^2
%            - c_2,
%     \qquad
%     c_2 := \sum_{i=1}^n |\Psi_i(0)|.
% \end{equation}

% Combining~\eqref{eq:H-coercive} and~\eqref{eq:sum-Psi-lower} in
% \eqref{eq:YLG-V-def} yields
% \begin{equation}\label{eq:V-coercive-raw}
%     V(x)
%     \;\ge\;
%     \frac12\lambda_H\|x\|^2
%     -\frac12\Bigl(\sum_{i=1}^n \widehat C_i\Bigr)\|x\|^2
%     -c_2.
% \end{equation}
% Since C3 is homogeneous in $H$, we may scale $H$ so that
% $\lambda_H > \sum_i \widehat C_i$.
% Setting
% \[
%     c_1 := \tfrac12\bigl(\lambda_H - \sum_i \widehat C_i\bigr)>0,
% \]
% we obtain the coercivity estimate
% \begin{equation}\label{eq:V-coercive}
%     V(x) \;\ge\; c_1\|x\|^2 - c_2,
%     \qquad \forall x\in\R^n.
% \end{equation}
% Therefore $V(x)\to+\infty$ as $\|x\|\to\infty$, i.e.\ $V$ is radially
% unbounded.

% \medskip
% \subsubsection{Quasi-gradient-likeness and convergence to the equilibrium}

% Properties~\eqref{eq:Vdot-neg} and~\eqref{eq:V-coercive} show that $V$
% satisfies the hypotheses of Lemma~1.5 in YLG:
% \begin{itemize}
%     \item $V$ is continuous, nonnegative and radially unbounded;
%     \item $V(x(t))$ is nonincreasing along any Filippov solution
%           of~\eqref{eq:HNN-shifted};
%     \item if $V(x(t))\equiv\const$ along a complete trajectory, then
%           (by~\eqref{eq:Vdot-neg}) we must have $x(t)\equiv x^\ast$
%           and $\eta(t)\equiv\eta^\ast$ with
%           $0\in Bx^\ast+T\,\conv g(x^\ast)$, i.e.\ the trajectory is
%           stationary.
% \end{itemize}
% Hence, by Lemma~1.5 of YLG, the Hopfield inclusion
% \eqref{eq:HNN-shifted} is \emph{quasi-gradient-like} and all its
% Filippov trajectories are bounded on $[0,+\infty)$.

% Let $\Lambda$ denote the stationary set of~\eqref{eq:HNN-shifted}.
% In Section~V-A we have already shown that $\Lambda=\{0\}$, i.e.\ the
% equilibrium is unique.  Applying Lemma~1.7 of YLG (with the shifted
% functional $V(x-c)$) we conclude that the singleton stationary set
% $\Lambda=\{0\}$ is \emph{pointwise globally stable}.  Equivalently,
% the origin is a globally asymptotically stable equilibrium of the
% Hopfield inclusion~\eqref{eq:HNN-shifted} in the sense of Filippov:
% for every $x_0\in\R^n$ and every Filippov solution $x(\cdot)$ of
% \eqref{eq:HNN-shifted} with $x(0)=x_0$ one has
% \[
%     \lim_{t\to+\infty} x(t;x_0) = 0,
% \]
% and Lyapunov stability holds in the usual sense.

% Since \eqref{eq:HNN-shifted} is obtained from~\eqref{eq:special} by an
% equilibrium shift, this proves that under Assumptions~C1–C3 the
% original discontinuous Hopfield network~\eqref{eq:special} possesses a
% unique equilibrium which is globally asymptotically stable for all
% Filippov solutions.

% \medskip
% Now, we have everything to formulate the main statement of this work.

% \begin{theorem}[Global asymptotic stability of a Hopfield network with
% discontinuous activations]
% Consider the continuous–time Hopfield neural network
% \begin{equation}
%   \dot x(t) = -B x(t) + T\,g(x(t)) + I, \qquad x(t)\in\mathbb{R}^n,
%   \label{eq:HNN}
% \end{equation}
% where
% \begin{itemize}
%   \item $B = \mathrm{diag}(b_1,\dots,b_n)\succ 0$ is a diagonal matrix
%         of leak rates;
%   \item $T\in\mathbb{R}^{n\times n}$ is the interconnection (synaptic) matrix;
%   \item $I\in\mathbb{R}^n$ is a constant input vector;
%   \item $g=(g_1,\dots,g_n)^{\top}:\mathbb{R}^n\to\mathbb{R}^n$ is the
%         componentwise activation mapping, each $g_i:\mathbb{R}\to\mathbb{R}$
%         possibly discontinuous.
% \end{itemize}
% Assume that $g$ belongs to the extended class $\mathcal{G}'$ of activation
% functions (Definition~1) and satisfies the following conditions.

% \begin{description}
%   \item[C1 (Linear growth).]
%   For each $i\in\{1,\dots,n\}$ there exist constants $a_i,b_i\ge 0$
%   such that
%   \[
%       |g_i(\xi)| \le a_i|\xi| + b_i, \qquad \forall\,\xi\in\mathbb{R}.
%   \]

%   \item[C2 (One–sided Lipschitz, componentwise).]
%   For each $i\in\{1,\dots,n\}$ there exists a constant $C_i\in\mathbb{R}$
%   such that for all $x,y\in\mathbb{R}^n$ and for all selections
%   $\eta\in\mathrm{conv}\,g(x)$, $\zeta\in\mathrm{conv}\,g(y)$ one has
%   \[
%     (\eta_i-\zeta_i)(x_i-y_i) \le C_i (x_i-y_i)^2 .
%   \]

%   \item[C3 (Diagonal Lyapunov stability of the linear part).]
%   There exist a diagonal matrix $H=\mathrm{diag}(h_1,\dots,h_n)\succ0$
%   and a constant $\alpha\ge0$ such that
%   \[
%     -HB - (HB)^{\top} + \alpha\,(T^{\top}H + HT) \prec 0 .
%   \]
% \end{description}

% Interpret system~\eqref{eq:HNN} in the sense of Filippov, i.e. as the
% differential inclusion
% \[
%    \dot x(t) \in -B x(t) + T\,\varphi(x(t)) + I,
% \]
% where $\varphi(x)=\mathrm{conv}\,g(x)$ is the Filippov convexification
% of $g$.

% Then the following properties hold.
% \begin{enumerate}
%   \item For every initial condition $x(0)=x_0\in\mathbb{R}^n$ there exists
%         at least one Filippov solution $x(\cdot;x_0)$ of~\eqref{eq:HNN}
%         defined on the whole half–line $[0,+\infty)$.

%   \item There exists a unique equilibrium point $e\in\mathbb{R}^n$,
%         i.e. a constant vector satisfying
%         \[
%            0 \in -B e + T\,\mathrm{conv}\,g(e) + I .
%         \]

%   \item This equilibrium is \emph{globally asymptotically stable in the
%         sense of Filippov}: for every Filippov solution
%         $x(\cdot;x_0)$ of~\eqref{eq:HNN} one has
%         \[
%            \lim_{t\to+\infty} x(t;x_0) = e ,
%         \]
%         and for every $\varepsilon>0$ there exists $\delta>0$ such that
%         $|x_0-e|<\delta$ implies $|x(t;x_0)-e|<\varepsilon$ for all
%         $t\ge0$.
% \end{enumerate}
% \end{theorem}


%%%%%%%%%%%%%%%%%%%%%%%%%%%%%%%%%%%%%%%%%%%%%%
%%%%%%%%%%%%%%%%%%%%%%%%%%%%%%%%%%%%%%%%%%%%%%

% \begin{IEEEproof}

% Summarizing all written above in one line:
% \begin{itemize}
%   \item C1 $\Rightarrow$ linear growth of $g$ and
%         $\Phi(x)\to\infty$ as $\|x\|\to\infty$;
%   \item C2 + class $\mathcal{G}'$ $\Rightarrow$ piecewise smooth
%         gradient nonlinearity, admissible for the YLG analysis in
%         §3.3;
%   \item C3 $\Rightarrow$ Yakubovich–type frequency inequalities for the
%         Lur'e pair $(B,T)$;
%   \item Theorem~3.14 (YLG) $\Rightarrow$ quasi-gradient-like behavior:
%         every trajectory of the Filippov inclusion tends to the
%         stationary set;
%   \item uniqueness of equilibrium (Section~V-A)
%         $\Rightarrow$ stationary set is $\{e\}$;
%   \item Lemmas~1.4--1.7 (YLG) $\Rightarrow$ this singleton stationary
%         set is pointwise globally stable.
% \end{itemize}

% Therefore, the Hopfield network with activation functions satisfying
% C1--C3 has a \emph{unique equilibrium}, and this equilibrium is
% \emph{globally asymptotically stable} for all Filippov solutions of the
% discontinuous system.
% \end{IEEEproof}

% \newpage
% \subsection{Lyapunov function as quadratic form and integral of nonlinearity}

% % Let $H=\alpha^{-1}\succ0$ and define
% % \begin{equation}\label{eq:lyap}
% %   V(x)=\frac12 x^{\!*}Hx+\sum_{i=1}^{n}\int_{0}^{x_i}\!\eta_i(\sigma)\,d\sigma,
% %   \qquad\eta_i(\sigma)\in\varphi_i(\sigma).
% % \end{equation}
% % The integrals are well defined by C1 and $V$ is radially unbounded.
% % For a Filippov trajectory of~\eqref{eq:lure} we have, using the chain rule
% % for regular functions~\cite{p9},
% % \[
% %  \dot V=x^{\!*}(HB+BH)x+\xi^{\!*}M\xi\leqslant0,
% % \]
% % because $HB+BH=-2H\operatorname{diag}(b_i)H\prec0$ and $M\prec0$ from C3.  Equality is
% % possible only when $x=0$ and $\xi=0$.

% Choose $H:=\alpha^{-1}\succ0$ and define the YLG Lyapunov function
% \begin{equation}\label{eq:Vdef}
% \begin{aligned}
%     &V(x)=\frac12 x^{\!*}Hx+\sum_{i=1}^{n}\Psi_i(x_i), \\
%     &\textbf{}\Psi_i(s):=\int\limits_{0}^{s}\eta_i(\sigma)\,d\sigma,\;\eta_i(\sigma)\in\varphi_i(\sigma).
% \end{aligned}
% \end{equation}
% By \textbf{C1} the integral grows at least quadratically, hence $V$ is radially
% unbounded ($V(x)$ tends to $+\infty$ if $\|x\|$ tends to $+\infty$). Also note that this function is regular, being a convex nonsmooth function in $\mathbb{R}^n$ \cite{516}. Differentiating along a Filippov solution of \eqref{eq:lure}
% (cf.\ the chain rule for regular functions \cite{p9}) yields
% \begin{align}\label{eq:Vdot}
%   \dot V &= x^{\!*}H\dot x+\xi^{\!*}\dot x
%           = x^{\!*}H(Px+q\xi)+\xi^{\!*}(Px+q\xi)\\
%          &= x^{\!*}(HP+P^{\!*}H)x+\xi^{\!*}(\alpha T+T^{\!*}\alpha)\xi.
% \end{align}
% Because $HP+P^{\!*}H=-2H\operatorname{diag}(b_i)H\prec0$ and $M~\prec~0$ by \textbf{C3}, we have
% \begin{equation}\label{eq:negVdot}
%   \dot V \leqslant -\lambda_1\|x\|^{2}-\lambda_2\|\xi\|^{2},\qquad \lambda_{1,2}>0.
% \end{equation}
% Hence $V$ is non‑increasing along every solution.

% \subsection{Global Asymptotic Stability}\label{sec:GAS}
% % \begin{theorem}[GAS]\label{thm:gas}
% % Under \textnormal{C1--C3}, the equilibrium $x^\star$ of the HNN is globally asymptotically stable.
% % \end{theorem}
% % \begin{IEEEproof}[Proof sketch]
% % Invoke a YLG functional $V(x)=x^{\top}Hx + \sum_i \int_{0}^{x_i}\phi_i(\tau)\,d\tau$ with appropriate $\phi_i$ associated to $g_i$, and show $\dot V\le -\lambda \|x-x^\star\|^2$ along Filippov trajectories. Full derivations appear in Appendix~\ref{app:proofs}.
% % \end{IEEEproof}
% % Because $\dot V\leqslant0$ and the set where $\dot V=0$ reduces to the singleton
% % $\{x^{\!*}\}$, LaSalle’s invariance principle for differential inclusions
% % (\cite[Cor.~4.2]{p223} together with \cite[\S1.3]{Leonov}) implies
% % \(\displaystyle\lim_{t\to\infty}x(t)=x^{\!*}\).  Thus the unique equilibrium
% % is globally asymptotically stable.
% \subsubsection{Boundedness of trajectories.}\label{ssec:bounded}

% In this subsection we prove that \emph{every} Filippov solution of the
% Hopfield inclusion \eqref{eq:lure} is bounded on $[0,\infty)$.  
% The proof splits into two lemmas.

% \begin{lemma}[radial unboundedness of $V$]\label{lem:radial}
% There exist positive constants $c_1,c_2$ such that for all $x\in\mathbb{R}^n$
% \begin{equation}\label{eq:radial}
%    V(x)\;\geqslant\; c_1\,\|x\|^{2}-c_2.
% \end{equation}
% Consequently $V(x)\to\infty$ as $\|x\|\to\infty$.
% \end{lemma}
% \begin{IEEEproof}
% Write $x=(x_1,\dots,x_n)$. By definition \eqref{eq:Vdef}
% \(V(x)=\frac12 x^{\!*}Hx+\sum_i \Psi_i(x_i)\).  Because $H\succ0$ there is
% $\lambda_H>0$ with $x^{\!*}Hx\geqslant \lambda_H\|x\|^{2}$.  
% It remains to lower bound $\Psi_i$.

% For each $i$ fix any measurable selection $\eta_i\in\varphi_i$.  
% By condition
% \textbf{C2} the one--sided Lipschitz property implies $|\eta_i(s)|\leqslant C_i|s|+|\eta_i(0)|$.
% Hence for $|s|\geqslant1$
% \(|\eta_i(s)|\le\hat C_i|s|\) with $\hat C_i:=C_i+|\eta_i(0)|$.
% Using the sign of $s$ we get
% \(
%   \Psi_i(s)=\int\limits_{0}^{s}\eta_i(\sigma)\,d\sigma\geqslant-\tfrac12 \hat C_i s^{2}.
% \)
% Summing over $i$ gives
% \[\sum_{i}\Psi_i(x_i)\geqslant-\tfrac12\bigl(\sum_i\hat C_i\bigr)\,\|x\|^{2}.
% \]
% Set $c_1:=\frac12\bigl(\lambda_H-\sum_i\hat C_i\bigr)$ (positive because $H$
% can be scaled if necessary) and $c_2:=\sum_i |\Psi_i(0)|$ to obtain~\eqref{eq:radial}.
% \end{IEEEproof}

% \begin{lemma}[non--increase of $V$]\label{lem:noninc}
% Along any Filippov solution $t\mapsto x(t)$ of \eqref{eq:lure}
% \begin{equation*}\label{eq:noninc}
%   t\mapsto V\bigl(x(t)\bigr)\quad\text{is non--increasing on }[0,\infty).
% \end{equation*}
% \end{lemma}
% \begin{IEEEproof}
% This is immediate from inequality \eqref{eq:negVdot}, which shows
% $\dot V\leqslant0$ for almost every~$t$.
% \end{IEEEproof}

% \begin{proposition}[boundedness]\label{prop:bounded}
% Let $x(\cdot)$ be a Filippov solution of \eqref{eq:lure}.  Then there exists
% $R>0$ (depending on $x(0)$) such that $\|x(t)\|\leqslant R$ for all $t\geqslant0$.
% \end{proposition}
% \begin{IEEEproof}
% Combine Lemma~\ref{lem:noninc} with Lemma~\ref{lem:radial}: for all
% $t\geqslant0$
% \[
%    V\bigl(x(t)\bigr)\leqslant V\bigl(x(0)\bigr)=:V_0.
% \]
% From \eqref{eq:radial} we deduce $c_1\|x(t)\|^{2}-c_2\leqslant V_0$, hence
% \(\|x(t)\|^{2}\leqslant (V_0+c_2)/c_1\).  Setting
% $R:=\sqrt{(V_0+c_2)/c_1}$ completes the proof.
% \end{IEEEproof}

% \begin{remark}
% Proposition~\ref{prop:bounded} 
% is the desired boundedness statement.  
% Note that no trajectory--specific information besides the initial energy $V_0$ enters the bound.
% \end{remark}


% % Inequality \eqref{eq:negVdot} and radial unboundedness of $V$ imply
% % \(V(x(t))\leqslant V(x(0))<\infty\), so $x(t)$ stays in the compact sublevel set
% % $\{V\leqslant V(x(0))\}$ for all $t\geqslant0$.  Therefore every Filippov trajectory is
% % bounded.

% \subsubsection{Invariance principle for inclusions.}

% We now show rigorously that every bounded Filippov trajectory converges to
% the equilibrium by applying the LaSalle--type dichotomy lemma of
% YLG \cite[Lem.~1.4]{Leonov}. One can underline the similarity with the Appendix B, where LaSalle invariance principle, modified for regular Lyapunov functions, was used.
% \paragraph{1.  The zero–derivative set.}
% Define the set of points where the Lyapunov derivative vanishes:
% \[
%  \mathcal E\;:=\;\bigl\{(x,\xi)\,\big|\, \dot V(x,\xi)=0\bigr\}.
% \]
% Using \eqref{eq:Vdot} we have
% \(
%  \dot V= x^{\!*}(HP+P^{\!*}H)x+\xi^{\!*}M\xi.
% \)
% Because $HP+P^{\!*}H\prec0$ and $M\prec0$, the quadratic form is the sum of
% two \emph{strictly negative} definite terms.  Consequently
% \[\dot V(x,\xi)=0\quad\Longleftrightarrow\quad x=0\;\text{and}\;\xi=0.
% \]
% Thus
% \begin{equation}\label{eq:E}
%   \boxed{\;\mathcal E=\{(0,0)\}.\;}
% \end{equation}

% \paragraph{2.  $\omega$--limit sets lie in $\mathcal E$.}
% Let $z(t)=(x(t),\xi(t))$ be any bounded Filippov solution of
% \eqref{eq:lure}.  Since $V$ is non‑increasing and bounded below, the limit
% $V_{\infty}:=\lim_{t\to\infty}V(z(t))$ exists.  Take an arbitrary sequence
% $t_k\to\infty$ such that $z(t_k)\to \bar z$.  Continuity of $V$ yields
% $V(\bar z)=V_{\infty}$.  Moreover,
% \[
%   0\leqslant \int\limits_{t_k}^{t_k+\tau}\!(-\dot V)\,dt=V(z(t_k))-V(z(t_k+\tau))\xrightarrow{k\to\infty}0
% \]
% for any fixed $\tau>0$.  
% Therefore $\dot V(z(t))\to0$ as $t\to\infty$ (in the
% \emph{Cesàro} sense).  Passing to a subsequence if necessary we may assume
% $\dot V(z(t_k))\to0$, hence $\bar z\in\mathcal E$.  This shows that every
% $\omega$–limit point of $z(\cdot)$ lies in $\mathcal E$.

% \paragraph{3.  Maximal weakly invariant subset of $\mathcal E$.}
% A set $S$ is \emph{weakly invariant} for a differential inclusion if every
% point of $S$ lies on a complete trajectory remaining in $S$.  Because
% $\mathcal E$ in \eqref{eq:E} is a single point, it is trivially weakly
% invariant and is the \emph{largest} such subset contained in itself.
% Applying the dichotomy lemma of YLG gives
% \[\omega\bigl(z(\cdot)\bigr)\subseteq\mathcal E.
% \]

% \paragraph{4.  Convergence of the state.}
% Since $\mathcal E$ projects onto $x=0$ in the state space, we have
% $\lim_{t\to\infty}x(t)=0$.  Denote this equilibrium by $e$.  Hence
% \begin{equation}\label{eq:conv}
%   \boxed{\;\displaystyle\lim_{t\to\infty}x(t)=e.\;}
% \end{equation}
% That is, the equilibrium is \emph{attractive}.

% % Let $\mathcal{E}:=\{(x,\xi)\st \dot V=0\}$, i.e.
% % $Hx=0$ and $\xi\in\ker M$.  Because $H\succ0$ we immediately get $x=0$.
% % Since $M\prec0$, its kernel is trivial, whence $\xi=0$.  Thus
% % \begin{equation}\label{eq:invariant-set}
% %   \mathcal{E}=\{(0,0)\}.
% % \end{equation}

% % By YLG Lemma~1.4 (dichotomy lemma) the \emph{\(\omega\)-limit set} of any
% % bounded trajectory is contained in the largest weakly invariant subset of
% % $\mathcal{E}$.  Via \eqref{eq:invariant-set} this subset is the singleton
% % $\{e\}$, where $e$ is the state satisfying $x=0$ and
% % \eqref{eq:lure}.  Consequently
% % \[\lim_{t\to\infty}x(t)=e.\]

% \subsubsection{Lyapunov stability.}

% Because $V$ is continuous and strictly increasing away from $e$, it
% serves as a Lyapunov function in the sense of Fillipov (see
% \cite[Ch.~1,\,§3]{Leonov}).  
% Any initial condition taken sufficiently close to
% $e$ remains close for all $t\geqslant0$.  
% Let us justify Lyapunov stability in detail in five explicit steps.

% \paragraph{Preliminaries.}  Recall the quadratic bound derived from the
% negative--definite inequality \eqref{eq:negVdot}:
% \begin{equation}\label{eq:quad-bound}
%   \dot V(x,\xi)\leqslant -\lambda_1\|x\|^{2}-\lambda_2\|\xi\|^{2}, \qquad
%   \lambda_{1,2}>0.
% \end{equation}
% Furthermore, by positive–definiteness of $H$ and linear growth of the
% integrals in $V$ there exist positive constants $c_{\min},c_{\max}$ such
% that
% \begin{equation}\label{eq:norm-to-V}
%   c_{\min}\|x\|^{2}\leqslant V(x)\leqslant c_{\max}\|x\|^{2}\quad(\forall x\in\mathbb{R}^{n}).
% \end{equation}

% \paragraph{Step 1.  Choose an energy radius.}  Fix an arbitrary
% $\varepsilon>0$.  Define
% \[
%   \gamma:=\frac12\min\{\lambda_1,\lambda_2\}\,\varepsilon^{2}.
% \]
% This is the maximal acceptable energy level for Lyapunov stability.

% \paragraph{Step 2.  Select $\delta$ via continuity of $V$.}  Because $V$ is
% continuous and $V(e)=0$, there exists $\delta>0$ such that
% $\|x-e\|<\delta$ implies $V(x)<\gamma$.

% \paragraph{Step 3.  Show invariance of the $V\le\gamma$ sublevel set.}
% Let $x(0)$ satisfy $\|x(0)-e\|<\delta$ so that $V(x(0))<\gamma$.
% Inequality \eqref{eq:quad-bound} gives $\dot V\leqslant0$ along every Filippov
% trajectory.  Therefore $V(x(t))\leqslant V(x(0))<\gamma$ for \emph{all} $t\geqslant0$.

% \paragraph{Step 4.  Translate the energy bound back to the state radius.}
% Applying the left inequality of \eqref{eq:norm-to-V} we obtain
% \[\|x(t)-e\|^{2}\leqslant \frac{V(x(t))}{c_{\min}}< \frac{\gamma}{c_{\min}}
%   =\frac{\min\{\lambda_1,\lambda_2\}}{2c_{\min}}\,\varepsilon^{2}.
% \]
% Set $k:=\sqrt{\min\{\lambda_1,\lambda_2\}/(2c_{\min})}\,>0$.  Then
% \(\|x(t)-e\|<k\varepsilon.\)  Since $k\leqslant1$ (because $c_{\min}$ may be
% chosen as $\lambda_1/2$), we have the stronger estimate
% \[\|x(t)-e\|<\varepsilon\quad\forall t\geqslant0.
% \]

% \paragraph{Step 5.  Conclude Lyapunov stability.}  We have shown that for
% any $\varepsilon>0$ there exists $\delta>0$ such that
% $\|x(0)-e\|<\delta$ implies $\|x(t)-e\|<\varepsilon$ for all
% $t\geqslant0$.  This is precisely the definition of Lyapunov stability of
% $e$.
% Combining Lyapunov stability with the attractivity result~\eqref{eq:conv} establishes \emph{global asymptotic stability} of the
% unique equilibrium.

% \subsection{Relation to the \textbf{C1--C3} conditions}\label{sec:rel}

% We now unpack how hypotheses \textbf{C1--C3} enforce every requirement
% used in Section \ref{sec:GAS}.

% \begin{enumerate}
%   \item \textbf{Existence of Filippov solutions.}  \textbf{C1} ensures linear growth
%     of $g$; together with local Lipschitz continuity of the linear part it
%     satisfies the Filippov existence theorem \cite[Th.~1,\,p.~77]{Leonov}.
%   \item \textbf{Upper semicontinuity and convexity.}  The extended
%     non‑linearity $\varphi$ inherits bounded closed convex values from the
%     convex hull; upper semicontinuity follows from finite jumps.  This matches
%     YLG assumptions~(i)--(ii).
%   \item \textbf{Negative definiteness of $M$.}  The first part of \textbf{C3} is
%     exactly the frequency‑domain negativity condition of YLG Theorem~3.15 for
%     gradient non‑linearities, guaranteeing \emph{dissipation} via the second
%     term in \eqref{eq:Vdot}.
%   \item \textbf{Small–gain inequality.}  The second clause in \textbf{C3} bounds the
%     contribution of the cross term in the equilibrium–uniqueness argument
%     (see \eqref{eq:uniq} above).  It turns the one‑sided Lipschitz estimate
%     from \textbf{C2} into a strict inequality.
% \end{enumerate}

% \subsubsection{Summary.}
%     The classical Yakubovich--Leonov--Gelig machinery provides a concise proof
% of the uniqueness and global stability of Hopfield equilibria under the
% non‑smooth conditions \textbf{C1--C3}.  Besides simplifying the argument, the method
% highlights how frequency–domain negativity of \(\alpha T+T^{\!*}\alpha\)
% serves as the decisive criterion, irrespective of boundedness or monotonicity
% of activation functions.

% \subsection{Tightness and Conservatism}
% State small counterexamples or boundary cases indicating near-necessity of parts of C1--C3.
\\ 

\todo[inline]{NEW BELOW}
\section{Nonunique Filippov trajectories:
A Low–Dimensional Example}\label{sec:nonunique}

In this section, we present a low-dimensional example demonstrating that, even when the standard assumptions \textbf{C1--C3} (linear growth, extended nonlinearity structure, and the YLG-type frequency condition) are satisfied, the corresponding Filippov inclusion may admit multiple distinct forward trajectories originating from the same initial state. This demonstrates that global asymptotic stability of a unique equilibrium does not, in general, guarantee forward uniqueness of Filippov solutions. Throughout this section and Section~\ref{sec:examples}, we will employ the diagonal Lyapunov stability condition \textbf{C3$'$} in place of the original frequency-domain criterion \textbf{C3}. Since the matrix inequality is rigorously implied by the spectral condition (Lemma~\ref{lem:C3prime-implies-C3}), this substitution preserves all theoretical guarantees while providing a more transparent algebraic framework for the examples under consideration.

Consider a two-dimensional Hopfield network
\begin{equation}\label{eq:example-system}
\begin{cases}
\dot{x}_{1}=-x_{1},\\
\dot{x}_{2}=-x_{2}+a\,\operatorname{sign}(x_{1}),
\end{cases}
\qquad x=(x_{1},x_{2})^{\top}\in\mathbb{R}^{2},
\end{equation}

In this example we take:
\begin{equation}\label{eq:choice}
B=I_{2},\qquad T=\begin{pmatrix}0&0\\ a&0\end{pmatrix},\qquad I=0,
\end{equation}
with a real parameter $a\in\mathbb{R}$, and discontinuous  activation:
\begin{equation}\label{eq:activation}
g_{1}(s)=\operatorname{sign}(s),\qquad g_{2}(s)\equiv 0,
\end{equation}
where 
\[
\operatorname{sign}(s)=
\begin{cases}
-1,&s<0,\\
0,&s=0,\\
+1,&s>0.
\end{cases}
\]

The system \eqref{eq:example-system} satisfies conditions \textbf{C1--C3} for $|a|<2$:
\begin{itemize}
\item \textbf{C1 (linear growth):} $|g_{1}(\xi)|\leq 1$, $g_{2}\equiv 0$.
\item \textbf{C2 (extended nonlinearity):} $g_{1}$ is a selection of the subdifferential $\partial\zeta_{1}$ of the convex function $\zeta_{1}(s)=|s|$; $g_{2}\equiv 0$ corresponds to $\zeta_{2}\equiv 0$.
\item \textbf{C3 $\Rightarrow$ C3$'$ (diagonal stability):} With $H=I_{2}$ and $\alpha=1$, the matrix inequality
\[
-2I_{2}+(T+T^{\top})\prec 0
\]
holds when $|a|<2$.
\end{itemize}

Due to the discontinuity, we interpret \eqref{eq:example-system} in the Filippov sense. The regularized inclusion reads
\begin{equation}\label{eq:filinc-example}
\begin{cases}
\dot{x}_{1}=-x_{1},\\
\dot{x}_{2}\in -x_{2}+a\,G_{1}(x_{1}),
\end{cases}
\end{equation}
where 
\[
G_{1}(x_{1})=
\begin{cases}
\{-1\},&x_{1}<0,\\
[-1,1],&x_{1}=0,\\
\{+1\},&x_{1}>0,
\end{cases}
\]
and $G_{2}(x_{2})\equiv\{0\}$.

The first equation in \eqref{eq:filinc-example} is a classical linear ODE: for any $x_{1,0}\in\mathbb{R}$
\[
x_{1}(t)=x_{1,0}e^{-t}
\]
is the unique solution. In particular, if $x_{1,0}=0$, then
\[
x_{1}(t)\equiv 0\quad\text{for all }t\geq 0.
\]
So for the initial condition $x(0)=(0,0)$, the first equation gives $x_{1}(t)\equiv 0$. 

Substituting $x_{1}(t)=0$ into the second inclusion in \eqref{eq:filinc-example} yields
\[
G_{1}(x_{1}(t))=G_{1}(0)=[-1,1],
\]
and the dynamics of $x_{2}$ reduce to the scalar differential inclusion
\begin{equation}\label{eq:scalarinc}
\dot{x}_{2}(t)\in -x_{2}(t)+a\,[-1,1]=\bigl[-x_{2}(t)-a,\,-x_{2}(t)+a\bigr].
\end{equation}
Thus, along the invariant line $\{x_{1}=0\}$ the 2D Filippov system \eqref{eq:filinc-example} reduces to the 1D inclusion \eqref{eq:scalarinc}.

Now we will show that for the initial condition $x(0)=(0,0)$ there exist distinct Filippov solutions of \eqref{eq:example-system}, i.e., forward uniqueness fails.

Fix $a\in(0,2)$ and consider the initial condition
\[
x(0)=(0,0).
\]
As noted above, the first coordinate is uniquely determined:
\[
x_{1}(t)\equiv 0\quad\text{for all }t\geq 0.
\]
Hence any Filippov solution with this initial condition must satisfy \eqref{eq:scalarinc} for $x_{2}(t)$.

\medskip
\textbf{Solution 1: upper branch.}
Choose the upper endpoint of the interval in \eqref{eq:scalarinc}, i.e. enforce
\[
\dot{x}_{2}(t)=-x_{2}(t)+a.
\]
Solving this ODE with $x_{2}(0)=0$ gives
\[
x_{2}^{(1)}(t)=a\bigl(1-e^{-t}\bigr),\qquad t\geq 0.
\]
Then
\[
\dot{x}_{2}^{(1)}(t)=-x_{2}^{(1)}(t)+a\in\bigl[-x_{2}^{(1)}(t)-a,\,-x_{2}^{(1)}(t)+a\bigr],
\]
so $(x_{1}^{(1)}(t),x_{2}^{(1)}(t))=(0,a(1-e^{-t}))$ is a Filippov solution of \eqref{eq:example-system} with $x^{(1)}(0)=(0,0)$.

\medskip
\textbf{Solution 2: lower branch.}
Choose the lower endpoint of the interval in \eqref{eq:scalarinc}, i.e. enforce
\[
\dot{x}_{2}(t)=-x_{2}(t)-a.
\]
Solving with $x_{2}(0)=0$ gives
\[
x_{2}^{(2)}(t)=-a\bigl(1-e^{-t}\bigr),\qquad t\geq 0.
\]
Then
\[
\dot{x}_{2}^{(2)}(t)=-x_{2}^{(2)}(t)-a\in\bigl[-x_{2}^{(2)}(t)-a,\,-x_{2}^{(2)}(t)+a\bigr],
\]
and $(x_{1}^{(2)}(t),x_{2}^{(2)}(t))=(0,-a(1-e^{-t}))$ is also a Filippov solution of \eqref{eq:example-system} with $x^{(2)}(0)=(0,0)$.

\begin{proposition}[Nonuniqueness of forward trajectories]\label{eq:example-system}
For any $a\in(0,2)$, system \eqref{eq:example-system} has at least two distinct Filippov solutions $x^{(1)}(\cdot)$ and $x^{(2)}(\cdot)$ defined on $[0,+\infty)$ with the same initial condition $x^{(1)}(0)=x^{(2)}(0)=(0,0)$.
\end{proposition}
\begin{IEEEproof}
See Appendix \ref{prop3_proof}.
\end{IEEEproof}

\begin{remark}
In fact, there exists a continuum of Filippov solutions starting from $(0,0)$. If we consider
\[
\dot{x}_{2}(t)=-x_{2}(t)+a\,\mu(t),\qquad x_{2}(0)=0.
\]
The unique solution of this linear ODE is
\[
x_{2}(t)=e^{-t}\int\limits_{0}^{t}e^{s}\,a\,\mu(s)\,ds,
\]
which satisfies \eqref{eq:scalarinc} almost everywhere. For each such $\mu$ we obtain a distinct Filippov solution $x(t)=(0,x_{2}(t))^{\top}$ of \eqref{eq:example-system} with $x(0)=(0,0)$.
\end{remark}


Thus, even though the YLG conditions \textbf{C1--C3} guarantee global asymptotic stability of the unique equilibrium (here the origin), they do not ensure forward uniqueness of Filippov trajectories. This example highlights that the regularity required for a well-posed semiflow is stronger than that needed for stability alone.
\todo[inline]{NEW ABOVE}

\section{Challenges in Numerical Simulation of HNNs with Discontinuous Activations}\label{sec:examples}

% In this section we will translate four optimization tasks into the
% discontinuous HNN template.

Let us reinforce the analytical results of our work 
(and Theorem 3 from Appendix \ref{app:proofs}) with an example.
Consider the system of differential equations:
\begin{equation}\label{example1}
\begin{cases}\frac{\mathbf{d}x_1}{\mathbf{d}{t}}=-x_1-\frac 1 4 g_1(x_1)-\frac 1 8 g_2(x_2),\\\frac{\mathbf{d}x_2}{\mathbf{d}{t}}=-x_2+\frac 1 8 g_1(x_1)-\frac 1 4 g_2(x_2).\end{cases}
\end{equation}

where the activation functions are as follows:
% \begin{equation*}
% f_i(x_i) = \begin{cases}\tanh(x) -x+1,&x>0,\\ \tanh(x)+x-1, &x\leqslant0,\end{cases} ~~~i\in 1:2.
%  \end{equation*}
\begin{equation*}
    g_i(x_i) = \begin{cases}1-x_i,&x_i>0,\\ e^{x_i}-1, &x_i\leqslant0,\end{cases} \quad i\in 1:2.
\end{equation*}

This function has a first-order discontinuity and is neither bounded nor monotonic (see Fig. \ref{fig2:left}).
Since the main conditions are satisfied, the origin $(0,0)$
is a globally asymptotically stable equilibrium point,
as can be seen from the phase portrait that we obtained by modeling system $(\ref{example1})$ using Python (see Fig. \ref{fig2:right}).

% \todo[inline]{New below}

% Next example aims to present a system with an activation function $h(x)$ (Fig. \ref{fig3:left}) that does not belong to the class $\mathcal{G'}$ (Definition \ref{G_n}), since condition 1 from the definition, which imposes the relation $h_i(0^+) > h_i(0^-)$ on the left-hand and right-hand limit values of the function, is clearly not satisfied. In the right half-plane for $x$, the function grows exponentially, meaning assumption C2 is also violated.

% \begin{equation*}
%     h_i(x_i) = \begin{cases} e^{x_i}-1, &x_i>0, \\ 
%     1-x_i, & 
%     x_i\leqslant 0, \end{cases} \quad i\in 1:2.
% \end{equation*}


Now let’s take function $s(x)$ (Fig. \ref{fig4:left}) as activation for the system $(\ref{example1})$. Assumptions \textbf{C1} (linear growth) and \textbf{C2} (unilateral Lipschitz-like condition) fail because, as stated in~\cite{unilateral_Lipschitz_like_condition}, this function possesses super linear growth and does not comply with unilateral Lipschitz-like condition. The restriction condition from definition \ref{G_n} is also violated $s_i((-1)^+) = 0.01 < 3 = s_i((-1)^-)$ and therefore this function does not belong to the class $\mathcal{G'}$ (Definition \ref{G_n}). 

\begin{equation*}
    s_i(x_i) = \begin{cases} 0.01, & |x_i| \leqslant 1, \\ 
    x_i^4+2, & 
    |x_i| > 1, \end{cases} \quad i\in 1:2.
\end{equation*}


Consequently, such a system fails to satisfy all the assumptions stated in section~\ref{sec:problem} and therefore topological structure of the phase space is drastically changed and this system no longer has a unique globally asymptotically stable equilibrium point., as is evident from the phase portrait (Fig. \ref{fig4:right}).

Lastly, let's present the counterexample to assumption \textbf{C3$'$} (DLS) by taking a skew-symmetric matrix $T$ ($T^{\mathrm{T}} = -T$):

\begin{equation}\label{example2}
\begin{cases}\frac{\mathbf{d}x_1}{\mathbf{d}{t}}=-x_1- s_2(x_2),\\\frac{\mathbf{d}x_2}{\mathbf{d}{t}}=-x_2 + s_1(x_1.\end{cases}
\end{equation}

Fig. \ref{fig5:left} shows that even though activation function is viable, assumptions on the coefficients of the system are also very important in establishing global asymptotic stability.

% \begin{figure*}[!ht]
% \begin{minipage}{0.49\linewidth}
% \includegraphics[width=1.05\textwidth]{g(x).png}
%   \caption{Activation function graph $g(x)$ for the network \eqref{eq:special}, example}\label{f_i}
% \end{minipage}
% \hfill
% \begin{minipage}{0.49\linewidth}
%  \includegraphics[width=1.05\textwidth]{draft_with_another_func.png}
%   \caption{Phase portrait of the differential equation system from Example 1.}\label{fig2:right}
% \end{minipage}
% \end{figure*}

\begin{figure*}[h]
\begin{minipage}{0.49\linewidth}
\includegraphics[width=1.05\textwidth]{Figures/g(x).png}
  \caption{Graph for the activation function $g(x)$.}\label{fig2:left}
\end{minipage}
\hfill
\begin{minipage}{0.49\linewidth}
 \includegraphics[width=1.05\textwidth]{draft2_with_another_func.png}
  \caption{Phase portrait of the system $(\ref{example1})$ with $g(x)$ activation.}\label{fig2:right}
\end{minipage}
\end{figure*}

% \begin{figure*}[h]
% \begin{minipage}{0.49\linewidth}
% \includegraphics[width=1.05\textwidth]{Figures/h(x).png}
%   \caption{Graph for the activation function $h(x)$.}\label{fig3:left}
% \end{minipage}
% \hfill
% \begin{minipage}{0.49\linewidth}
%  \includegraphics[width=1.05\textwidth]{Figures/h(x) good sys.png}
%   \caption{Phase portrait of the system $(\ref{example1})$ with $h(x)$ activation.}\label{fig3:right}
% \end{minipage}
% \end{figure*}

\begin{figure*}[h]
\begin{minipage}{0.49\linewidth}
\includegraphics[width=1.05\textwidth]{Figures/s(x).png}
  \caption{Graph for the activation function $s(x)$.}\label{fig4:left}
\end{minipage}
\hfill
\begin{minipage}{0.49\linewidth}
 \includegraphics[width=1.05\textwidth]{Figures/f_1(x) good sys.png}
  \caption{Phase portrait of the system $(\ref{example1})$ with $s(x)$ activation.}\label{fig4:right}
\end{minipage}
\end{figure*}

\begin{figure*}[h]
\begin{minipage}{0.49\linewidth}
\includegraphics[width=1.05\textwidth]{Figures/g(x) bad sys.png}
  \caption{Phase portrait of the system $(\ref{example2})$ with $g(x)$ activation.}\label{fig5:left}
\end{minipage}
\hfill
\begin{minipage}{0.49\linewidth}
 \includegraphics[width=1.05\textwidth]{Figures/f_1(x) bad sys.png}
  \caption{Phase portrait of the system $(\ref{example2})$ with $s(x)$ activation.}\label{fig5:right}
\end{minipage}
\end{figure*}

% \begin{figure*}[!ht]
% \begin{minipage}{0.49\linewidth}
% \includegraphics[width=1.05\textwidth]{Figures/f(x).png}
%   \caption{Activation function graph $g(x)$ for the network \eqref{eq:special}, example}\label{f_i}
% \end{minipage}
% \hfill
% \begin{minipage}{0.49\linewidth}
%  \includegraphics[width=1.05\textwidth]{draft.png}
%   \caption{Phase portrait of the differential equation system from Example 1.}\label{fig2:right}
% \end{minipage}
% \end{figure*}

% \begin{figure*}[!ht]
% \begin{minipage}{0.49\linewidth}
% \includegraphics[width=1.05\textwidth]{Figures/f(x).png}
%   \caption{Activation function graph $g(x)$ for the network \eqref{eq:special}, example}\label{f_i}
% \end{minipage}
% \hfill
% \begin{minipage}{0.49\linewidth}
%  \includegraphics[width=1.05\textwidth]{draft2.png}
%   \caption{Phase portrait of the differential equation system from Example 1.}\label{fig2:right}
% \end{minipage}
% \end{figure*}

% \clearpage
%%%%%%%%%%%%%%%%%%%%%%%%%%%%%%%%%%%%%%%%%%%%%%%%%%%%%%%%%%%%%%%%%%%%%%%%%%%%%
\section{Application I: Quadratic Unconstrained Binary Optimization via Discontinuous HNNs}
\label{sec:qubo-hnn}
\todo[inline]{New below}

We show that under a \emph{strong field} regime, 
a discontinuous HNN with sign activations globally recovers the unique QUBO optimizer.

\subsection{QUBO with Strong Field and Uniqueness}
Consider
\begin{equation}\label{eq:qubo-energy}
  \min_{s\in\{-1,1\}^n}\ E(s)=\tfrac12 s^\top Q s - h^\top s,\quad Q=Q^\top.
\end{equation}
\begin{definition}[Strong field]\label{def:strong-field}
\((Q,h)\) satisfies the strong field condition if
\(
|h_i|>\sum_{j\neq i}|Q_{ij}|\ \ \forall i.
\)
\end{definition}
\begin{lemma}[Unique minimizer]\label{lem:qubo-unique}
Under the strong field condition, \(s^\star=\mathrm{sign}(h)\) 
is the unique global minimizer of \eqref{eq:qubo-energy}. 
Proof in Appendix~\ref{app:qubo-proofs}.
\end{lemma}

\subsection{HNN Embedding and Global Convergence}
Use the Lur’e form \(\dot x=-Bx+Tg(x)+I\) with
\begin{equation}\label{eq:hnn-qubo}
  \dot x = -B x - \beta Q_{\mathrm{shift}}\, g(x) + \beta h,\quad
  g_i(u)=\mathrm{sign}(u),
\end{equation}
where \(B=\mathrm{diag}(b_i)\succ0\), \(\beta>0\), and
\begin{equation}\label{eq:Qshift}
  Q_{\mathrm{shift}}:=Q+\mu I,\quad \mu>\lambda^+_{\max}(Q)\ \Rightarrow\ Q_{\mathrm{shift}}\succ0.
\end{equation}
Then \(T=-\beta Q_{\mathrm{shift}}\) yields \(T+T^\top\prec0\), so \textbf{C3} holds (with \(\alpha=I\)); \textbf{C1-C2} hold for the sign activation.
\begin{theorem}[Global convergence to the planted optimizer]\label{th:hnn-solves-strong-field-qubo}
Under the strong field condition and \textbf{C1-C3}, \eqref{eq:hnn-qubo} has a unique Filippov equilibrium \(x^\star\), globally asymptotically stable, and \(g(x^\star)=\mathrm{sign}(x^\star)=s^\star\), the unique QUBO optimizer. 
Proof in Appendix~\ref{app:qubo-proofs}.
\end{theorem}

\subsection{Planted MAX-CUT with Strong Field}\label{subsec:planted-maxcut}
Let \(G=(V,E,w)\) be a weighted graph and \(Q^{\mathrm{MC}}\) the MAX-CUT matrix
\(
Q^{\mathrm{MC}}_{ij}=-\tfrac12 w_{ij}\ (i\ne j),\ Q^{\mathrm{MC}}_{ii}=0.
\)
Partition \(V=A\cup B\), and define a field \(h_i=\gamma\) for \(i\in A\), \(h_i=-\gamma\) for \(i\in B\).
If the weights satisfy
\begin{equation}\label{eq:gamma-condition}
  \gamma>\tfrac12\max\!\big\{(n_A-1)w_{\mathrm{in}}\!+\!n_B w_{\mathrm{out}},\ (n_B-1)w_{\mathrm{in}}\!+\!n_A w_{\mathrm{out}}\big\},
\end{equation}
then \((Q^{\mathrm{MC}},h)\) meets Definition~\ref{def:strong-field}. With \(Q_{\mathrm{shift}}=Q^{\mathrm{MC}}+\mu I\succ0\) and the HNN \eqref{eq:hnn-qubo}, Theorem~\ref{th:hnn-solves-strong-field-qubo} implies:
\begin{corollary}\label{cor:hnn-planted-maxcut}
The HNN output converges globally to \(s^\star\) aligned with the planted partition \((A,B)\), which is the unique minimizer of the MAX-CUT QUBO with field. Proof in Appendix~\ref{app:qubo-proofs}.
\end{corollary}
%%%%%%%%%%%%%%%%%%%%%%%%%%%%%%%%%%%%%%%%%%%%%%%%%%%%%%%%%%%%%%%%%%%%%%%%%%%%%

%%%%%%%%%%%%%%%%%%%%%%%%%%%%%%%%%%%%%%%%%%%%%%%%%%%%%%%%%%%%%%%%%%%%%%%%%%%%%
\section{Application II: Constrained Optimization via HNNs and the Penalty Method}
\label{sec:penalty-HNN}

We consider a nonlinear program
\begin{equation}\label{eq:orig-P}
  \min_{x\in\mathbb{R}^n} f(x)\quad\text{s.t.}\quad g_i(x)\le 0,\ i=1,\dots,m,
\end{equation}
and the quadratic exterior penalty \(p(x)=\sum_i (\max\{0,g_i(x)\})^2=\|g^+(x)\|_2^2\).
For \(c>0\), define the penalized subproblem
\begin{equation}\label{eq:penalty-Pc}
  P(c):\quad \min_{x} q_c(x):=f(x)+c\,p(x).
\end{equation}
Classical exact-penalty theory ensures that for \(c\) large enough, global minimizers of \(q_c\) approach feasible solutions of \eqref{eq:orig-P}.

\subsection{Embedding the Penalty Subproblem in Discontinuous~HNN}
Represent \(q_c\) with a convex–nonsmooth potential \(\phi\) whose subdifferential aggregates the smooth gradient and the penalty kink (via \(g_i^+\)); add a diagonal leakage to form
\begin{equation}\label{eq:penalty-V}
  V(x)=\phi(x)+\tfrac12 x^\top A x - \ell^\top x,\qquad A\succ0,
\end{equation}
and run \(\dot x\in -\partial V(x)\). This realizes \(\dot x=-Bx+T g(x)+I\) 
with finitely many discontinuities in the activation (from \(\max\{0,\cdot\}\)), 
satisfying our conditions after a standard diagonal shift of the quadratic part (YLG/Lur’e spirit). 
Then the unique GAS equilibrium \(x^\star\) solves \(P(c)\). 
Construction and checks are in Appendix~\ref{app:penalty-proofs}.

\subsection{Example: Soft–Margin SVM in Hopfield Form}\label{sec:SVM-HNN}
The \(\ell_2\) soft-margin SVM
\begin{equation}\label{eq:SVM-primal}
  \min_{w,b,\xi}\ \tfrac12\|w\|^2 + C\sum_i \xi_i\quad \text{s.t.}\quad 
  y_i(w^\top z_i + b) \ge 1-\xi_i,\ \xi_i\ge 0
\end{equation}
can be written as an exact-penalty objective
\begin{equation}\label{eq:SVM-penalty}
  q_\rho(u,\xi)=\tfrac12 u^\top u + C\sum_i \xi_i + \rho\sum_i \max\{0,\,1-\xi_i-y_i u^\top\bar z_i\},
\end{equation}
with \(u=(w^\top,b)^\top\), \(\bar z_i=(z_i^\top,1)^\top\), and sufficiently large~\(\rho\).
Let \(x=(u^\top,\xi^\top)^\top\). The piecewise-linear kinks in the hinge term define a componentwise discontinuous activation (single upward jump). Choosing a diagonal leakage and a shifted quadratic synaptic matrix (so \(T+T^\top\prec0\)) yields a discontinuous HNN of the Lur’e form satisfying \textbf{C1-C3}. Therefore:
\begin{proposition}[Unique globally convergent SVM solver]\label{prop:SVM}
The Hopfield network built from \eqref{eq:SVM-penalty} by \(\dot x\in -\partial V(x)\) has a unique Filippov equilibrium \(x^\star=(u^\star,\xi^\star)\) and \(x(t)\to x^\star\) from any initialization. 
The classifier parameters \((w^\star,b^\star)\) coincide with the optimal soft-margin SVM solution. 
\end{proposition}
Proof in Appendix~\ref{app:penalty-proofs}.

\paragraph*{Summary}
For constrained problems, exact/exterior penalties produce discontinuous activations (via \(\max\{0,\cdot\}\)) that squarely fit the Filippov/YLG toolbox. Under \textbf{C1-C3}, the associated HNN has a unique GAS equilibrium equal to the penalized optimizer, and—in the SVM case—recovers the primal solution without schedules or coordination mechanisms.
%%%%%%%%%%%%%%%%%%%%%%%%%%%%%%%%%%%%%%%%%%%%%%%%%%%%%%%%%%%%%%%%%%%%%%%%%%%%%


%%%%%%%%%%%%%%%%%%%%%%%%%%%%%%%%%%%%%%%%%%%%%%%%%%%%%%%%%%%%%%%%%%%%%%%%%%%%%
\section{Application III: Discontinuous HNN as a Real-Time Solver for Constrained Optimal Control Problem}
\label{sec:application_control}

We illustrate why a \emph{unique globally asymptotically stable} equilibrium is essential when the HNN serves as a \emph{real-time} optimizer inside a feedback loop.

\subsection{Setup: Linear Plant and Box-Constrained Quadratic Law}
Consider the LTI plant
\begin{equation}\label{eq:plant}
  \dot z = Fz + Gu, \qquad z\in\mathbb{R}^n,\ u\in\mathbb{R}^m,
\end{equation}
with componentwise box constraints \(u_{\min}\le u\le u_{\max}\).
At each frozen state \(z\), the control \(u^\ast(z)\) is the unique solution of the strictly convex QP
\begin{equation}\label{eq:opt-problem}
  \min_{u\in\mathbb{R}^m}\ J_z(u) := \tfrac12 u^\top H u + f(z)^\top u
  \quad\text{s.t.}\quad u_{\min}\le u\le u_{\max},
\end{equation}
where \(H\succ0\) and \(f(\cdot)\) is affine. The solver must operate continuously as \(z\) evolves; multi-valued or initial-state-dependent outcomes are unacceptable.

\subsection{HNN Structure via a Convex Nonsmooth Potential}
Let \(x\in\mathbb{R}^m\) denote the HNN state (interpreted as \(u\)).
Encode the cost and constraints with
\begin{multline}\label{eq:phi-def}
  \phi(x):=\tfrac12 x^\top Hx + \sum_{i=1}^m \iota_{[u_{\min,i},u_{\max,i}]}(x_i),
  \\
  \partial\phi(x)=Hx + N_{[u_{\min},u_{\max}]}(x),
\end{multline}
and define the parameterized energy
\begin{equation}\label{eq:Vz-def}
  V_z(x):=\phi(x)-b(z)^\top x+\tfrac12 x^\top A x,
\end{equation}
with \(A\succ0\) diagonal and \(b(z)\) affine in \(z\).
We run the \emph{subgradient flow}
\begin{equation}\label{eq:hnn-subgrad}
  \dot x \in -\partial V_z(x).
\end{equation}
This is a Filippov inclusion with finitely many first-order discontinuities (due to the box normal cone). It fits the Lur’e form \(\dot x=-Bx+Tg(x)+I\) and satisfies linear growth, one-sided slope and YLG-type frequency condition under the choices below.

\begin{proposition}[Real-time solver]\label{prop:rt-solver}
Under Assumptions~\textbf{C1--C3} and the model \eqref{eq:hnn-subgrad}–\eqref{eq:Vz-def}, for each frozen \(z\) the HNN admits a unique Filippov equilibrium \(x^\ast(z)\), globally asymptotically stable, and \(x^\ast(z)=u^\ast(z)\) solves \eqref{eq:opt-problem}. \emph{(Proof in Appendix~\ref{app:control-proofs}.)}
\end{proposition}

\paragraph*{Why uniqueness matters}
Embedding the HNN into the loop
\begin{equation}\label{eq:closed-loop}
  \dot z = Fz + Gx, \qquad \dot x\in -\partial V_z(x),
\end{equation}
requires a single-valued, robust control law \(u=\kappa(z)\). If multiple stable equilibria coexisted for a given \(z\), the closed-loop action would depend on internal initial conditions \(x(0)\), jeopardizing predictability and stability.

\subsection{1D Illustrative Example}\label{subsec:illustrative-example-control}
Plant: \(\dot z=-z+u\), \(|u|\le 1\). Instantaneous QP
\begin{equation}\label{eq:ex-qp}
  \min_{u\in[-1,1]}\ J_z(u):=\tfrac12 h u^2 + \alpha z u, \quad h>0,\ \alpha>0,
\end{equation}
has solution \(u^\ast(z)=\Pi_{[-1,1]}(-\alpha z/h)\).
Choose
\[
  \phi(x)=\tfrac12 h x^2 + \iota_{[-1,1]}(x),\quad
  V_z(x)=\phi(x)-\alpha zx+\tfrac12 a x^2,\ a>0,
\]
and evolve \(\dot x\in -\partial V_z(x) = -(h+a)x+\alpha z - N_{[-1,1]}(x)\).
Then
\[
  \dot V_z(x)\le -a x^2<0\ \ (x\neq x^\ast(z)),
\]
so \textbf{C1-C3} hold; \(x^\ast(z)=u^\ast(z)\) 
and the closed loop has the unique equilibrium \((z^\ast,x^\ast)=(0,0)\). 
See corresponding derivations in Appendix~\ref{app:control-proofs}.

\subsection{2D Illustrative Example}\label{subsec:2d-example-control}
Plant: \(\dot z = Az+Bu\) with
\[
A=\begin{bmatrix}-1&1\\0&-2\end{bmatrix},\quad
B=I_2,\quad u\in[-1,1]^2.
\]
QP:
\begin{multline*}    
  \min_{u\in[-1,1]^2}\ J_z(u)=\tfrac12 u^\top H u + q(z)^\top u, \\
  H=\mathrm{diag}(2,1),\ 
  q(z)=\begin{bmatrix}z_1+0.5 z_2\\ 2z_2\end{bmatrix}.
\end{multline*}
Then \(u^\ast(z)=\Pi_{[-1,1]^2}\big(-H^{-1}q(z)\big)\).
Let
\begin{multline*}
  \phi(x)=\tfrac12 x^\top Hx+\iota_{[-1,1]^2}(x),\\
  V_z(x)=\phi(x)-q(z)^\top x+\tfrac12 x^\top A_\ell x,\\ 
  A_\ell=\mathrm{diag}(a_1,a_2)\succ0,
\end{multline*}
and \(\dot x\in -\partial V_z(x)=-(H+A_\ell)x+q(z)-N_{[-1,1]^2}(x)\).
Then \(\dot V_z\le -x^\top A_\ell x\), hence \textbf{C1--C3} hold and \(x^\ast(z)=u^\ast(z)\) globally. 
See details in Appendix~\ref{app:control-proofs}.
% \vspace{-1mm}
\begin{remark}
The convex normal-cone implements saturation as a discontinuous activation, placing the design squarely in the Filippov/YLG setting used by our main theorem.
\end{remark}
%%%%%%%%%%%%%%%%%%%%%%%%%%%%%%%%%%%%%%%%%%%%%%%%%%%%%%%%%%%%%%%%%%%%%%%%%%%%%
\todo[inline]{New above}

% \section{Applications and Embeddings}\label{sec:apps}
% \todo[inline]{Map optimization tasks into the HNN template and verify C1--C3 per case.}


% \section{Numerical Illustrations}\label{sec:num}
% Describe the experimental setup (solver, step size/tolerance, initialization, random seeds). Mention Filippov integration strategy (hysteresis/smoothing or event handling).

% \begin{figure}[!t]
% \centering
% \includegraphics[width=\linewidth]{fig1.png}
% \caption{Illustrative phase portrait / convergence trace for a toy instance consistent with C1--C3.}
% \label{fig:demo}
% \end{figure}

% \todo[inline]{Reformat two column table into one column}

\section{Discussion and Conclusion}\label{sec:concl}
The classical Yakubovich--Leonov--Gelig machinery provides a concise proof
of the uniqueness and global stability of Hopfield equilibria under the
non‑smooth conditions \textbf{C1--C3}.  Besides simplifying the argument, the method
highlights how frequency–domain negativity of \(\alpha T+T^{\!*}\alpha\)
serves as the decisive criterion, irrespective of boundedness or monotonicity
of activation functions. For the Hopfield model new results on global stability
in the case of neither monotonic, nor bounded activation functions,
which extend the previous results from~\cite[see Theorems~1-3]{p1} 
and \cite[see Theorems~1-3]{our_paper}, 
were derived and proven.
Numerical simulations approving the obtained results are demonstrated.

\section*{Acknowledgment}

\noindent\textit{Dedication.} 
This paper is dedicated to the memory of V.A.~Yakubovich (21~Oct~1926--17~Aug~2012) 
and G.A.~Leonov (2~Feb~1947--23~Apr~2018). 
The centenary year of Yakubovich (2026) and the approaching 80th anniversary of Leonov (2027) 
commemorate their outstanding contributions to nonlinear control theory---especially to absolute stability and to systems with discontinuous nonlinearities---and to the development of cybernetics and early directions of artificial intelligence. 
G.A.~Leonov was a fomer student of V.A.~Yakubovich; both were our teachers and fathers in science.

\medskip
This work is supported by the Russian Science Foundation (RSF, Russia) 
project 25-11-00147.

% \printbibliography 


% \section{Detailed Proofs}\label{app:proofs}
% Provide complete derivations for Theorems~\ref{thm:uniq}--\ref{thm:gas}, including any auxiliary lemmas and technical facts about Filippov solutions.

% \section{Algorithmic Tests and Implementation Notes}
% Give explicit LMI/QMI or frequency-domain checks for C3 and pseudocode for verification.

% \section{Additional Lemmas}
% State and prove any auxiliary results not in the main text.

% \section{Biological and Circuit Intuition (Informative)}\label{app:biocircuit}
% This section is provided for context and does not participate in the proofs.
% Following the classical Kirchhoff circuit interpretation, neuron states are voltages across
% amplifier-like units with refractory dynamics modeled by a diagonal matrix $B=\mathrm{diag}(b_i)>0$,
% interconnections by a conductance matrix $T$, and external inputs by $I$.
% Under standard circuit laws, one obtains the HNN in Lur\'e form
% \begin{equation}
% \dot{x} = -Bx + T\,g(x) + I,
% \end{equation}
% where $g(x)=(g_1(x_1),\ldots,g_n(x_n))^\top$ is the activation map (possibly discontinuous).
% This viewpoint motivates the diagonal Lyapunov test used in C3.

% \bibliographystyle{IEEEtran}
% \bibliography{ref}

% \acksection 


%----------------------------------------------------------------------------------------
%	 REFERENCES
%----------------------------------------------------------------------------------------

\printbibliography 

\newpage
\appendices

%%%%%%%%%%%%%%%%%%%%%%%%%%%%%%%%%%%%%%%%%%%%%%%%%%%%%%%%%%%%%%%%%%%%%%%%%%%%%%%%%%%%%%%%%%%%%%%%
%%%%%%%%%%%%%%%%%%%%%%%%%%%%%%%%%%%%%%%%%%%%%%%%%%%%%%%%%%%%%%%%%%%%%%%%%%%%%%%%%%%%%%%%%%%%%%%%
\section{Yakubovich--Leonov--Gelig Functionals and Auxiliary Results for Lur\'e Systems.} \label{app:YLG}

% \todo[inline]{Add definitions from YLG theory}

% \noindent\smallskip
% \textbf{A. Existence and global continuation of Filippov solutions}
\subsection{Existence and global continuation of Filippov solutions}

\begin{theorem}[Local existence for differential inclusions (see~{\cite[Thm.~1.1]{YLG}})]
Let $t_0\in\mathbb{R}$, $a\in\mathbb{R}^N$, and let
\[
D := \bigl\{(x_1,t_1)\in\mathbb{R}^N\times\mathbb{R}
       \;\big|\; |t_1-t_0|\le \alpha,\ |x_1-a|\le\rho \bigr\}
\]
for some $\alpha>0$, $\rho>0$.
Suppose that for every $(x_1,t_1)\in D$ the multivalued map
$f:\mathbb{R}^N\times\mathbb{R}\rightrightarrows\mathbb{R}^N$
satisfies:
\begin{enumerate}
  \item $f$ is upper semicontinuous at $(x_1,t_1)$;
  \item the value $f(x_1,t_1)$ is nonempty, closed, bounded and convex;
  \item there exists $c>0$ such that
        \begin{equation}\label{eq:f-bound-local}
           \sup\{\,|y| : y\in f(x_1,t_1)\,\} \;=\; c,
           \qquad (x_1,t_1)\in D .
        \end{equation}
\end{enumerate}
Consider the differential inclusion
\begin{equation}\label{eq:incl-local}
   \dot x(t) \in f(x(t),t).
\end{equation}
Then, for any initial condition
\begin{equation}\label{eq:initial-local}
   x(t_0) = a,
\end{equation}
there exists at least one absolutely continuous solution $x(\cdot)$
of~\eqref{eq:incl-local}--\eqref{eq:initial-local}
on the interval
\[
   |t-t_0|\le\tau, \qquad
   \tau := \min\{\alpha,\ \rho/c\}.
\]
\end{theorem}

\begin{theorem}[Wintner--type continuation criterion (see~{\cite[Thm.~1.3]{YLG}})]
Let the multivalued right-hand side $f:\mathbb{R}^N\times(0,+\infty)
\rightrightarrows\mathbb{R}^N$ of the inclusion
\begin{equation}\label{eq:incl-global}
   \dot x(t) \in f(x(t),t)
\end{equation}
be defined for all $x\in\mathbb{R}^N$ and all $t\in(0,+\infty)$
(\emph{respectively,} for all $t\in\mathbb{R}$).
Assume that there exists a function $L:[0,+\infty)\to(0,+\infty)$
such that
\begin{equation}\label{eq:L-bound}
   |\xi| \le L(|x|),
   \qquad \forall\,\xi\in f(x,t),\ \forall\,x\in\mathbb{R}^N,\
          t\in(0,+\infty),
\end{equation}
and $L$ is positive, continuous and satisfies
\begin{equation}\label{eq:Wintner-integral}
   \int\limits_0^{+\infty} \frac{dr}{L(r)} = +\infty.
\end{equation}
Then every solution of~\eqref{eq:incl-global} can be continued to the
entire half-line $0<t<+\infty$
(\emph{respectively,} to the whole line $-\infty<t<+\infty$).
\end{theorem}

% \bigskip\noindent
% \textbf{B. Representing Filippov solutions via measurable extended nonlinearities}

\subsection{Representing Filippov solutions via measurable extended nonlinearities}

\begin{lemma}[Makarov measurable selection (see~{\cite[Lem.~1.2]{YLG}})]\label{lemma:makarov}
Let $\mathcal{H}$ and $Q$ be compact metric spaces and
$\Psi:\mathcal{H}\to Q$ be a continuous mapping onto $Q$.
Then there exists a Borel-measurable mapping
$\theta:Q\to\mathcal{H}$ such that
\[
   \Psi(\theta(q)) = q, \qquad \forall\,q\in Q.
\]
\end{lemma}

\begin{theorem}[Makarov measurable extended nonlinearity (see~{\cite[Thm.~1.7]{YLG}})]\label{thm:makarov}
Let $F:[a,b]\times\mathbb{R}^N\times\mathbb{R}^m\to\mathbb{R}^N$ be a
vector function continuous in all arguments.
Let $\sigma:[a,b]\times\mathbb{R}^N\to\mathbb{R}^\ell$ be a continuous
vector function and let $\varphi:[a,b]\times\mathbb{R}^\ell
\rightrightarrows\mathbb{R}^m$ be a multivalued map which is
semicontinuous and whose values are closed bounded subsets of
$\mathbb{R}^m$.
Let $x_0:[a,b]\to\mathbb{R}^N$ be an absolutely continuous function
satisfying, for almost all $t\in[a,b]$,
\begin{equation}\label{eq:makarov-incl}
   \dot x_0(t) \in \{\,F(t,x_0(t),\xi)\;|\; \xi\in A(t)\,\},
\end{equation}
where
\[
   A(t) := \varphi\bigl(t,\sigma(t,x_0(t))\bigr).
\]
Then there exists a Lebesgue-measurable vector-valued function
$\xi_0:[a,b]\to\mathbb{R}^m$ such that
\begin{equation}\label{eq:makarov-realization}
   \dot x_0(t) = F\bigl(t,x_0(t),\xi_0(t)\bigr), \qquad
   \xi_0(t)\in A(t),
\end{equation}
for almost all $t\in[a,b]$.
\end{theorem}

\begin{remark}%[On measurable selections and extended nonlinearities]
Lemma~\ref{lemma:makarov} and Theorem~\ref{thm:makarov} are concrete instances of general measurable
selection results of Kuratowski and Ryll--Nardzewski
(see~\cite{kuratowski-ryll-nardzewski-1965} 
and also~\cite{wagner-1977-measurable-selection-survey,castaing-valadier-1977,srivastava-1998-borel-sets}).

Lemma~\ref{lemma:makarov} deals with a continuous surjection
$\Psi:\mathcal{H}\to Q$ between compact metric spaces and asserts the
existence of a Borel-measurable right inverse
$\theta:Q\to\mathcal{H}$ with $\Psi(\theta(q))=q$ for all $q\in Q$.
If one defines a multifunction $F(q):=\Psi^{-1}(\{q\})$, then
$F(q)$ is a nonempty compact subset of $\mathcal{H}$ and the map
$q\mapsto F(q)$ is weakly measurable. The Kuratowski--Ryll-Nardzewski
measurable selection theorem ensures the existence of a measurable
selection $\theta(q)\in F(q)$; 
Lemma~\ref{lemma:makarov} is precisely this special
case, specialized to compact metric spaces and a single continuous
surjection.

Theorem~\ref{thm:makarov} provides a ``pathwise'' measurable selection result
for differential inclusions. Given an absolutely continuous solution
$x_0(\cdot)$ of an inclusion of the form
\[
   \dot x_0(t) \in \{\,F(t,x_0(t),\xi)\mid \xi\in A(t)\,\},
\]
with $F$ continuous and $A(t)$ obtained from a semicontinuous
multifunction with closed bounded values, 
Theorem~\ref{thm:makarov} states that there
exists a Lebesgue-measurable function $\xi_0(t)\in A(t)$ such that
\[
   \dot x_0(t) = F\bigl(t,x_0(t),\xi_0(t)\bigr)
   \quad\text{for a.e. }t.
\]
In other words, along any absolutely continuous solution of the
set-valued system one can select a measurable, single-valued
``extended nonlinearity'' which realizes the dynamics exactly. This is
again an application of measurable selection theorems (now to the
multifunction
$t\mapsto\{\xi\in A(t)\mid F(t,x_0(t),\xi)=\dot x_0(t)\}$).

In the Hopfield-network setting considered in this paper, the
right-hand side of the discontinuous HNN is interpreted in the Filippov
sense. For any Filippov solution $x(\cdot)$ of the HNN, Theorem~\ref{thm:makarov}
yields a measurable ``extended activation''
$\eta(t)\in\operatorname{conv} g(x(t))$ such that, almost everywhere,
the trajectory satisfies an ordinary differential equation of Lur'e
type,
\[
   \dot x(t) = -B x(t) + T \eta(t) + I.
\]
Thus every Filippov trajectory of the discontinuous HNN can be viewed
as a trajectory of a single-valued Lur'e system with a measurable
input $\eta(\cdot)$ constrained to the convexified activation graph.
This realization is crucial for applying the frequency-domain
Yakubovich--Leonov--Gelig functionals and the Barabanov-type
dichotomy/quasi-gradient-like criteria in the stability analysis.
\end{remark}

% \bigskip\noindent
% \textbf{D. Lyapunov-type lemmas for dichotomy and stability of stationary sets}
\subsection{Lyapunov-type lemmas for dichotomy and stability of stationary sets}

Consider now the autonomous differential inclusion
\begin{equation}\label{eq:aut-incl}
   \dot x(t) \in f(x(t)), \qquad x(t)\in\mathbb{R}^n,
\end{equation}
with stationary set
\[
   \Lambda := \{\,x\in\mathbb{R}^n : 0\in f(x)\,\}.
\]

\begin{lemma}[Lyapunov function and dichotomy {\cite[Lem.~1.4]{YLG}}]\label{lem:YLGl1-4}
Suppose there exists a continuous function $V:\mathbb{R}^n\to\mathbb{R}$
having the properties:
\begin{enumerate}
  \item $V[x(t)]$ is nonincreasing in $t$ for any solution $x(t)$
        of~\eqref{eq:aut-incl};
  \item if the identity $V[x(t)]\equiv\mathrm{const}$ holds for all
        $-\infty<t<\infty$ along some solution $x(t)$ that is bounded
        for $-\infty<t<\infty$, then this solution $x(t)$ is stationary.
\end{enumerate}
Then system~\eqref{eq:aut-incl} is dichotomic.
\end{lemma}

\begin{lemma}[Quasi-gradient-likeness {\cite[Lem.~1.5]{YLG}}]\label{lem:YLGl1-5}
Suppose there exists a continuous function $V:\mathbb{R}^n\to\mathbb{R}$
which satisfies the conditions \emph{(1)} and \emph{(2)} of
Lemma~\ref{lem:YLGl1-4} and, in addition,
\begin{enumerate}
  \setcounter{enumi}{2}
  \item $V(x)\to +\infty$ as $|x|\to +\infty$.
\end{enumerate}
Then system~\eqref{eq:aut-incl} is quasi-gradient-like.
\end{lemma}

\begin{lemma}[Lyapunov stability of the stationary set {\cite[Lem.~1.6]{YLG}}]\label{lem:YLGl1-6}
Suppose that the stationary set $\Lambda$ of~\eqref{eq:aut-incl} is
bounded and that there exists an $\varepsilon$-neighbourhood
$\Lambda_\varepsilon$ of $\Lambda$ such that there is a continuous
function $V:\Lambda_\varepsilon\times\Lambda\to\mathbb{R}$ with the
properties:
\begin{enumerate}
  \item $V(x,c)>0$ for all $x\in\Lambda_\varepsilon\setminus\Lambda$ and
        all $c\in\Lambda$;
  \item $V(c,c)=0$ for all $c\in\Lambda$;
  \item for any solution $x(t)$ of~\eqref{eq:aut-incl},
        $V[x(t),c]$ is nonincreasing in $t$ whenever $x(t)\in\Lambda_\varepsilon$.
\end{enumerate}
Then the stationary set $\Lambda$ is Lyapunov stable.
\end{lemma}

\begin{lemma}[Pointwise global stability {\cite[Lem.~1.7]{YLG}}]\label{lem:YLGl1-7}
Assume that system~\eqref{eq:aut-incl} is quasi-gradient-like, all the
hypotheses of Lemma~\ref{lem:YLGl1-6} are satisfied, and the equality
$V(x,c)=0$ can hold only for $x=c$.
Then the stationary set $\Lambda$ of system~\eqref{eq:aut-incl} is
pointwise globally stable.
\end{lemma}

\subsection{YLG frequency–domain theorem for several nonlinearities}

Consider the Lur'e–type system
\begin{equation}
  \dot x = P x + q\,\varphi(\sigma), 
  \qquad \sigma = r^{\top} x,
  \label{eq:YLG-31-system}
\end{equation}
where $P\in\mathbb{R}^{n\times n}$, $q\in\mathbb{R}^{n\times m}$, 
$r\in\mathbb{R}^{n\times m}$ are constant real matrices, 
$\sigma=(\sigma_1,\dots,\sigma_m)^{\top}$ and
$\varphi(\sigma)=(\varphi_1(\sigma_1),\dots,\varphi_m(\sigma_m))^{\top}$
is an $m$–vector of scalar nonlinearities.  
Each $\varphi_i(\sigma_i)$ is assumed to be \emph{piecewise single–valued}
in the sense of~\cite{YLG}: it is single–valued and continuous everywhere,
except at some isolated points (which do not accumulate on a finite
interval), and at each such point its value is a closed segment containing
the right and left limits.

Define the transfer matrix
\[
  X(\lambda) := r^{\top}(P-\lambda I)^{-1}q .
\]

\begin{theorem}[Barabanov {\cite[Thm.~3.1]{YLG}}]\label{th:YLG-3-1}
Assume that the following conditions hold.
\begin{enumerate}%[(i)]
\item The transfer function $X(\lambda)$ is \emph{irreducible} in the
sense of Definition~2.6 in~\cite{YLG}.

\item There exist diagonal matrices
\begin{multline*}
   T=\operatorname{diag}(\tau_1,\dots,\tau_m)>0,\\
   M=\operatorname{diag}(\mu_1,\dots,\mu_m)\ge0,\\
   \theta=\operatorname{diag}(\vartheta_1,\dots,\vartheta_m)      
\end{multline*}
such that
\begin{equation}
   T M + \Re\bigl[(T+i\omega\theta)\,X(i\omega)\bigr] \;\succeq\; 0,
   \quad \forall\,\omega\in(-\infty,+\infty),
   \label{eq:YLG-31-freq}
\end{equation}
and
\begin{equation}
   \det(T+\lambda\theta)\neq0
   \quad\text{whenever}\quad
   \det(P-\lambda I)=0 .
   \label{eq:YLG-31-det}
\end{equation}

\item For each $i=1,\dots,m$ the nonlinearity $\varphi_i$ satisfies
\begin{align}
 (\sigma_i-\mu_i\xi_i)\,\xi_i &\ge 0,
   &&\forall\,\xi_i\in\varphi_i(\sigma_i),
   \label{eq:YLG-31-sector}\\
 \sigma_i &\notin \mu_i\varphi_i(\sigma_i),
   &&\text{for all }\sigma_i\neq0 ,
   \label{eq:YLG-31-sigma-notin}
\end{align}
and, in addition, if $\mu_i\neq0$ and $\vartheta_i<0$ then
\begin{equation}
   \int\limits_0^{+\infty} \bigl[\lambda-\mu_i\varphi_i(\lambda)\bigr]\,d\lambda
   = +\infty .
   \label{eq:YLG-31-integral}
\end{equation}

\item There exist numbers $\mu_i^0$, $1\le i\le m$, such that
$\mu_i^0>\mu_i$ and the matrix
\[
   P + q M_0^{-1} r^{\top}, 
   \qquad M_0 := \operatorname{diag}(\mu_1^0,\dots,\mu_m^0),
\]
is Hurwitz.

\item The stationary set $\Lambda$ of system~\eqref{eq:YLG-31-system}
is bounded. Moreover, if a solution $(x(t),\xi(t))$ of
\eqref{eq:YLG-31-system} satisfies
\begin{equation}
   \xi_i(t)\,\sigma_i(t)\equiv 0\quad(1\le i\le m), 
   \qquad x^{\top}(t)H x(t)\equiv\mathrm{const},
   \label{eq:YLG-31-orthogonality}
\end{equation}
for all $-\infty<t<+\infty$ and some positive definite matrix
$H\succ0$, then $x(t)\equiv x^0\in\mathbb{R}^n$, 
$\xi(t)\equiv\xi^0\in\mathbb{R}^m$ are constant, and
\[
   P x^0 + q\xi^0 = 0,
\]
i.e., the solution $(x(t),\xi(t))$ is stationary.

\end{enumerate}
Then system~\eqref{eq:YLG-31-system} is \emph{quasi-gradient-like}.

Suppose, in addition, that
\begin{enumerate}%[(i)]\setcounter{enumi}{5}
\item there exists a number $\varepsilon>0$ such that, whenever
$\rho(x,\Lambda)<\varepsilon$, then for any stationary solution
$(x^0,\xi^0)$ the inequalities
\begin{multline}
  (\xi_i-\xi_i^0)\bigl[\sigma_i-\sigma_i^0-\mu_i(\xi_i-\xi_i^0)\bigr]\ge0,
  \\ \forall\,\xi_i\in\varphi_i(\sigma_i), \quad 1\le i\le m,
  \label{eq:YLG-31-local}
\end{multline}
hold, where $\sigma_i$ and $\sigma_i^0$ are the $i$th components of the
vectors $r^{\top}x$ and $r^{\top}x^0$, respectively. In the case
$\vartheta_i<0$, the expression in square brackets in
\eqref{eq:YLG-31-local} is nonzero whenever $\xi_i\neq\xi_i^0$ and
$x\neq x^0$.
\end{enumerate}
Then the stationary set $\Lambda$ of system~\eqref{eq:YLG-31-system}
is \emph{pointwise globally stable}.
\end{theorem}

\todo[inline]{NEW BELOW}
\subsection{Proof of Proposition 3}\label{prop3_proof}
\begin{proof}
Define
\[
x^{(1)}(t):=\bigl(0,\,a(1-e^{-t})\bigr),\qquad x^{(2)}(t):=\bigl(0,\,-a(1-e^{-t})\bigr).
\]
As shown above, both $x^{(1)}(\cdot)$ and $x^{(2)}(\cdot)$ are absolutely continuous, satisfy \eqref{eq:example-system} almost everywhere and $x^{(1)}(0)=x^{(2)}(0)=(0,0)$, hence are Filippov solutions with the same initial data.

For any $t>0$,
\[
x^{(1)}(t)\neq x^{(2)}(t),
\]
since $x_{2}^{(1)}(t)=a(1-e^{-t})>0$ and $x_{2}^{(2)}(t)=-a(1-e^{-t})<0$ when $a>0$. Thus the solutions are distinct, and forward uniqueness fails.
\end{proof}
\todo[inline]{NEW ABOVE}

% \bigskip\noindent
% \textbf{C. Frequency–domain criteria and YLG functionals for gradient Lur'e systems}
% \subsection{Frequency–domain criteria and YLG functionals for gradient Lur'e systems}

% Consider the Lur'e-type system
% \begin{align}
%     \dot{x} &= P x + q \,\xi, \qquad \xi = \nabla \zeta(\sigma),
%     \label{eq:grad-lure}\\
%     \sigma &= r^{\top} x,
%     \label{eq:sigma-def}
% \end{align}
% where $P\in\mathbb{R}^{n\times n}$, $q\in\mathbb{R}^{n\times m}$,
% $r\in\mathbb{R}^{n\times m}$ are constant real matrices ($m\le n$), and
% $\zeta:\mathbb{R}^{m}\to\mathbb{R}$ is a scalar potential function.
% The vector $\nabla\zeta(\sigma)$ consists of the partial derivatives of
% $\zeta$ with respect to the components of $\sigma$; these derivatives
% are assumed to be piecewise continuous. The system is understood in the
% sense of differential inclusions when $\nabla\zeta$ is discontinuous.

% The transfer matrix of the linear part is
% \begin{equation}
%     W(\lambda) := r^{\top} (P - \lambda I)^{-1} q,
% \end{equation}
% and on the imaginary axis we write
% \[
%    W(i\omega) = \Re W(i\omega) + i\,\Im W(i\omega),
% \]
% where
% \begin{align*}
%    \Re W(i\omega) &= \tfrac{1}{2}\bigl(W(i\omega) + W^{\ast}(i\omega)\bigr),\\
%    \Im W(i\omega) &= \tfrac{1}{2i}\bigl(W(i\omega) - W^{\ast}(i\omega)\bigr),
% \end{align*}
% with $\ast$ denoting conjugate transpose.

% \begin{theorem}[Barabanov, 1982; dichotomy criterion (see~{\cite[Thm.~3.13]{YLG}})]\label{th:Barabanov-dichotomy}
% Let the matrix pair $(P,q)$ be controllable and assume that there exists
% a real number $\omega_\ast$ such that
% \[
%    \Im W(i\omega_\ast) \prec 0.
% \]
% Then system~\eqref{eq:grad-lure}, \eqref{eq:sigma-def} is dichotomic in
% the class of nonlinearities $\zeta(\sigma)$ with piecewise continuous
% derivatives if and only if the following conditions hold:
% \begin{align}
%    \det(P - i\omega I) &\neq 0, \qquad \forall\,\omega>0,
%    \label{eq:Barabanov-3154}\\[2mm]
%    \Im W(i\omega) &\prec 0, \qquad \forall\,\omega>0.
%    \label{eq:Barabanov-3155}
% \end{align}
% \end{theorem}

% \begin{theorem}[Barabanov, 1982; quasi-gradient-like criterion (see~{\cite[Thm.~3.14]{YLG}})]\label{th:Barabanov-quasigrad}
% Let the matrix pair $(P,q)$ be controllable, let there exist a number
% $\omega_\ast$ such that $\Im W(i\omega_\ast)\prec 0$, and assume that
% the matrix $P + \varepsilon q r^{\top}$ is Hurwitz for all sufficiently
% small $\varepsilon>0$. Then the frequency conditions
% \eqref{eq:Barabanov-3154}, \eqref{eq:Barabanov-3155} are necessary and
% sufficient for system~\eqref{eq:grad-lure}, \eqref{eq:sigma-def} to be
% quasi-gradient-like in the class of nonlinearities $\zeta(\sigma)$
% satisfying
% \begin{equation}
%    \zeta(\sigma) \;\to\; +\infty \quad \text{as}\quad |\sigma|\to\infty.
%    \label{eq:Barabanov-3160}
% \end{equation}
% \end{theorem}

\section{Global stability results for HNNs using D.W.~Shevitz and B.~Paden's approach.} \label{app:proofs}

% \subsection*{Tools for the Stability Analysis}

To analyze the global asymptotic stability of the equilibrium point of system~\eqref{eq:special}, 
an \emph{alternative} approach is presented, where the 
generalized Lyapunov method \cite{p9} is employed, 
involving the study of the sign-definiteness of the derivative of a specially chosen Lyapunov function of the Lurie-Postnikov type.
This function also serves as the energy function for Hopfield's neural network~\cite{p10}.
Using the generalized Clarke gradient~\cite{p14} and the chain rule for regular functions~\cite{p9},
we can modify the LaSalle invariance principle \cite[Cor. 4.2, p. 129]{p223}
and apply it to a regular Lyapunov function that is differentiable almost everywhere and radially unbounded.

%%%%%%%%%%%%%%%%%%%%%%%%%%%%%%%%%%%%%%%%%%%%%%%%%%%%%%%%%%%%%%%%%%%%%%%%%%%%%%%%%%%%%%%%%%%%%%%%
\subsection{Existence and Continuation of Solutions}
Before investigating global asymptotic stability, it is necessary to ensure that the solutions of the system \eqref{eq:special} exist on the interval $[0, +\infty)$.

\begin{theorem}[Existence and Continuation of Solutions]\label{max_int}

Let the activation function of the neural network belong to the extended class $g(x)\in \mathcal{G'}$, and condition \textbf{C1} is satisfied. Then the interval of existence for each solution of the system \eqref{eq:special} with initial conditions $x(0) = x_0$ is $[0, +\infty)$.

\end{theorem}

\begin{remark} [The idea of proof]
Similarly to the way such result was proven in \cite{eq_ref_11, GAS_of_delayed, GAS_discrete_time_varying}, this theorem can be proven using Gronwall’s inequality and the existence and continuation theorems for differential inclusions by A.F. Filippov \cite{p2}. Therefore such proof mostly depends on the theory of upper semicontinuous set-valued mappings and the fact that the set-valued map from the right hand side of inclusion \eqref{eq:special_conv} satisfies all conditions of Filippov's theorems.
\end{remark}
\begin{IEEEproof}
To prove this theorem, we will use the fact that the existence and continuation of solutions \( x(t) \) of differential inclusion \( \dot{x} \in F(x) \) with the initial condition \( x(0) = x_0 \) on the interval \([0, t_0]\) for \( t_0 > 0 \) is guaranteed if the set-valued mapping \( F(x) \) is upper semicontinuous and has nonempty, convex, and compact values. In other words, it satisfies the so-called basic requirements of A.F. Filippov's theorems on the existence and continuation of solutions \cite[Theorems 1 and 2, pp. 77-78]{p2}. 

Since the activation function $g(x(t))$ in this case belongs to the extended (in comparison with such a class presented in \cite{p1}) class of functions
$\mathcal{G'}$, we have:
\[
\begin{aligned}
    \varphi(x(t)) &= \overline{\operatorname{conv}}\big(Bx(t) + T g(x(t)) + I\big) = \\
    &= Bx(t) + T ~\overline{\operatorname{conv}}\big(g(x(t))\big) + I,
\end{aligned}
\]
Then we can write:
\[
\dot{x}(t) \in \varphi\big(x(t)\big) = Bx(t) + T ~\overline{\operatorname{conv}}\big(g(x(t))\big) + I. \tag{+}\label{2.2.6}
\]
Where:
\[\overline{\operatorname{conv}}\big(g(x)\big) = \big(\overline{\operatorname{conv}}(g_1(x_1)), ..., \overline{\operatorname{conv}}(g_n(x_n))\big),
\]
and for functions belonging to the class $\mathcal{G'}$, the following holds:
\[
\overline{\operatorname{conv}}(g_i(x_i)) = [g_i(x_i^-), g_i(x_i)^+].
\]
From the form of \eqref{2.2.6}, it follows that this set-valued mapping has 
nonempty, convex, and compact values. Moreover, because this mapping was defined as part of the concept of Filippov solutions it also implies that  by Theorem 2.2.1 from \cite[p. 40]{p16}, it is upper semicontinuous. With these conditions satisfied, we can apply theory presented by Filippov, which implies that solutions to the system exist, and the existence interval for each  solution of system \eqref{eq:special} with initial data $x(0) = x_0$ is $[0, +\infty)$.
\end{IEEEproof}
\subsection*{Existence of an equilibrium point of the Network}
Now, based on assumptions \textbf{C2} and \textbf{C3$'$} (Diagonal Lyapunov stability), we formulate Lemma 1 which will be used both in Theorem 2 and 3 to prove existence and stability of the equilibrium point. It can be considered an attempt to rewrite these assumptions to align proofs presented in this work with those presented by M. Forti and P. Nistiri \cite{p1}. Since the proof of the Lemma consists only of trivial mathematical transformations and directly follows from \textbf{C2} and \textbf{C3$'$}, we do not provide it in this work.  \\
% (CHANGE d -> c to better align)
\begin{lemma}\label{main_lemma}
    \textnormal{Let \textbf{C2} and \textbf{C3$'$} hold, then there exist positive $d$ and $\varepsilon$ such that:}
    \begin{equation*}\label{2.3.2}
        C_i \leqslant \frac {1 -\varepsilon b_i}{d\alpha_i b_i}
    \end{equation*}
    \textnormal{also:}
    \begin{equation}\label{2.3.3}   
     \|-B^{-1} T\|^2_2 < {-\lambda_M b_m}d  \varepsilon.  
    \end{equation}
   
    
\end{lemma}

\begin{theorem}[Existence of an equilibrium point of the Network]\label{exists}

Let $g\in \mathcal{G'}$ and conditions \textbf{C1-C3$'$} are satisfied, then the neural network \eqref{eq:special} has an equilibrium point:
\begin{equation*}
	0\in\varphi(e)=Be+T ~ \overline{\operatorname{conv}}\big(g(e)\big)+I.
\end{equation*}

\end{theorem}

\begin{remark} [The idea of proof]
The proof of this theorem can be carried out using the fixed-point theorem 3.2.8 \cite[p. 90]{p3}. A similar result was proven in the article by M. Forti and P. Nistri \cite{p1} using Kakutani's fixed-point theorem 3.2.3 \cite[p. 87]{p3}. However, the use of this theorem will no longer give the desired result, as the boundedness of the activation function was required for its application.
\end{remark}

\begin{IEEEproof}
    Let us recall the set-valuep mapping from the right-hand side of the differential inclusion in the definition of a Cauchy problem solution for a system of differential equations with a discontinuous right-hand side in the sense of Filippov (formula \eqref{eq:special_conv}):
\[\begin{aligned}
        \varphi(x(t))&=\overline{\operatorname{conv}}\big(Bx(t) + T g(x(t)) + I\big)=\\
        &=Bx(t)+T ~\overline{\operatorname{conv}}\big(g(x(t))\big)+I.
    \end{aligned} \]
   
    By repeating the reasoning given in the proof of theorem \ref{max_int}, we conclude that this is an upper semicontinuous set-valued mapping with nonempty, convex, and compact images. 
    
    To apply theorem 3.2.8 \cite[p. 90]{p3}, 
    we introduce a dependence on $\lambda \in [0,1]$ into the system, thereby defining new mapping $\varPhi(\lambda,x(t))$:
    \[\varPhi(\lambda,x(t)) = Bx + \lambda T~ \overline{\operatorname{conv}}(g(x))+\lambda I.\]
    As the set $K$ (from the theorem 3.2.8), we choose a closed sphere in $\mathbb{R}^n$:
    \[K = \{x:~x^\mathrm{T}x \leqslant R^2 \} = \{ x:~\|x\| \leqslant R\}.\]
Now we determine the tangent cone\footnotemark[2] to $K$:
\[T_K(x) = \{y : y^\mathrm{T}x \leqslant 0\}, \quad \forall x \in \partial K.\]
Since $\varPhi(0,x) = Bx$, and $B =\operatorname{diag}(-b_1,...,-b_n)\in M_n(\mathbb{R})$, $b_i>0$, we obtain:
\[\varPhi(0,x)\cap T_k(x) = \{Bx\} \ne \emptyset, \quad \forall x \in \partial K.\]
The first condition of theorem 3.2.8 is satisfied; it remains to show that for all $x \in \partial K$ and $\lambda \in [0,1]$, the condition $0 \notin \varPhi(\lambda,x)$ holds.
\footnotetext[2]{Let $K$ be a convex subset of $\mathbb{R}^n$. The \textit{tangent cone} to $K$ at a point $x\in K$ is defined as the closed convex cone generated by the vector $K-x$: $T_K(x)=\overline{\bigcup\limits_{h>0}\frac{K-x}{h}}.$}
% Тут раньше была не теорема Маркова, а теорема 8.1.3 из \cite[стр. 308]{p3} нужно уточнить подходит ли теоерма Маркова и её источник 
Let $\gamma \in \overline{\operatorname{conv}}(g(0))$ be a constant vector. Since $\varPhi(\lambda,x)$ is upper semicontinuous and, therefore, a measurable mapping, then the Makarov’s theorem \cite[p. 38, Th. 1.7]{Leonov} implies that if $x(t)$ is a solution of system \eqref{eq:special} on the interval $[t_0,t_1]$, where $t_0<t_1\leqslant+\infty$, then there exists an almost everywhere measurable mapping $\eta(t) \in \overline{\operatorname{conv}}\big(g(x(t))\big)$ such that the equality holds almost everywhere on $[t_0, t_1]$:
\[\dot x(t) = Bx(t) + \lambda T \eta(t) + \lambda I \tag{a.e.}.\]
Consider $f \in \varPhi(\lambda, x)$ such that:
\[f = Bx(t) + \lambda T \eta(t) + \lambda I .\]
Since $\gamma \in \overline{\operatorname{conv}}(g(0))$, the following equality holds:
\[0 =   0 + \lambda T \gamma + \lambda  I .\]
Subtracting one from the other, we obtain: 
\begin{equation}\label{2.3.4}
    f = Bx(t) + \lambda T \widetilde\eta(t) + \lambda \widetilde I,
\end{equation}
where $\widetilde\eta(t) = \eta(t) - \gamma$ and $\widetilde I = T \gamma + I$.
From condition \textbf{C2} and subsequently lemma \ref{main_lemma} it follows that:
    \[
    1 - C_i d  \alpha_i b_i \geqslant \varepsilon b_i. 
    \]
Rewriting in vector form and applying the statement of the lemma 1:
\begin{equation}\label{2.3.5}
x^\mathrm{T} x - d \widetilde{\eta}^{\mathrm{T}}  \alpha  Bx \geqslant \varepsilon x^\mathrm{T} (-B) x . 
\end{equation}

Multiplying equation \eqref{2.3.4} by the same bracket:
\begin{equation}\label{2.3.6}
\begin{aligned}&
\underbrace{(-2B^{-1}x+2d \alpha\widetilde\eta)^\mathrm{T}}_{=\xi} f  =\\&= \underbrace{(-2B^{-1}x+2d \alpha\widetilde\eta)^\mathrm{T}}_{=\xi}(Bx + \lambda T \widetilde{\eta} + \lambda \widetilde I ).    
\end{aligned}
\end{equation}
Here, $(-2B^{-1}x+2d\alpha\widetilde{\eta}) = \xi$ can be interpreted as an element of the Clarke generalized gradient \cite{p14} for the neural networks energy function (this term will be further investigated in the proof of the theorem 3). Expanding the brackets, equation \eqref{2.3.6} becomes:
\begin{equation*}\label{2.3.7} 
\begin{aligned}\xi f  = & \big[-2 x^{\mathrm{T}} x + 2 d \widetilde{\eta}^\mathrm{T}  \alpha Bx\big] - 2 \lambda x^{\mathrm{T}} B^{-1} T \widetilde{\eta} +\\ & +2 \lambda d  \widetilde{\eta}^\mathrm{T} \alpha T \widetilde{\eta} - 2 \lambda x^\mathrm{T} B^{-1} \widetilde I + 2\lambda d \widetilde\eta^\mathrm{T} \alpha \widetilde I .\end{aligned}   
\end{equation*}
Using equation \eqref{2.3.5}, we derive the inequality:
\begin{equation*}\label{2.3.8} 
\begin{aligned} \xi f  \leqslant & \big[-2 \varepsilon b_m x^\mathrm{T} x \big]- 2 \lambda x^{\mathrm{T}} B^{-1} T \widetilde{\eta} +\\ & +2 \lambda d  \widetilde{\eta}^\mathrm{T} \alpha T \widetilde{\eta} - 2 \lambda x^\mathrm{T} B^{-1} \widetilde I + 2\lambda d \widetilde\eta^\mathrm{T} \alpha \widetilde I.
\end{aligned}      
\end{equation*} 
Rewriting $-2\varepsilon b_m z^{\mathrm{T}}z$ as the sum of two terms and multiplying the second term by $\lambda$, we obtain the following inequality (since $\lambda \in [0,1]$):  
\begin{equation}\label{2.3.9}
\begin{aligned}\xi f \leqslant & - \varepsilon b_m x^\mathrm{T} - \lambda \varepsilon b_m x^\mathrm{T} x - 2 \lambda x^{\mathrm{T}} B^{-1} T \widetilde{\eta} -\\ & -\big[ \lambda d \widetilde\eta^\mathrm{T}\big (\alpha(-T) +(-T)^\mathrm{T}\alpha\big)\widetilde\eta\big]- \\ &-2 \lambda x^\mathrm{T} B^{-1} \widetilde I + 2\lambda d \widetilde\eta^\mathrm{T} \alpha \widetilde I.
\end{aligned}
\end{equation}
Next, by adding and subtracting the term \\$\widetilde{\eta}^{\mathrm{T}}\big(\frac{\lambda}{\varepsilon b_m}(B^{-1}T)^{\mathrm{T}}(B^{-1}T)\big)\widetilde{\eta}$, we transform \eqref{2.3.9} into:  
\begin{multline}\label{2.3.10}
% \begin{aligned}
\xi f \leqslant  -\varepsilon b_m x^{\mathrm{T}}x -\lambda\bigg\|\sqrt{\varepsilon b_m} x - \frac{1}{\sqrt{\varepsilon b_m}} B^{-1}T \widetilde{\eta} \bigg\|^2-\\
- \widetilde{\eta}^{\mathrm{T}}\bigg( \lambda d \big(\alpha(-T)+(-T)^{\mathrm{T}}\alpha\big) - \\
-\frac{\lambda}{\varepsilon b_m}(B^{-1}T)^{\mathrm{T}}(B^{-1}T)\bigg) \widetilde{\eta} - \\
-2 \lambda x^\mathrm{T} B^{-1} \widetilde I + 2\lambda d \widetilde\eta^\mathrm{T} \alpha \widetilde I.
% \end{aligned}
\end{multline}
Applying condition \textbf{C3$'$}, we further transform \eqref{2.3.10} into:  
\begin{equation*}\label{2.3.11}
\xi f \leqslant -\varepsilon b_m x^{\mathrm{T}} x - 2 \lambda x^{\mathrm{T}} B^{-1} \widetilde{I} + 2 \lambda d \widetilde{\eta}^\mathrm{T} \alpha  \widetilde{I}.
\end{equation*}

Applying condition \textbf{C1} and the fact that $\widetilde\eta(t) \in \overline{\operatorname{conv}}\big(g(x(t))\big)$:
\begin{equation*}\label{2.3.12}
\begin{aligned}
    \xi f   \leqslant -\varepsilon b_m\|x\|^2&+2\left(\|-B^{-1} \widetilde{I}\| + a d  \|\alpha \widetilde{I}\|\right)\|x\|+ \\ 
    &+2 bd  \|\alpha \widetilde{I}\|.   
\end{aligned} 
\end{equation*}
Thus, if the radius $R$ of the closed sphere $K$ is sufficiently large, it can be asserted that for all $f\in \varPhi(\lambda,x)$:
\begin{equation*}\label{2.3.13}
\begin{aligned}
\xi f  &\leqslant -\varepsilon b_m R^2+2\left(\|-B^{-1} \widetilde{I}\| + a d  \|\alpha \widetilde{I}\|\right)R+\\&~~~~+2 bd  \|\alpha \widetilde{I}\|\\&<0 .     
\end{aligned}
\end{equation*}
Hence, zero cannot belong to the image of $\varPhi(\lambda, x)$ for all $x \in \partial K$ and $\lambda \in [0,1]$. And finally by the theorem 3.2.8 \cite[p. 90]{p3}, the neural network has an equilibrium point.
\end{IEEEproof}

%%%%%%%%%%%%%%%%%%%%%%%%%%%%%%%%%%%%%%%%%%%%%%%%%%%%%%%%%%%%%%%%%%%%%%%%%%%%%%%%%%%%%%%
\subsection{Global Asymptotic Stability of the Unique equilibrium point of the Network}
\begin{theorem}[Global Asymptotic Stability of the Unique equilibrium point of the Network]\label{main}


Let the activation function $g\in \mathcal{G'}$ and conditions \textbf{C1-C3$'$} are satisfied. Then for any input $I\in \mathbb{R}^n$, the network \eqref{eq:special} has a unique globally asymptotically stable equilibrium point.

\end{theorem}

\begin{remark}[The idea of proof] This theorem can be proven using the LaSalle invariance principle, modified for regular Lyapunov functions \cite{p9}. To define the derivative of the Lyapunov function, the concept of the generalized gradient by Clarke and the chain rule theorem for regular functions can be used.
\end{remark}

\begin{IEEEproof} 
By Theorem \ref{exists}, the neural network \eqref{eq:special} possesses an equilibrium point (state).
Let $e$ be an equilibrium point of the set-valued mapping $\varphi(x(t))$, i.e., by Definition \ref{equilib}:
\[
0 \in \varphi\big(e(t)\big)=Be+T~\overline{\operatorname{conv}}(g(e))+I ,
\]
and let $\eta^* \in \overline{\operatorname{conv}}(g(e))$ be the corresponding equilibrium point of the neural network's output, such that the following equality holds almost everywhere:
\[
0 = Be + T \eta^*+ I. \tag{a.e.}
\]

\begin{remark}
{The existence of $\eta^*$ can be established using reasoning analogous to that in the proof of Theorem \ref{exists}, i.e., by Makarov’s theorem \cite[p. 38, Th. 1.7]{Leonov}, for any solution $x(t)$, there exists a measurable bounded mapping $\eta(t) \in \overline{\operatorname{conv}}g(x(t))$ such that the equality in the inclusion holds almost everywhere.}
\end{remark}
Thus, $e$ is an equilibrium point of the neural network \eqref{eq:special} (since there is $\eta^*$ corresponding to it). 
Next we perform a change of variables $z=x-e$ to shift the equilibrium point to the origin and transform system \eqref{1.2} (see Definition \ref{fillipovsol}) into system \eqref{2.4.1}:
\[
\begin{aligned}
    \dot z(t) &\in Bz+ T~ \overline{\operatorname{conv}}(g(z+e))+Be+I=\\
    &= Bz+T~ \overline{\operatorname{conv}}(g(z+e)) - T\eta^*.\\
\end{aligned}
\]

Henceforth, we consider system \eqref{2.4.1}:
\begin{equation}\label{2.4.1}
\dot z(t) \in Bz+T~ \overline{\operatorname{conv}}(G(z)),
\end{equation}

where $G(z)=g(z+e)-\eta^*$.

Taking into account the validity of the theorem's conditions for the function $g\in \mathcal{G'}$, we reformulate them for the function $G$:
\begin{enumerate}
    \item $G\in \mathcal{G'}$ (see Definition \ref{G_n});
    
    \item Since $\overline{\operatorname{conv}}(G(0))=\overline{\operatorname{conv}}(g(e)) -\eta^*$, and $\eta^* \in \overline{\operatorname{conv}}(g(e))$, 
    
    it follows that $0 \in \overline{\operatorname{conv}}(G(0))$. Thus, by Definition \ref{equilib}, the origin is an equilibrium point of system \eqref{2.4.1};
    \item $ \frac{ \widetilde{\eta_i}}{z} \leqslant C_i $, $\forall x \ne 0, \quad \forall \widetilde{\eta_i} \in  \overline{\operatorname{conv}}[G_i(z)]$ (a consequence of condition \textbf{C2}).
\end{enumerate}

    It is important to note that such a measurable $\widetilde{\eta_i}$ exists by the remark above. Moreover, analogously to $e(t)$ and $\eta^*(t)$, for $z(t)$ and $\widetilde{\eta(t)} \in \overline{\operatorname{conv}}G(z(t))$, the following holds almost everywhere:
\begin{equation}\label{2.4.2}
\dot z (t) = Bz(t) + T \widetilde{\eta(t)}, \quad \widetilde{\eta(t)} = T ^{-1}\big(\dot z(t) - Bz(t)\big).     
\end{equation}

By Lemma \ref{main_lemma} and property 3 of the function $G$, there exist positive constants $d$ and $\varepsilon$ such that:
\begin{equation}\label{2.4.3}
\frac{\widetilde{\eta_i}}{z} \leqslant C_i \leqslant \frac{1 - \varepsilon b_i}{d \alpha_i b_i},\quad \forall z \ne 0,\quad \forall \widetilde{\eta_i} \in  \overline{\operatorname{conv}}[G_i(z)].    
\end{equation}


where $\alpha = \operatorname{diag}\{\alpha_1,...,\alpha_n\},$ $\alpha_i>0$ is the matrix from condition \textbf{C3$'$}.
Moreover, inequality \eqref{2.3.3} holds.

To analyze the stability of system \eqref{2.4.1}, we consider a Lur'e-Postnikov-type energy function \cite{p17} for the Hopfield neural network, which the author originally proposed as a Lyapunov function for stability analysis in his 1986 paper \cite{p10}.

\begin{equation}\label{2.4.4}
V(z) ~= ~\sum\limits_{i=1}^n\frac {1}{b_i}z_i^2 +2d\sum\limits_{i=1}^n\alpha_i\int\limits_0^{z_i}G_i(s)ds.    
\end{equation}

By condition \textbf{C2}, the integral of the nonlinearity satisfies:
\[
\int\limits_0^{z_i}G_i(s)ds \leqslant\frac {1}{2} C_iz_i^2.
\]\\
Let us elaborate on this fact. For $z_i \geqslant 0 $, condition \textbf{C2} implies $G_i(s) \leqslant C_i s$ for $s\in [0,z_i]$. Integrating this inequality yields:
\[
\int\limits_0^{z_i} G_i (s) ds \leqslant\int\limits_0^{z_i}C_i sds = \frac{1}{2}C_iz_i^2.
\]
For $z_i < 0$: $G_i(s) \geqslant C_i s$, where $s\in [z_i,0]$. Integrating gives:
\[
\int\limits_0^{z_i}G_i(s)ds = - \int\limits_{z_i}^0G_i(s)ds \leqslant - \int\limits_{z_i}^0C_isds= \frac{1}{2}C_iz_i^2.
\]
From inequality \eqref{2.4.3}, we have $\forall i \in 1 : n$:
\[
\frac{1}{b_i}z_i^2+ 2 d \alpha_i \int\limits_0^{z_i} G_i(s)ds \geqslant \varepsilon z_i^2.
\]
Thus, for the function $V(z)$, the following holds:
\[
V(z) = \sum_{i=1}^n\frac{1}{b_i}z_i^2 + 2 d \sum_{i=1}^n\alpha_i \int\limits_0^{z_i}G_i(s)ds \geqslant \varepsilon \|z\|^2.
\]
Note that this function is regular, being a convex nonsmooth function in $\mathbb{R}^n$ \cite{516}. Additionally, $V(0)=0$. From the above derivations, it is also positive definite ($V(x)>0, \quad \forall x\ne 0$ and $V(0)=0$) and radially unbounded ($V(x)\to +\infty ~if ~\|x\| \to +\infty$).

From formula \eqref{2.4.4}, the Clarke generalized gradient \cite{p14} is given by:
\[
\partial V(z) = -2B^{-1}z + 2d\alpha~\overline{\operatorname{conv}}\big(G(z(t))\big).
\]
Now, applying the chain rule for regular functions \cite{p9}, we express the derivative of $V(z(t))$ along the trajectories of system \eqref{2.4.1} as a scalar product:
\[
\frac{\mathbf{d}V(z(t))}{\mathbf{d}t} = \big( \xi(t), \dot z(t) \big), \tag{a.e.}
\]
where $\xi(t) \in \partial V(z(t))$. Let us choose:
\[
\xi(t) = -2B^{-1}z(t)+2d\alpha~\widetilde{\eta(t)}, \quad \widetilde{\eta(t)} \in \overline{\operatorname{conv}}G(z(t)).
\]
Taking into account the choice of $\xi (t)$ and formula \eqref{2.4.2}, we expand:
\begin{equation}\label{2.4.5}
\begin{aligned} & \frac{\mathbf{d}V(z(t))}{\mathbf{d}t} = \big(-2B^{-1}z+2d\alpha \widetilde{\eta}\big)^{\mathrm{T}}\big(Bz+T\widetilde{\eta}) =\\ 
% & = -2z^{\mathrm{T}}z+2d\widetilde{\eta}^{\mathrm{T}}\alpha B z - 2z^{\mathrm{T}}B^{-1}T\widetilde{\eta} + 2d \widetilde{\eta}^{\mathrm{T}} \alpha T \widetilde{\eta} =  \\
& = -2z^{\mathrm{T}}z+2d\widetilde{\eta}^{\mathrm{T}}\alpha B z -  2z^{\mathrm{T}}B^{-1}T\widetilde{\eta} - \\
&~~~~~~~~~~~~~~~~~~~~~- \big[2d \widetilde{\eta}^{\mathrm{T}} \alpha (-T) \widetilde{\eta} \big]= \\
& = -2z^{\mathrm{T}}z+2d\widetilde{\eta}^{\mathrm{T}}\alpha B z - 2z^{\mathrm{T}}B^{-1}T\widetilde{\eta} - \\
& ~~~~~~~~~~~~~~~~~~~~~- \big[ d \widetilde{\eta}^{\mathrm{T}}\big(\alpha(-T)+(-T)^{\mathrm{T}}\alpha\big) \widetilde{\eta}\big] \leqslant\\
& \leqslant -2\varepsilon b_m z^{\mathrm{T}}z - 2z^{\mathrm{T}}B^{-1}T\widetilde{\eta} -\\&~~~~~~~~~~~~~~~~~~~~~- d \widetilde{\eta}^{\mathrm{T}}\big(\alpha(-T)+(-T)^{\mathrm{T}}\alpha\big) \widetilde{\eta}.
\end{aligned}
\end{equation}
We represent $-2\varepsilon b_m z^{\mathrm{T}}z$ as two terms, then add and subtract $\widetilde{\eta}^{\mathrm{T}}\big(\frac{1}{\varepsilon b_m}(B^{-1}T)^{\mathrm{T}}(B^{-1}T)\big)\widetilde{\eta}$, transforming \eqref{2.4.5} into:
\begin{equation*}\label{2.4.6}
\begin{aligned} 
% & = -\varepsilon b_m z^{\mathrm{T}}z  -\varepsilon b_m z^{\mathrm{T}}z - 2z^{\mathrm{T}}B^{-1}T\widetilde{\eta} + d \widetilde{\eta}^{\mathrm{T}}\big(\alpha(-T)+(-T)^{\mathrm{T}}\alpha\big) \widetilde{\eta} =\\
&\frac{\mathbf{d}V(z(t))}{\mathbf{d}t}  =  \\
&=-\varepsilon b_m z^{\mathrm{T}}z -\bigg\|\sqrt{\varepsilon b_m} z - \frac{1}{\sqrt{\varepsilon b_m}} B^{-1}T \widetilde{\eta} \bigg\|^2-\\
& \quad- \widetilde{\eta}^{\mathrm{T}}\bigg(d \big(\alpha(-T)+(-T)^{\mathrm{T}}\alpha\big) - \\&~~~~~~~~~~~~~~~~~~~~~~~~~~~-\frac{1}{\varepsilon b_m}(B^{-1}T)^{\mathrm{T}}(B^{-1}T)\bigg) \widetilde{\eta} \leqslant\\
& \leqslant-\varepsilon b_m z^\mathrm{T}z - \Lambda \widetilde{\eta}^{\mathrm{T}}\widetilde{\eta},
\end{aligned} 
\end{equation*}
where $\Lambda = d \lambda_M - \frac{1}{\varepsilon b_m}\|-B^{-1}T\|^2_2>0$ by inequality \eqref{2.3.3}. Since $b_m$ and $\varepsilon$ are positive constants, $\frac{\mathrm{d}}{\mathrm{d} t} V(z(t))$ is negative definite, and by LaSalle invariance principle \cite[Cor. 4.2, p. 129]{p223} applied it to a Lyapunov function that is differentiable almost everywhere, the zero equilibrium point of system \eqref{2.4.1} is globally asymptotically stable. Consequently, the same holds for the equilibrium point $e$ of system \eqref{eq:special}. 

The uniqueness of the equilibrium point can be shown analogously to the reasoning in \cite{eq_ref_11, period_din, eq_ref_2}, where discontinuous activation functions for neural networks were also considered. Suppose that in the transformed system \eqref{2.4.1}, for which the origin is an equilibrium point, there exists another equilibrium point $z^* \ne 0$, and $\eta^*$ is the corresponding output equilibrium point (see Definition \ref{equilib}). Again by chain rule for regular functions~\cite{p9}, the expression for the derivative of the Lyapunov function along the system trajectories is:
\[
\frac{\mathbf{d}V(z^*)}{\mathbf{d}t} = \big(-2B^{-1}z^*+2d\alpha {\eta^*}\big)^{\mathrm{T}}\big(Bz^*+T {\eta}^*) = 0 .
\]
By definition of the system's equilibrium point, the right-hand term vanishes, which contradicts the negative definiteness ($\frac{\mathbf{d}V(z(t))}{\mathbf{d}t} < 0 $ for $z \ne 0$).
\end{IEEEproof}

\section{Proofs for Application I}\label{app:control-proofs}

\section{Proofs for Application II (QUBO/MAX-CUT)}\label{app:qubo-proofs}

\begin{IEEEproof}[Proof of Lemma~\ref{lem:qubo-unique}]
Fix $i \in \{1,\dots,n\}$ and let $s^{(i)}$ be the configuration obtained
from $s^\star$ by flipping the $i$th spin:
\[
  s^{(i)}_j =
  \begin{cases}
    -s_i^\star, & j = i,\\
    s_j^\star,  & j \neq i.
  \end{cases}
\]
We compute the energy difference
$\Delta_i := E(s^{(i)}) - E(s^\star)$.
Write $E(s) = \frac{1}{2} s^\top Q s - h^\top s$ and
let $\delta := s^{(i)} - s^\star$.
Since only the $i$th component flips, we have
$\delta = -2 s_i^\star e_i$, where $e_i$ is the $i$th canonical basis vector.
Then
\begin{align*}
  \Delta_i
  &= \frac{1}{2}\bigl((s^\star+\delta)^\top Q (s^\star+\delta)
      - (s^\star)^\top Q s^\star\bigr) - h^\top \delta \\
  &= s^{\star\top} Q \delta + \frac{1}{2} \delta^\top Q \delta - h^\top \delta.
\end{align*}
Using $\delta_j = 0$ for $j \neq i$ and $\delta_i = -2 s_i^\star$, we obtain
\[
  s^{\star\top} Q \delta = (Q s^\star)_i \delta_i,
  \qquad
  \delta^\top Q \delta = \delta_i^2 Q_{ii} = 4 Q_{ii},
  \qquad
  h^\top \delta = h_i \delta_i.
\]
Hence
\begin{align*}
  \Delta_i
  &= (Q s^\star)_i (-2 s_i^\star) + 2 Q_{ii} - h_i (-2 s_i^\star) \\
  &= -2 s_i^\star (Q s^\star)_i + 2 Q_{ii} + 2 h_i s_i^\star.
\end{align*}
Next expand $(Q s^\star)_i = Q_{ii} s_i^\star + \sum_{j \neq i} Q_{ij} s_j^\star$.
Then
\[
  s_i^\star (Q s^\star)_i
  = Q_{ii} + s_i^\star \sum_{j \neq i} Q_{ij} s_j^\star,
\]
and therefore
\begin{align*}
  \Delta_i
  &= -2\Bigl[Q_{ii} + s_i^\star \sum_{j \neq i} Q_{ij} s_j^\star\Bigr]
     + 2 Q_{ii} + 2 h_i s_i^\star \\
  &= -2 s_i^\star \sum_{j \neq i} Q_{ij} s_j^\star + 2 h_i s_i^\star \\
  &= 2 s_i^\star \Bigl(h_i - \sum_{j \neq i} Q_{ij} s_j^\star\Bigr).
\end{align*}
Using $s_i^\star = \operatorname{sign}(h_i)$ and the strong field condition
\eqref{eq:strong-field}, we estimate
\[
  \Bigl| \sum_{j \neq i} Q_{ij} s_j^\star \Bigr|
  \leq \sum_{j \neq i} |Q_{ij}|
  < |h_i|.
\]
Hence
\[
  s_i^\star\Bigl(h_i - \sum_{j \neq i} Q_{ij} s_j^\star\Bigr)
  \geq |h_i| - \Bigl|\sum_{j \neq i} Q_{ij} s_j^\star\Bigr|
  > 0,
\]
and thus $\Delta_i > 0$ for all $i=1,\dots,n$:
flipping any single spin in $s^\star$ strictly increases the energy $E$.
Therefore $s^\star$ is a strict local minimizer of~$E$.

Now take any $s \in \{-1,1\}^n$ with $s \neq s^\star$.
There exists a finite sequence of single-spin flips transforming $s^\star$
into $s$:
\[
  s^\star = s^{(0)} \to s^{(1)} \to \dots \to s^{(k)} = s.
\]
At each step the energy strictly increases,
$E(s^{(\ell+1)}) > E(s^{(\ell)})$, $\ell=0,\dots,k-1$.
Hence
\[
  E(s) = E(s^{(k)}) > E(s^{(k-1)}) > \dots > E(s^{(0)}) = E(s^\star),
\]
which shows that $s^\star$ is the unique global minimizer of~$E$.
\end{IEEEproof}

\medskip
Lemma~\ref{lem:qubo-unique} identifies a nontrivial subclass of QUBO problems
(``QUBO with strong field'') for which the optimizer is unique and given by
$s^\star = \operatorname{sign}(h)$.

\medskip
\begin{IEEEproof}[Proof of Theorem~\ref{th:hnn-solves-strong-field-qubo}]
By Lemma~\ref{lem:qubo-unique} and the strong field
condition~\eqref{eq:strong-field} the QUBO problem~\eqref{eq:qubo-energy}
has a unique global minimizer $s^\star = \operatorname{sign}(h)$.

The activation~\eqref{eq:sign-activation} satisfies \textbf{C1--C2} and, by construction,
$T = -\beta Q_{\mathrm{shift}}$ with $Q_{\mathrm{shift}} \succ 0$ implies
$T+T^\top \prec 0$, so that \textbf{C3} holds (with $\alpha=I$).
Therefore, by Theorem~\ref{th:global-stability}, system~\eqref{eq:hnn-qubo}
has a unique globally asymptotically stable Filippov equilibrium $x^\star$,
and any Filippov solution $x(\cdot)$ of~\eqref{eq:hnn-qubo} satisfies
$x(t) \to x^\star$ as $t \to +\infty$.

At a Filippov equilibrium one has
\[
  0 \in B x^\star - \beta Q_{\mathrm{shift}} G^\star + \beta h,
\]
where $G^\star$ is the Filippov set of possible activation values at $x^\star$.
Under the strong field condition~\eqref{eq:strong-field} the external field
$h_i$ cannot be balanced by the couplings $\{Q_{ij}\}_{j \neq i}$ if $x_i^\star = 0$:
for any convex combination of neighbor contributions we still have
$|h_i| > \sum_{j \neq i} |Q_{ij}|$, so $0$ cannot belong to the right-hand side
of the $i$th scalar equilibrium inclusion.
Hence no equilibrium can have $x_i^\star = 0$, and therefore
$g(x^\star) = \operatorname{sign}(x^\star)$ is single-valued and lies in
$\{-1,1\}^n$.

Denote $s := g(x^\star) = \operatorname{sign}(x^\star)$.
The equilibrium condition for~\eqref{eq:hnn-qubo} can be written as
\[
  0 = B x^\star - \beta Q_{\mathrm{shift}} s + \beta h,
\]
that is,
\[
  x^\star = B^{-1} \beta (Q_{\mathrm{shift}} s - h).
\]
Thus any equilibrium uniquely determines a spin vector $s \in \{-1,1\}^n$.
Conversely, if $s \in \{-1,1\}^n$ and we define
$x(s) := B^{-1} \beta (Q_{\mathrm{shift}} s - h)$, then $(x(s),s)$ would be
an equilibrium pair provided that $\operatorname{sign}(x(s)) = s$.
The same single-spin flip argument as in Lemma~\ref{lem:qubo-unique} and the
strong field condition~\eqref{eq:strong-field} imply that the only spin vector
consistent with such a relation is precisely $s^\star$.
Therefore $g(x^\star) = s^\star$.

Finally, by Lemma~\ref{lem:qubo-unique}, $s^\star$ is the unique global minimizer
of the QUBO energy~\eqref{eq:qubo-energy}.
This completes the proof.
\end{IEEEproof}

\section{Proofs for Application III (Penalty \& SVM)}\label{app:penalty-proofs}



\end{document}
